\documentclass[11pt,a4paper]{report}
\usepackage{fontspec,amsmath,amssymb,amsthm,microtype,xcolor}
\setmainfont{lmroman10-regular.otf}[BoldFont=lmroman10-bold.otf,ItalicFont=lmroman10-italic.otf,BoldItalicFont=lmroman10-bolditalic.otf]
\usepackage[margin=25mm]{geometry}
\usepackage[unicode,hidelinks]{hyperref}
\hypersetup{pdfauthor={LangevinLean Project},bookmarksdepth=2}
\setlength{\emergencystretch}{3em}
\setlength{\parskip}{0.4em}
\setlength{\parindent}{0pt}
\setcounter{tocdepth}{1}
\makeatletter
\renewcommand*\l@section{\@dottedtocline{1}{1.5em}{3.8em}}
\renewcommand*\@pnumwidth{2.5em}
\renewcommand*\@tocrmarg{3.5em}
\makeatother
\setcounter{secnumdepth}{2}
\allowdisplaybreaks[2]
\raggedbottom
\title{\texorpdfstring{Unadjusted Langevin under LSI}{Unadjusted Langevin under LSI}\\[0.6em]\large Mathematical proofs}
\author{LangevinLean Project}
\date{Draft report LL-02\\[0.4em]\small Content revision e64aa62fb836}
\begin{document}
\pagenumbering{roman}
\maketitle
\chapter*{About this document}
\addcontentsline{toc}{chapter}{About this document}
These draft notes present the assumptions, convergence theorems and detailed mathematical proofs for this development. The supporting arguments are arranged so that prerequisites precede their applications. They use ordinary mathematical notation throughout.

Prepared by the LangevinLean Project with AI assistance, these notes form part of an ongoing effort to formalize Langevin algorithms comprehensively. This is draft work and may contain mistakes.

\section*{Literature and scope}
\begin{enumerate}
\item Santosh S. Vempala and Andre Wibisono (2019). \emph{Rapid Convergence of the Unadjusted Langevin Algorithm: Isoperimetry Suffices}. Advances in Neural Information Processing Systems 32, 8094–8106. \url{https://proceedings.neurips.cc/paper/2019/hash/65a99bb7a3115fdede20da98b08a370f-Abstract.html}.
\end{enumerate}
These proofs use results established in the earlier project reports \cite{LangevinLean2026ULASC}. Each application below gives the section containing the required result. The earlier proofs are not reproduced in this report.

\tableofcontents
\clearpage
\pagenumbering{arabic}
\chapter{Assumptions and main results}
\section{\texorpdfstring{KL convergence of the Gaussian ULA algorithm}{KL convergence of the Gaussian ULA algorithm}}


Let \(d\ge1\) and \(\nu(dx)=Z^{-1}e^{-U(x)}dx\) be a probability measure on \(\mathbb R^d\). Assume \(U\) is differentiable with globally \(L\)-Lipschitz gradient, \(L>0\), and \(\nu\) satisfies a log-Sobolev inequality with constant \(\alpha>0\):


\[2\alpha H_\nu(\rho)\le I_\nu(\rho),\qquad
H_\nu(\rho)=\int\log\frac{d\rho}{d\nu}\,d\rho,\quad
I_\nu(\rho)=\int\left\|\nabla\log\frac{d\rho}{d\nu}\right\|^2d\rho.\]


Take an initial probability law \(\rho_0\) with finite relative entropy. Let \(Z_k\) be independent standard Gaussian vectors, independent of initialization. For


\[X_{k+1}=X_k-h\nabla U(X_k)+\sqrt{2h}\,Z_k,
\qquad 0<h\le\frac{\alpha}{4L^2},\]


the laws \(\rho_k=\operatorname{Law}(X_k)\) satisfy, for every \(k\ge0\),


\[\boxed{H_\nu(\rho_k)\le e^{-\alpha hk}H_\nu(\rho_0)+\frac{8hdL^2}{\alpha}.}\]


This is the KL estimate of Vempala–Wibisono (2019), Theorem 2. The first term is the initial entropy contracted at an exponential rate. The second is the discretization floor. Convexity of \(U\) is not assumed. The proof uses the compatible inequality \(\alpha\le L\), derived below in positive dimension. If the initial entropy is infinite the extended-valued inequality is immediate; the finite case is the substantive assertion.


{\small\color{gray}Supporting proof: Section~\ref{proof-module-UnadjustedLangevin-LSI-UnconditionalTheorem-lean}.\par}
\chapter{Mathematical argument}
\section{\texorpdfstring{The Gibbs law, the actual update, and the LSI test class}{The Gibbs law, the actual update, and the LSI test class}}


Write \(\nu(dy)=J(y)dy\), where \(J=Z^{-1}e^{-U}>0\). A positive C\ensuremath{{}^{1}} relative density \(r\), normalized by \(\int r\,d\nu=1\), defines the law \(r\nu\). The functional inequality is


\[2\alpha\,\mathrm{KL}(r\nu\mid\nu)\le\int\|\nabla\log r\|^2r\,d\nu.\]


Both quantities initially take values in the extended nonnegative reals. Their finiteness is established before ordinary real integral manipulations are used. It applies this same LSI definition to the actual Gaussian mixture, whose strict positivity and C\ensuremath{{}^{1}} regularity are proved below.




For one step, take \(X\sim\mu\) independently of \(Z_0\sim\mathcal N(0,I_d)\), and set


\[Y_t=X-t\nabla U(X)+\sqrt{2t}\,Z_0,\quad\rho_t=\operatorname{Law}(Y_t).\]


The independent product measure makes independence explicit. At zero the law is \(\mu\), and at \(h\) it is its actual ULA update. The discrete laws are repeated applications of this transition.



The relative density is distinct from a Lebesgue density. If \(p_t\) is the density of \(\rho_t\) with respect to volume, then \(r_t=p_t/J\); consequently \(\log r_t=\log p_t+U+\log Z\) and the relative score is \(S_t=\nabla\log p_t+\nabla U\). For the exact Gibbs density, these two score contributions cancel and Fisher information is zero. This explains why the functional inequality controls distance from Gibbs rather than merely spatial variation of \(p_t\).



The interpolation is local to a numerical step: its initial random variable remains fixed in the frozen drift. It is therefore not the exact Langevin diffusion during that step. The proof will quantify the discrepancy between \(\nabla U(X)\) and \(\nabla U(Y_t)\), while retaining the exact law equality at both endpoints. Those endpoint equalities are what allow a continuous entropy estimate to become a recurrence for the actual algorithm.

\section{\texorpdfstring{Derive \ensuremath{\alpha} \ensuremath{\le} L with bounded tests}{Derive α ≤ L with bounded tests}}


We first derive compatibility, rather than assume it. For bounded C\ensuremath{{}^{1}} \(f\) with \(\|\nabla f\|\le D\), let \(\Lambda(s)=\log\int e^{sf}d\nu\). The normalized density \(e^{sf-\Lambda(s)}\) is an admissible LSI test, and its log-gradient is \(s\nabla f\). Its entropy is \(s\Lambda'(s)-\Lambda(s)\). Consequently


\[2\alpha\bigl(s\Lambda'(s)-\Lambda(s)\bigr)\le s^2D^2.\]


Boundedness permits differentiating the moment generating function near zero. Divide by \(s^2\) and let \(s\downarrow0\): the quotient \((s\Lambda'(s)-\Lambda(s))/s^2\) tends to \(\Lambda''(0)/2=\operatorname{Var}_\nu(f)/2\), so the entire left side tends to \(\alpha\operatorname{Var}_\nu(f)\). Thus \(\alpha\operatorname{Var}_\nu(f)\le D^2\).




Choose a unit vector \(v\). Translation invariance of the partition function and the smooth upper model \cite[Section~3.9.1]{LangevinLean2026ULASC} give a centered sub-Gaussian directional score \(g_v=\langle\nabla U,v\rangle\) \cite[Section~3.27.3]{LangevinLean2026ULASC}, with \(\int g_v^2d\nu\le L\) by the directional second-moment estimate \cite[Section~3.27.5]{LangevinLean2026ULASC}. Integration by parts for bounded C\ensuremath{{}^{1}} tests gives


\[\int f g_v\,d\nu=\int\langle\nabla f,v\rangle d\nu.\]


Use \(f_t(x)=\sin(t\langle v,x\rangle)\). Its gradient norm is at most \(t\), and the identity reads \(\int(f_t-\nu f_t)g_vd\nu=t\int\cos(t\langle v,x\rangle)d\nu\). Cauchy–Schwarz and the bounded-test variance bound yield


\[\alpha\left(\int\cos(t\langle v,x\rangle)d\nu\right)^2\le L.\]


The cosine is bounded by one and converges pointwise to one. Dominated convergence gives \(\alpha\le L\). This uses no position second moment, Hessian, or coordinate approximation. If the state space is a singleton, all probability measures coincide and KL is zero; that case is handled separately.



To see the cancellation of the small parameter explicitly, write \(D=t\) for the sine test. The variance bound gives \(\operatorname{Var}(f_t)\le t^2/\alpha\), while the score has mean zero and second moment at most \(L\). Thus

\[t^2\left(\int\cos(t\langle v,x\rangle)d\nu\right)^2
=\left(\int(f_t-\nu f_t)g_vd\nu\right)^2
\le\frac{Lt^2}{\alpha}.\]

For \(t>0\), divide by \(t^2\) and multiply by \(\alpha\). The limiting cosine integral is one without any assumption on \(\langle v,X\rangle\), because its integrand is bounded. Positive dimension enters only through the existence of a unit direction; all scalar integrations remain bounded-test calculations.

\section{\texorpdfstring{Obtain the initial gradient moment directly from entropy}{Obtain the initial gradient moment directly from entropy}}


The directional score estimate \cite[Section~3.27.2]{LangevinLean2026ULASC} also gives, for every vector \(v\),


\[\int e^{\langle\nabla U(x),v\rangle}d\nu(x)\le e^{L\|v\|^2/2}.\]


Average this inequality over an auxiliary centered Gaussian vector. Completing the square in its density gives an exponential square moment. Tonelli applies because every integrand is nonnegative. With the Gaussian variance chosen to match \(4L\),


\[\int\exp\!\left(\frac{\|\nabla U\|^2}{4L}\right)d\nu
\le 2^{d/2}\le e^{d/2}.\]


For any finite-KL probability law \(\mu\), the entropy variational inequality then gives


\[G_0:=\int\|\nabla U\|^2d\mu\le4L\,H_\nu(\mu)+2dL
\le\frac{4L^2}{\alpha}H_\nu(\mu)+2dL.\]


The last inequality uses \(\alpha\le L\) and nonnegative KL. Integrability is part of the conclusion. Its integrability is derived before using the real integral: for the Radon–Nikodym density \(a=d\mu/d\nu\), the pointwise inequality \(ab\le a\log a-a+e^b\) controls the nonnegative weighted integrand. Subtracting the certified log-moment bound inside the Fenchel inequality gives the sharper variational bound. Thus finite KL alone supplies the score moment required by the Gaussian calculations.



The entropy variational inequality can be read directly as a density calculation. If \(a=d\mu/d\nu\) and \(F\) has a finite exponential moment, normalize \(e^F\nu\) by \(M=\int e^Fd\nu\). Nonnegativity of relative entropy with respect to this tilted probability gives

\[0\le\operatorname{KL}\left(\mu\,\middle\Vert\,M^{-1}e^F\nu\right)
=H_\nu(\mu)-\int Fd\mu+\log M.\]

The pointwise convex-duality estimate first supplies the absolute integrability required to write this identity. For the nonnegative score-square observable, monotone truncations also justify passage from bounded tests to the full \(F\). Taking \(F=\|\nabla U\|^2/(4L)\) and \(\log M\le d/2\) yields the stated coefficient after multiplication by \(4L\).

\section{\texorpdfstring{Compute displacement and drift-\allowbreak{}error energies}{Compute displacement and drift-error energies}}


On the same independent product probability space,


\[Y_t-X=-t\nabla U(X)+\sqrt{2t}\,Z_0.\]


Expand its squared norm. The cross term integrates to zero because the Gaussian mean is zero and the two factors are independent. The Gaussian second moment is \(\mathbb E\|Z_0\|^2=d\), hence


\[\mathbb E\|Y_t-X\|^2=t^2G_0+2td.\]


The Lipschitz gradient now gives the exact error estimate used in production:


\[\mathbb E\|\nabla U(Y_t)-\nabla U(X)\|^2
\le L^2(t^2G_0+2td).\]


All terms are integrable by the score moment just proved and the Gaussian second moment. The joint law of \((X,Y_t)\) is also identified as \(K_t(x,y)\mu(dx)dy\), where \(K_t\) is the Gaussian transition density. This transfers the same sharp bound to the kernel-weighted pair integral used in the weak equation.



In more detail, independence and Cauchy–Schwarz justify

\[\mathbb E\langle\nabla U(X),Z_0\rangle
=\left\langle\mathbb E\nabla U(X),\mathbb EZ_0\right\rangle=0.\]

The initial gradient is integrable because it is square-integrable under a probability law. Hence this is an ordinary expectation of an integrable mixed term, rather than a formal cancellation. No term involving \(\mathbb E\|X\|^2\) appears: the position cancels before any square is expanded. Multiplying the displacement energy by \(L^2\) preserves its two distinct orders, \(t^2G_0\) from the frozen drift and \(2td\) from Gaussian diffusion. These orders determine the separate initial-entropy and dimensional forcing terms later.

\section{\texorpdfstring{Prove density regularity, Fisher finiteness, and the weak equation}{Prove density regularity, Fisher finiteness, and the weak equation}}


For \(t>0\), write


\[K_t(x,y)=(4\pi t)^{-d/2}e^{-\|y-x+t\nabla U(x)\|^2/(4t)},\quad
p_t(y)=\int K_t(x,y)d\mu(x),\quad r_t=p_t/J.\]


Gaussian polynomials times their decaying exponential are bounded. They dominate the first spatial derivative independently of the center, so differentiation under the probability integral proves that \(p_t\), and therefore \(r_t\), are positive C\ensuremath{{}^{1}} functions. The posterior probability on \(x\), proportional to \(K_t(x,y)\mu(dx)\), gives


\[\nabla\log p_t(y)=-\frac1{2t}\mathbb E[y-X+t\nabla U(X)\mid Y_t=y].\]


Jensen and the Gaussian displacement moment imply \(\int\|\nabla\log p_t\|^2d\rho_t\le d/(2t)\). Since \(S_t:=\nabla\log r_t=\nabla\log p_t+\nabla U\),


\[I_t:=\int\|S_t\|^2d\rho_t\le d/t+2\int\|\nabla U\|^2d\rho_t<\infty.\]


Lipschitzness and the displacement identity bound the last integral by \(2(1+L^2t^2)G_0+4tdL^2\). Applying the original extended-valued LSI now proves finite actual KL for every \(t>0\), as well as \(2\alpha H(t)\le I_t\).




Directly differentiating the kernel in time gives


\[\dot K_t=K_t\left[-\frac d{2t}+\frac{\|y-x\|^2}{4t^2}-\frac{\|\nabla U(x)\|^2}{4}\right].\]


Its absolute integral is controlled on each positive compact time interval by \(G_0\) and Gaussian moments. Set \(F_t(x,y)=\nabla_y K_t(x,y)+K_t(x,y)\nabla U(x)\). Its spatial divergence equals \(\dot K_t\). Coordinate integration by parts, with integrable derivatives and flux, proves for every bounded C\ensuremath{{}^{1}} test \(\psi\) with \(\nabla\psi\in L^2(\rho_t)\),


\[\int\dot p_t\psi\,dy=-\iint\langle F_t(x,y),\nabla\psi(y)\rangle d\mu(x)dy.\]


This is the actual weak equation; no weak-equation hypothesis is passed to the final theorem. The jointly measurable score field equals \(S_t\) at positive times, so its spacetime integrals are well-defined.



The smoothing Fisher bound follows by integrating posterior Jensen with its denominator canceled:

\[\begin{aligned}
\int p_t(y)\|\nabla\log p_t(y)\|^2dy
&\le\iint\frac{\|y-x+t\nabla U(x)\|^2}{4t^2}
K_t(x,y)d\mu(x)dy\\
&=\frac{2td}{4t^2}=\frac d{2t}.
\end{aligned}\]

The inner Gaussian moment is the same for every initial point, so no location moment of \(\mu\) is introduced. The bound diverges as \(t\downarrow0\), which is why the entropy identity is first proved on intervals bounded away from zero. For the weak equation, a bounded test makes the flux-test product integrable, finite gradient energy makes the flux-gradient product integrable, and the explicit quadratic Gaussian envelope makes the flux derivative integrable. These are the precise global integration-by-parts conditions replacing an unproved boundary-decay assertion.

\section{\texorpdfstring{Pass bounded log tests to the entropy identity}{Pass bounded log tests to the entropy identity}}


An unbounded logarithm cannot be inserted into the weak equation without justification. For \(m\ge0\), put


\[\begin{aligned}
\chi_m(z)&=(m+1)\arctan\frac z{m+1},\\
w_m(z)&=\chi_m'(z)=\frac1{1+(z/(m+1))^2},\\
\Phi_m(r)&=\int_1^r\chi_m(\log u)du.
\end{aligned}\]


These functions satisfy \(0\le w_m\le1\), \(w_m\to1\), and \(0\le\Phi_m(r)\le r\log r-r+1\), with pointwise convergence to the upper bound. The test \(\psi_{m,t}=\chi_m(\log r_t)\) is bounded C\ensuremath{{}^{1}}, with \(\nabla\psi_{m,t}=w_m(\log r_t)S_t\). Its gradient is square-integrable because \(I_t<\infty\).




The time chain rule is first applied pointwise to \(J\Phi_m(r_t)\): its derivative is \(\chi_m(\log r_t)\dot p_t\). Boundedness of \(\chi_m\) and the spacetime L\ensuremath{{}^{1}} estimate for \(\dot p_t\) justify Fubini and the fundamental theorem of calculus on \([a,b]\subset(0,\infty)\). The weak equation identifies this bounded-entropy flux with a weighted Gaussian pair integral.




The Fisher bound and the Gaussian velocity energy are integrable uniformly on every such interval. Cauchy–Schwarz dominates the absolute cross term by half the sum of these energies. Therefore dominated convergence sends \(m\to\infty\) in both the entropy endpoints and the time-integrated flux. The result is local absolute continuity of actual KL and an entropy derivative almost everywhere. Averaging the Gaussian component score gives the mixture score, so the identity becomes


\[H'(t)=-I_t+\mathbb E\langle\nabla U(Y_t)-\nabla U(X),S_t(Y_t)\rangle.\]


The averaging is an integrable Fubini identity for the actual kernel; it does not require choosing a conditional drift jointly in time. Apply \(\langle a,b\rangle\le\|a\|^2+\|b\|^2/4\) and the sharp error energy:


\[\boxed{H'(t)\le-\tfrac34I_t+t^2L^2G_0+2tdL^2}\quad\text{for almost every }t>0.\]


This is sufficient for scalar integration. No pointwise entropy derivative at every time is assumed.



It is useful to keep the two limiting measures distinct. At a fixed endpoint time, the dominating function is the nonnegative entropy integrand \(r_t\log r_t-r_t+1\), integrated against Gibbs. For the time flux, the domination is instead a spacetime energy bound on one independent state-noise probability. Passing both limits in the integrated bounded-entropy identity yields

\[H(b)-H(a)=\int_a^b\left[-I_t+
\mathbb E\langle\nabla U(Y_t)-\nabla U(X),S_t(Y_t)\rangle\right]dt.\]

The integrand is integrable on every positive compact interval. This identity itself proves local absolute continuity and gives the derivative almost everywhere by the fundamental theorem of calculus. Applying Young's inequality spends exactly one quarter of \(I_t\); the remaining three quarters can still be converted into a strictly dissipative entropy term through LSI.

\section{\texorpdfstring{Control the initial endpoint for every finite-\allowbreak{}entropy law}{Control the initial endpoint for every finite-entropy law}}


A one-sided entropy bound suffices to send \(a\downarrow0\); we can construct it directly without assuming an initial smooth density. Data processing for the same independent Gaussian transition gives


\[\mathrm{KL}(\mu P_t\mid\nu P_t)\le\mathrm{KL}(\mu\mid\nu).\]


For \(\delta=1-4tL>0\), set \(s=t/\delta\). The first-order smoothness estimate \cite[Section~3.9.1]{LangevinLean2026ULASC} and gradient Lipschitzness yield the pointwise square-completion bound


\[U(y)-U(x)-\frac{\|y-x+t\nabla U(x)\|^2}{4t}
\le-\frac{\|x-y+s\nabla U(y)\|^2}{4s}
+\frac s4\|\nabla U(y)\|^2.\]


Integrate in \(x\). The ratio of Gaussian normalizers is \((s/t)^{d/2}=\delta^{-d/2}\), hence


\[\frac{d(\nu P_t)}{d\nu}(y)\le\delta^{-d/2}
\exp\!\left(\frac t{4\delta}\|\nabla U(y)\|^2\right).\]


The chain rule for log likelihood ratios, followed by the propagated gradient bound, proves


\[H(t)\le M(t):=H(0)-\frac d2\log\delta
+\frac t{4\delta}\bigl[2(1+L^2t^2)G_0+4tdL^2\bigr].\]


The proof establishes the needed likelihood-ratio integrability, including the negative part, before adding the entropies. The explicit majorant satisfies \(M(0)=H(0)\) and is continuous at zero. Consequently integrating any affine entropy differential inequality on \([a,h]\) and then sending \(a\downarrow0\) is valid. Finite initial KL is the only initial-law regularity assumption. The proof works directly for differentiable potentials with Lipschitz gradient, so no potential-smoothing limit is needed either.



To apply this endpoint control, suppose the positive-time comparison reads \(H'\le-aH+B\) with \(a>0\). Integration on \([s,h]\), where \(s>0\), gives

\[H(h)\le e^{-a(h-s)}H(s)+\frac Ba\bigl(1-e^{-a(h-s)}\bigr).\]

Insert \(H(s)\le M(s)\) and let \(s\downarrow0\). Continuity of \(M\) gives precisely the same formula with \(H(0)\) on the right. This argument requires an upper endpoint envelope rather than two-sided entropy continuity. It therefore accommodates every finite-entropy initial law covered by the theorem, while the density differentiations remain confined to strictly positive elapsed time.

\section{\texorpdfstring{Absorb the forcing and obtain the coefficient six}{Absorb the forcing and obtain the coefficient six}}


Combine LSI, the score moment, and production:


\[H'\le-\tfrac32\alpha H+t^2L^2\left(\frac{4L^2}{\alpha}H(0)+2dL\right)+2tdL^2.\]


From \(h\le\alpha/(4L^2)\) and \(\alpha\le L\), we have \(hL\le1/4\) and \(\alpha h\le1/4\). In particular,


\[h^2L^2\frac{4L^2}{\alpha}\le\frac\alpha4,\qquad
h^2L^2(2dL)\le\tfrac12hdL^2.\]


Thus the forcing for \(0<t<h\) is at most \(B=\alpha H(0)/4+(5/2)hdL^2\). Local absolute continuity permits multiplication by the integrating factor and integration from \(a>0\). The continuous initial majorant just proved permits \(a\downarrow0\), giving


\[H(h)\le e^{-3\alpha h/2}H(0)+\frac B{3\alpha/2}(1-e^{-3\alpha h/2})
\le e^{-3\alpha h/2}H(0)+Bh.\]


For \(0\le x\le1/4\), \(e^{-x}\ge e^{-3x/2}+x/4\). Apply this with \(x=\alpha h\), multiply by nonnegative \(H(0)\), and absorb its forcing contribution. The remaining \((5/2)h^2dL^2\) is at most \(6h^2dL^2\). Therefore


\[\boxed{H_\nu(\mu P_h)\le e^{-\alpha h}H_\nu(\mu)+6h^2dL^2.}\]


The scalar comparison needs derivatives only almost everywhere and no density extension through zero.



The signs used in absorption are part of the comparison. Nonnegative KL allows multiplication of the exponential-gap estimate by \(H(0)\), and nonnegative \(B\) allows the forcing kernel to be bounded by its interval length. Explicitly,

\[\int_0^he^{-(3\alpha/2)(h-t)}Bdt\le Bh
=\frac{\alpha h}{4}H(0)+\frac52h^2dL^2.\]

The first term is paid for by replacing \(e^{-3\alpha h/2}\) with the slower factor \(e^{-\alpha h}\). The second is a genuine discretization error, retained as the additive forcing in the discrete recurrence. The coefficient six is a convenient enlargement of \(5/2\); the proof makes no claim that this intermediate coefficient is optimal.

\section{\texorpdfstring{Prove finite KL at every iterate and sum the recurrence}{Prove finite KL at every iterate and sum the recurrence}}


Finite KL at each iterate is proved first by induction: finite current KL gives a square-integrable gradient, the Gaussian Fisher estimate gives finite relative Fisher after a positive step, and LSI then gives finite next KL. Thus the real-valued sequence \(H_k\) faithfully represents the actual divergences.




Write \(q=e^{-\alpha h}\), \(c=6h^2dL^2\). The actual one-step estimate and the finite-horizon affine recurrence \cite[Section~3.15.2]{LangevinLean2026ULASC} imply


\[H_k\le q^kH_0+c\frac{1-q^k}{1-q}.\]


For \(0\le x\le1/4\), the exponential tangent bound gives


\[1-e^{-x}\ge\frac{x}{1+x}\ge\tfrac34x.\]


Use \(x=\alpha h\), discard \(1-q^k\le1\), and simplify:


\[\frac{6h^2dL^2}{1-e^{-\alpha h}}\le\frac{8hdL^2}{\alpha}.\]


This proves the displayed convergence bound with the original potential, actual independent-Gaussian ULA laws, and every finite-KL initial probability law. The final theorem assumes smoothness, LSI, the step restriction, and finite initial KL. All density, moment, weak-equation, entropy-regularity, compatibility, and recurrence facts have been proved within its dependency chain.



The induction can be written in its exact finite-horizon form. Starting from \(H_0\), each insertion of \(c\) is damped once for every remaining transition. The geometric-sum form \cite[Section~3.15.1]{LangevinLean2026ULASC} is

\[H_k\le e^{-\alpha hk}H_0+
6h^2dL^2\sum_{j=0}^{k-1}e^{-\alpha hj}.\]

The geometric denominator is strictly positive because both \(\alpha\) and \(h\) are positive. Its lower bound is proportional to \(h\), turning the local \(h^2\) error into a global \(h\) floor. At fixed step, increasing \(k\) removes only the initial-entropy contribution. Reducing the asymptotic floor requires a smaller step; the separate dependences on dimension, smoothness, and LSI constant remain explicit.

\chapter{Supporting proofs}
We now supply the detailed supporting arguments. Each section introduces the objects and estimates needed by the results that follow it.

\section{\texorpdfstring{The relative-\allowbreak{}density log-\allowbreak{}Sobolev convention}{The relative-density log-Sobolev convention}}
\label{proof-module-UnadjustedLangevin-LSI-FunctionalInequality-lean}


Write \(\rho=r\nu\), \(H_\nu(\rho)=\int r\log r\,d\nu\), and \(J_\nu(\rho)=\int\|\nabla\log r\|^2d\rho\). This module defines the explicit smooth positive-density LSI convention \(2\alpha H\le J\) and guards real-valued conversions by finite divergence and Fisher information.



A relative density \(r\) is measured against the target \(\nu\), so its induced law is \(\rho_r=r\nu\). For a positive normalized density, \(H_\nu(r)=\int r\log r\,d\nu\) when the integrand is integrable, and \(I_\nu(r)=\int r\|\nabla\log r\|^2d\nu\). The convention \(2\alpha H_\nu(r)\le I_\nu(r)\) fixes the meaning of the LSI constant. These are initially nonnegative extended quantities; finiteness must precede ordinary real integral manipulations.

\subsection{\texorpdfstring{The law specified by a relative density}{The law specified by a relative density}}
\label{report-step-UnadjustedLangevin-LSI-FunctionalInequality-lean-1}


For a real relative density \(r\), define \(\rho_r(dx)=r(x)\nu(dx)\) using a nonnegative scalar \ensuremath{\circ} r). Positivity and normalization are later hypotheses, rather than part of this raw measure definition. A normalized positive \(r\) is the density of the current law relative to the target.



The density specifies a measure before it specifies a probability. To use entropy as a nonnegative divergence, require \(r\ge0\) and \(\int r\,d\nu=1\). Then integration against the induced law is ordinary weighted integration against \(\nu\). Strict positivity is added for logarithmic derivatives, rather than being needed merely to form the weighted measure.

\subsection{\texorpdfstring{Relative Fisher information}{Relative Fisher information}}
\label{report-step-UnadjustedLangevin-LSI-FunctionalInequality-lean-2}


Define \(I_\nu(r)=\int\|\nabla\log r\|^2d\rho_r\) as a nonnegative extended integral. The specified positive density is differentiated pointwise. Subsequent energy estimates establish finiteness before this quantity is used as an ordinary real number.



The relative score differentiates the chosen positive density \(r\) with respect to the reference \(\nu\). For a Gibbs reference with Lebesgue density \(J\), if \(r=p/J\), it is \(\nabla\log p+\nabla U\). Thus Fisher information measures departure from the target score, rather than just roughness of the Lebesgue density.

\subsection{\texorpdfstring{The log-\allowbreak{}Sobolev inequality}{The log-Sobolev inequality}}
\label{report-step-UnadjustedLangevin-LSI-FunctionalInequality-lean-3}


The target is a probability and satisfies \(2\alpha\operatorname{KL}(\rho_r\mid\nu)\le I_\nu(r)\) for every strictly positive \(C^1\) normalized relative density. The factor \(2\alpha\) fixes the convention used throughout the proof. This test class and its normalization are explicit; the definition contains no algorithm or entropy-decay claim.



The hypothesis is a functional inequality for every member of the displayed positive \(C^1\) normalized test class. It is not an assertion that every probability law already has a smooth density. In the ULA application, Gaussian smoothing proves that the actual positive-time relative density belongs to this class before LSI is invoked.

\subsection{\texorpdfstring{The finite real log-\allowbreak{}Sobolev bound}{The finite real log-Sobolev bound}}
\label{report-step-UnadjustedLangevin-LSI-FunctionalInequality-lean-4}


Apply LSI to the positive \(C^1\) normalized density. If KL and Fisher information are finite, this is an inequality between ordinary nonnegative numbers, giving \(2\alpha H_\nu(r)\le I_\nu(r)\) with the same normalization.



The extended inequality does not by itself justify ordinary subtraction or multiplication when infinite values are possible when infinity is possible. Once KL and Fisher are finite, both are ordinary nonnegative real quantities and the coefficient \(2\alpha\) has its usual value. The displayed real LSI then follows without changing its normalization.

\section{\texorpdfstring{Integrate a forced entropy-\allowbreak{}dissipation inequality}{Integrate a forced entropy-dissipation inequality}}
\label{proof-module-UnadjustedLangevin-LSI-EntropyDecay-lean}


An integrating factor solves \(H\prime\le-aH+b\) as \(H(t)\le e^{-at}H(0)+(b/a)(1-e^{-at})\), under the displayed positive-rate, continuity and one-sided derivative assumptions. The entropy result supplies the appropriate rate from LSI and dissipation.



The integrating factor for \(H'\le-aH+b\) is \(e^{at}\). Integrating the weighted inequality gives \(H(t)\le e^{-at}H(0)+b\int_0^te^{-a(t-s)}ds\), and the integral equals \((1-e^{-at})/a\) for \(a>0\). The same comparison follows from the stated one-sided derivative hypotheses. In the unforced case \(b=0\), LSI with coefficient \(2\alpha\) gives decay at that exact rate.

\subsection{\texorpdfstring{A scalar affine comparison from right derivatives}{A scalar affine comparison from right derivatives}}
\label{report-step-UnadjustedLangevin-LSI-EntropyDecay-lean-1}


For continuous \(H\) with right derivative bounded by \(-aH+b\), \(a>0\), compare it to the exact scalar ODE solution. The one-sided Grönwall theorem accepts the right derivative through its right-slope liminf inequality. Simplifying its explicit solution gives \(H(t)\le e^{-at}H(0)+(b/a)(1-e^{-at})\). Neither \(H\) nor \(b\) needs to be nonnegative.



Subtract the exact affine ODE solution from \(H\). Its upper right derivative is bounded by \(-a\) times the difference, and the difference starts at zero. The one-sided comparison gives a nonpositive difference throughout the interval.

\subsection{\texorpdfstring{Scalar LSI and dissipation give rate 2\ensuremath{\alpha}}{Scalar LSI and dissipation give rate 2α}}
\label{report-step-UnadjustedLangevin-LSI-EntropyDecay-lean-2}


The supplied dissipation identity is \(H'=-I\), and the scalar LSI premise is \(2\alpha H\le I\). Therefore \(H'\le-2\alpha H\). Apply the affine comparison theorem with \(a=2\alpha\) and \(b=0\), obtaining \(H(t)\le e^{-2\alpha t}H(0)\). This scalar theorem does not itself establish the dissipation identity for a law-valued evolution.



The coefficient two is inherited from the stated normalization of LSI. The unforced dissipation identity spends the entire Fisher term, so its scalar decay rate is \(2\alpha\); the later frozen-drift estimate spends part of that term on error absorption.

\section{\texorpdfstring{Relative entropy along an identified density path}{Relative entropy along an identified density path}}
\label{proof-module-UnadjustedLangevin-LSI-EntropyFlow-lean}


Density entropy is the real conversion of the actual KL divergence of the tilted measure. With finite-entropy guards, LSI and a supplied dissipation inequality give exponential entropy decay along this path.



For an admissible density path, the exact entropy identity \(H'=-I\) and the inequality \(I\ge2\alpha H\) imply \(H'\le-2\alpha H\). Multiplication by \(e^{2\alpha t}\) then shows that the weighted entropy is nonincreasing. This continuous-flow calculation explains the role of LSI. A frozen-drift numerical interpolation has an additional error pairing, so its later decay rate and bias require a separate estimate rather than this unforced identity alone.

\subsection{\texorpdfstring{Define actual density entropy with its finiteness guard}{Define actual density entropy with its finiteness guard}}
\label{report-step-UnadjustedLangevin-LSI-EntropyFlow-lean-1}


Define the real entropy of a specified relative density as \(H_\nu(r)=\operatorname{KL}(r\nu\mid\nu)\). It uses the same the probability law specified by a relative density and actual relative entropy as the LSI definition. Applications must separately establish finite KL before using it as an ordinary entropy value.



For a normalized density, the familiar entropy integral agrees with KL only after the needed integrability is established. The present definition retains the actual divergence as its starting point. Consequently later derivative or comparison statements must first identify a finite entropy integral with this quantity, rather than differentiate an arbitrary real-converted extended value.

\subsection{\texorpdfstring{Apply the finite real LSI at every time}{Apply the finite real LSI at every time}}
\label{report-step-UnadjustedLangevin-LSI-EntropyFlow-lean-2}


At each interpolation time, positivity, \(C^1\) regularity, normalization, and finite KL/\allowbreak{}Fisher allow the inequality \(2\alpha H\le I\) to give \(2\alpha H(t)\le I(t)\). Combine this with the explicit right-dissipation identity and entropy continuity in the integrating-factor comparison for \(H\prime\le-2\alpha H\). The conclusion is decay at rate \(2\alpha\) for that admissible density path; identifying an actual evolution with the path remains an application obligation.



At each time LSI is applied to the same density defining the law of the path. Together with the supplied dissipation identity it gives \(H'+2\alpha H\le0\). Integration of \((e^{2\alpha t}H(t))'\le0\) yields the displayed decay. The argument proves a scalar consequence for the admissible path; establishing that a particular stochastic evolution has this path is a separate law-identification step.

\section{\texorpdfstring{Identify density entropy and Fisher integrals}{Identify density entropy and Fisher integrals}}
\label{proof-module-UnadjustedLangevin-LSI-KLOfDensity-lean}


For a strictly positive relative density, with the stated differentiability and integrability, the actual likelihood ratio equals \(\log r\). This identifies the actual KL divergence with \(\int r\log r\,d\nu\), proves its finiteness and identifies the real Fisher value with its score-energy integral.



Changing from \(\rho_r\) to \(\nu\) inserts the factor \(r\) into every integral. In particular, \(\int|\log r|d\rho_r=\int r|\log r|d\nu\), so the integrability conditions in the two measures are equivalent. The logarithmic Radon–Nikodym derivative is \(\log r\) almost everywhere under \(\rho_r\). For probability measures the mass correction in the general relative-entropy formula vanishes, leaving the usual \(\int r\log r\,d\nu\).

\subsection{\texorpdfstring{The log likelihood of a positive relative density}{The log likelihood of a positive relative density}}
\label{report-step-UnadjustedLangevin-LSI-KLOfDensity-lean-1}


The weighted-density measure is absolutely continuous with respect to \(\nu\), and its Radon–Nikodym derivative equals a nonnegative scalar r \(\nu\)-almost everywhere. Absolute continuity transfers that equality to \(\rho_r\)-almost everywhere. Taking the real logarithm gives the actual log-likelihood ratio \(\log r\).



The equality is needed under the tilted law, where the KL integrand is integrated. Absolute continuity transfers the reference-almost-everywhere density identity to that law, allowing the logarithm to be identified on the correct full-measure set.

\subsection{\texorpdfstring{Finiteness of KL from weighted log integrability}{Finiteness of KL from weighted log integrability}}
\label{report-step-UnadjustedLangevin-LSI-KLOfDensity-lean-2}


Absolute continuity follows from the weighted-density construction. Integrability of \(r\log r\) gives integrability of \(\log r\) under \(\rho_r\), and the log-likelihood identity identifies that with the KL integrand. The standard finite-KL criterion therefore proves \(\operatorname{KL}(\rho_r\mid\nu)<\infty\).



Weighted integrability gives \(\int|\log r|d\rho_r<\infty\). With absolute continuity, this is the finite-KL criterion for the actual measure pair, not merely finiteness of a formally defined entropy expression.

\subsection{\texorpdfstring{The ordinary integral formula for entropy}{The ordinary integral formula for entropy}}
\label{report-step-UnadjustedLangevin-LSI-KLOfDensity-lean-3}


For normalized \(\nu\) and \(\rho_r\), the finite-KL integral formula has mass correction \(\nu(E)-\rho_r(E)=0\). Replace the log-likelihood ratio by \(\log r\), then change density to obtain \(H_\nu(r)=\int r\log r\,d\nu\). The integrability premise is used before invoking that formula.



Both measures have mass one, so the mass-correction term disappears. The density-change identity converts \(\int\log r\,d\rho_r\) into \(\int r\log r\,d\nu\), yielding the usual formula with its integrability already justified.

\subsection{\texorpdfstring{The ordinary integral formula for Fisher information}{The ordinary integral formula for Fisher information}}
\label{report-step-UnadjustedLangevin-LSI-KLOfDensity-lean-4}


Nonnegativity and integrability identify the real value of the extended Fisher integral with the ordinary integral under \(\rho_r\). The weighted-measure formula then rewrites it as \(I_\nu(r)=\int r\|\nabla\log r\|^2d\nu\). Both representations refer to the same supplied density.



The squared score is nonnegative. Its finite integral against \(\rho_r\) therefore has the same real value in the weighted reference representation \(\int r\|\nabla\log r\|^2d\nu\), which is the energy appearing in integration by parts.

\subsection{\texorpdfstring{Changing measure by a relative density}{Changing measure by a relative density}}
\label{report-step-UnadjustedLangevin-LSI-KLOfDensity-lean-5}


For measurable \(r\ge0\), the Bochner weighted-density formula gives \(\int g\,d\rho_r=\int r(x)g(x)d\nu\). The embedded density is finite pointwise and has real value \(r\), so the theorem's extended-to-real scalar action simplifies to ordinary multiplication.



For a nonnegative density, changing measure multiplies the integrand by \(r\). The statement applies to signed or vector-valued tests once the corresponding absolute integrability has been checked.

\subsection{\texorpdfstring{Equivalent log-\allowbreak{}integrability conditions}{Equivalent log-integrability conditions}}
\label{report-step-UnadjustedLangevin-LSI-KLOfDensity-lean-6}


Apply the weighted-density integrability criterion to \(g=\log r\). Since the density is measurable, finite, and nonnegative, it yields \(\log r\in L^1(\rho_r)\iff r\log r\in L^1(\nu)\). This is the exact entropy-integrability bridge used before using the finite KL integral formula.



The criterion concerns absolute values: \(\int|\log r|d\rho_r=\int r|\log r|d\nu\). This verifies both signs of the log likelihood before the KL formula is written as an ordinary signed integral.

\subsection{\texorpdfstring{Finiteness of relative Fisher information}{Finiteness of relative Fisher information}}
\label{report-step-UnadjustedLangevin-LSI-KLOfDensity-lean-7}


The score square is nonnegative and measurable. If its ordinary integral is finite, its nonnegative extended integral has the same finite value. That integral is precisely relative Fisher information, establishing the finiteness needed by the real LSI comparison.



A finite ordinary integral of a nonnegative measurable score square equals its nonnegative extended integral. Thus the latter cannot be infinite, proving the Fisher finiteness required before scalar LSI comparison.

\section{\texorpdfstring{Relative scores turn the weak equation into entropy production}{Relative scores turn the weak equation into entropy production}}
\label{proof-module-UnadjustedLangevin-LSI-ScoreCalculus-lean}


The chain rule \(\nabla\log r=\nabla r/r\) converts the entropy test \(\log r+1\) into the relative score. The weak identity supplied for this test then gives Fisher dissipation, with an extra score/\allowbreak{}error pairing when drift forcing is present.



For a positive differentiable density, its relative score is \(s=\nabla\log r=\nabla r/r\). Hence \(\langle\nabla r,\nabla(\log r+1)\rangle=r\|s\|^2\). This algebra explains why the entropy test produces Fisher information. The test \(\log r+1\) need not have compact support or be bounded, so its use in a weak equation requires the explicit admissibility and integrability assumptions stated here; later bounded logarithmic tests establish the actual Gaussian identity without assuming this insertion.

\subsection{\texorpdfstring{The gradient of a positive density logarithm}{The gradient of a positive density logarithm}}
\label{report-step-UnadjustedLangevin-LSI-ScoreCalculus-lean-1}


The logarithm chain rule gives \(D\log r(x)=r(x)^{-1}Dr(x)\) when \(r(x)>0\). The real Hilbert-space gradient is the Riesz representative of the derivative; linearity of that identification gives \(\nabla\log r(x)=r(x)^{-1}\nabla r(x)\).



Strict positivity supplies the reciprocal and permits the ordinary logarithm chain rule. The derivative-to-gradient identification is linear, so the scalar factor \(1/r(x)\) multiplies the gradient without any additional metric term.

\subsection{\texorpdfstring{Adding one does not change the log-\allowbreak{}test gradient}{Adding one does not change the log-test gradient}}
\label{report-step-UnadjustedLangevin-LSI-ScoreCalculus-lean-2}


Adding a constant to a scalar function leaves its Fréchet derivative unchanged. Apply the gradient's derivative-based definition to obtain \(\nabla(\log r+1)=\nabla\log r\). The explicit differentiability/\allowbreak{}positivity arguments align the statement with the entropy-test calculus.



The additive constant affects the entropy derivative as a scalar integrand but has zero spatial gradient. Thus it disappears from the Dirichlet pairing, leaving exactly the relative score.

\subsection{\texorpdfstring{The Fisher integrand in the entropy test}{The Fisher integrand in the entropy test}}
\label{report-step-UnadjustedLangevin-LSI-ScoreCalculus-lean-3}


Rearrange the positive-density score identity to \(\nabla r=r\nabla\log r\). Pair this with \(\nabla(\log r+1)=\nabla\log r\) and use \(\langle s,s\rangle=\|s\|^2\). The Dirichlet integrand is therefore \(\langle\nabla r,\nabla(\log r+1)\rangle=r\|\nabla\log r\|^2\).



Multiplying \(\nabla\log r=\nabla r/r\) by the positive density recovers \(\nabla r\). Pairing with the log-test gradient then gives the nonnegative Fisher integrand pointwise.

\subsection{\texorpdfstring{Entropy production from the weak identity}{Entropy production from the weak identity}}
\label{report-step-UnadjustedLangevin-LSI-ScoreCalculus-lean-4}


The strictly positive \(C^1\) density makes \(\psi=\log r+1\) a \(C^1\) test. The entropy derivative integrability is supplied, and the pointwise gradient identity converts its Dirichlet integrability to the supplied weighted Fisher integrability. Substitute this admissible test into the explicit weak equation to obtain \(\int\dot r(\log r+1)d\nu=-\int r\|\nabla\log r\|^2d\nu\).



The weak equation must admit this particular unbounded test, with all its required pairings integrable. Under that explicit premise, the pointwise score identity converts the Dirichlet term into actual Fisher dissipation.

\subsection{\texorpdfstring{Admissibility of the noncompact entropy test}{Admissibility of the noncompact entropy test}}
\label{report-step-UnadjustedLangevin-LSI-ScoreCalculus-lean-5}


Use the same \(C^1\) entropy test in the forced weak equation. Rewrite the Dirichlet term by the identity \(\langle\nabla r,\nabla(\log r+1)\rangle=r\|\nabla\log r\|^2\) and the forcing gradient by the equality \(\nabla(\log r+1)=\nabla\log r\). The three explicit integrability conditions justify the substitution, giving entropy production \(-I+\int r\langle e,\nabla\log r\rangle d\nu\). Admissibility of this noncompact test remains part of the weak-equation premise.



The forcing contributes \(\int r\langle e,\nabla\log r\rangle d\nu\). Its measurability and absolute integrability are separate requirements from finite entropy, and are retained when passing from the weak equation to this production identity.

\section{\texorpdfstring{Why the LSI constant cannot exceed smoothness}{Why the LSI constant cannot exceed smoothness}}
\label{proof-module-UnadjustedLangevin-LSI-Compatibility-lean}


Write \(\nu(dx)=Z^{-1}e^{-U(x)}dx\), with \(\nabla U\) globally \(L\)-Lipschitz, in a nonzero finite-dimensional real Hilbert space. The original positive \(C^1\) density formulation of LSI implies \(\alpha\le L\). The proof first derives a variance inequality for bounded smooth tests, then combines bounded sinusoidal tests with Gibbs integration by parts and the previously proved score bound. No position moment, transport inequality, Hessian, or strong convexity is assumed.



For bounded \(C^1\) tests, exponential tilting and the second-order expansion of the log moment generating function turn LSI into \(\alpha\operatorname{Var}_\nu(f)\le\|\nabla f\|_\infty^2\). Apply this to \(f_t(x)=\sin(t\langle v,x\rangle)\) for a unit vector \(v\). A Gibbs integration-by-parts identity and the directional score variance bound \(L\) give \(\alpha(\int\cos(t\langle v,x\rangle)d\nu)^2\le L\). Bounded convergence as \(t\downarrow0\) proves \(\alpha\le L\) without a position second moment.



Earlier results used here: \cite[Section~3.27]{LangevinLean2026ULASC}.
\subsection{\texorpdfstring{Integrability of bounded exponential tests}{Integrability of bounded exponential tests}}
\label{report-step-UnadjustedLangevin-LSI-Compatibility-lean-1}


If \(|f|\le B\), then \(tf\le |t|B\), so \(0<e^{tf}\le e^{|t|B}\). The right side is an integrable constant under every finite measure. Measurability follows from that of \(f\); integrable domination then proves actual integrability of the exponential. This establishes the whole real line as the domain of the moment generating function used below.



Boundedness of the test gives \(e^{sf}\le e^{|s|\|f\|_\infty}\). This constant is integrable under the target probability, uniformly for \(s\) in a compact interval, allowing derivatives of its moment generating function.

\subsection{\texorpdfstring{Testing LSI with a normalized exponential tilt}{Testing LSI with a normalized exponential tilt}}
\label{report-step-UnadjustedLangevin-LSI-Compatibility-lean-2}


Let \(M(t)=\int e^{tf}\,d\nu\), \(\psi(t)=\log M(t)\), and \(r_t=e^{tf}/M(t)\). Boundedness of \(f\) proves \(0<M(t)<\infty\); the normalized tilt is therefore a probability measure. The proof constructs strict positivity and \(C^1\) regularity of r before testing the original LSI definition.



The normalized density is \(r_s=e^{sf}/\int e^{sf}d\nu\). Its logarithmic gradient is \(s\nabla f\), and its entropy is \(s\Lambda'(s)-\Lambda(s)\), giving the displayed LSI inequality for every scalar tilt.

\subsection{\texorpdfstring{The small-\allowbreak{}tilt entropy quotient}{The small-tilt entropy quotient}}
\label{report-step-UnadjustedLangevin-LSI-Compatibility-lean-3}


For a \(C^3\) scalar function \(\psi\) with \(\psi(0)=0\), both numerator and denominator of \((t\psi'(t)-\psi(t))/t^2\) tend to zero as \(t\downarrow0\). The product rule cancels the two copies of \(\psi'(t)\), leaving numerator derivative \(t\psi''(t)\); the denominator derivative is \(2t\). On the punctured right neighborhood, their quotient equals \(\psi''(t)/2\). Continuity of the second derivative and one application of l'Hopital's rule prove the limit \(\psi''(0)/2\). The proof explicitly supplies the nonzero denominator derivative and both zero limits.



The constant and first-order terms cancel in \(s\Lambda'(s)-\Lambda(s)\). Its leading term is \(s^2\Lambda''(0)/2\), where \(\Lambda''(0)\) is the variance of the bounded test.

\subsection{\texorpdfstring{A variance bound for bounded differentiable tests}{A variance bound for bounded differentiable tests}}
\label{report-step-UnadjustedLangevin-LSI-Compatibility-lean-4}


For bounded \(f\), its moment and cumulant generating functions are analytic on all real parameters. the standard second derivative formula at zero is \(\psi''(0)=\int f^2\,d\nu-(\int f\,d\nu)^2\), since \(M(0)=1\). Thus no unproved assertion that LSI implies a spectral gap is inserted.



Divide the exponential-test inequality by \(t^2>0\), for \(t>0\), and use the second-order small-tilt entropy expansion. The left side tends to \(\alpha\operatorname{Var}_\nu(f)\), while the right side remains \(D^2\). Passing to the limit gives \(\alpha\operatorname{Var}_\nu(f)\le D^2\). The boundedness and gradient-bound hypotheses here describe this test function; they are not additional assumptions on the target.



Dividing the tilted LSI by \(s^2\) and taking the small-tilt limit gives \(\alpha\operatorname{Var}(f)\le D^2\). The bounded test class makes this step independent of unknown position moments.

\subsection{\texorpdfstring{Weighted Gibbs integration}{Weighted Gibbs integration}}
\label{report-step-UnadjustedLangevin-LSI-Compatibility-lean-5}


Unfold \(\nu=Z^{-1}(e^{-U}dx)\). Scalar multiplication of a measure multiplies its real Bochner integral by \((Z^{-1})\). The weighted-density integral formula then gives \(\int g\,d\nu=(Z^{-1})\int e^{-U}g\,dx\). Continuity of \(U\) makes the weight measurable, and its pointwise value is finite. This equality concerns total integrals; integrability for its later uses is established separately.



The Gibbs change of measure inserts its ordinary positive density into the Lebesgue integral. This is the form in which the weighted derivative identity can be integrated by parts.



Earlier results used here: \cite[Section~3.26.1]{LangevinLean2026ULASC}; \cite[Section~3.26.2]{LangevinLean2026ULASC}; \cite[Section~3.26.3]{LangevinLean2026ULASC}.
\subsection{\texorpdfstring{Weighted Gibbs integrability}{Weighted Gibbs integrability}}
\label{report-step-UnadjustedLangevin-LSI-Compatibility-lean-6}


Strict positivity of the Gibbs partition follows from differentiability of \(U\); the hypothesis here also excludes an infinite partition. Consequently multiplication of the weighted Lebesgue measure by \(Z^{-1}\) preserves integrability in both directions. The weighted-density integrability equivalence then gives \(g\in L^1(\nu)\) exactly when \(e^{-U}g\in L^1(dx)\). This is the justification for moving the bounded-test integration-by-parts calculation to Lebesgue measure.



The bounded test and gradient bounds, together with finite score energy, control the weighted test-score products. These integrability checks justify the Gibbs integration-by-parts identity on the noncompact state space.



Earlier results used here: \cite[Section~3.26.1]{LangevinLean2026ULASC}; \cite[Section~3.26.2]{LangevinLean2026ULASC}; \cite[Section~3.26.3]{LangevinLean2026ULASC}; \cite[Section~3.26.9]{LangevinLean2026ULASC}.
\subsection{\texorpdfstring{The Gibbs integration-\allowbreak{}by-\allowbreak{}parts identity for bounded tests}{The Gibbs integration-by-parts identity for bounded tests}}
\label{report-step-UnadjustedLangevin-LSI-Compatibility-lean-7}


For bounded \(C^1\) \(\psi\) with bounded gradient and any fixed \(v\), prove \(\mathbb E_\nu[\psi\langle\nabla U,v\rangle]=\mathbb E_\nu\langle\nabla\psi,v\rangle\). The already established score MGF implies integrability of \(\langle\nabla U,v\rangle\). Multiplying it by bounded \(\psi\) preserves integrability. Bounded \(\nabla\psi\) also makes its directional derivative integrable.



Differentiating the Gibbs weight in direction \(v\) gives \(-J\langle\nabla U,v\rangle\). Integration by parts therefore identifies the test-score expectation with the expected directional derivative of the test, with a positive sign.



Earlier results used here: \cite[Section~3.26.2]{LangevinLean2026ULASC}; \cite[Section~3.26.3]{LangevinLean2026ULASC}; \cite[Section~3.26.9]{LangevinLean2026ULASC}; \cite[Section~3.27.3]{LangevinLean2026ULASC}.
\subsection{\texorpdfstring{The sine-\allowbreak{}score covariance bound}{The sine-score covariance bound}}
\label{report-step-UnadjustedLangevin-LSI-Compatibility-lean-8}


For real \(L^2\) functions \(f,g\), Holder's inequality proves \(fg\in L^1\); their squares are integrable as well. For every real \(c\), integrate \((cf+g)^2\ge0\) to obtain the quadratic \((\int f^2)c^2+2(\int fg)c+\int g^2\ge0\). Its discriminant is nonpositive. Dividing the resulting algebraic inequality by four gives \((\int fg)^2\le(\int f^2)(\int g^2)\). This elementary squared Cauchy-Schwarz form avoids square roots in the compatibility argument.



Center the bounded sine test before applying Cauchy–Schwarz against the centered directional score. The score variance bound and the bounded-test variance estimate then control the square of their covariance.

\subsection{\texorpdfstring{Compatibility of LSI and smoothness constants}{Compatibility of LSI and smoothness constants}}
\label{report-step-UnadjustedLangevin-LSI-Compatibility-lean-9}


Choose \(\|v\|=1\), possible because the space is nontrivial. Put \(s(x)=\langle\nabla U(x),v\rangle\). The score MGF gives \(\mathbb E s^2\le L\); the bounded Stein identity with \(\psi=1\) gives \(\mathbb E s=0\).



For \(t>0\), use the bounded test \(f_t(x)=\sin(t\langle v,x\rangle)\). Its gradient is \(t\cos(t\langle v,x\rangle)v\), so its norm is at most \(|t|\). The bounded-test variance theorem gives \(\alpha\operatorname{Var}(f_t)\le t^2\). Stein's identity and score centering give \(\mathbb E[s(f_t-\mathbb Ef_t)]=t\mathbb E\cos(t\langle v,X\rangle)\). Squared Cauchy-Schwarz now yields \(t^2(\mathbb E\cos(t\langle v,X\rangle))^2\le L\operatorname{Var}(f_t)\).



Multiply by \(\alpha\ge0\), insert the variance bound, and cancel \(t^2>0\). This proves \(\alpha(\mathbb E\cos(t\langle v,X\rangle))^2\le L\). Dominated convergence applies because cosine is bounded by one and tends pointwise to one as \(t\downarrow0\). Its expectation tends to one, proving \(\alpha\le L\). No integrability of \(X\) itself is needed.



For the unit direction, the derivative of the sine test is \(t\cos(t\langle v,x\rangle)\). Cancelling the positive factor \(t^2\) gives the cosine inequality, and bounded convergence sends its expectation to one.



Earlier results used here: \cite[Section~3.26.2]{LangevinLean2026ULASC}; \cite[Section~3.26.3]{LangevinLean2026ULASC}; \cite[Section~3.27.3]{LangevinLean2026ULASC}; \cite[Section~3.27.5]{LangevinLean2026ULASC}.
\subsection{\texorpdfstring{The positive-\allowbreak{}dimensional compatibility conclusion}{The positive-dimensional compatibility conclusion}}
\label{report-step-UnadjustedLangevin-LSI-Compatibility-lean-10}


This result removes two temporary restrictions from the previous proof. If \(\alpha<0\), then \(\alpha\le0\le L\) immediately. Otherwise, LSI includes the assertion that the Gibbs law is a probability measure. An infinite partition would make its normalizing scalar zero and hence make the Gibbs measure zero, contradicting mass one. Thus the partition is finite internally, and the nonnegative-constant compatibility theorem applies. The final hypotheses are precisely LSI, differentiability, Lipschitz gradient, finite dimension, and a nontrivial state space.



In positive dimension a unit vector exists, so the preceding directional argument proves \(\alpha\le L\). In the singleton case no unit direction is available; the final entropy theorem handles that case directly with zero KL.



Earlier results used here: \cite[Section~3.26.2]{LangevinLean2026ULASC}; \cite[Section~3.26.3]{LangevinLean2026ULASC}.
\section{\texorpdfstring{Finite entropy supplies the gradient moment directly}{Finite entropy supplies the gradient moment directly}}
\label{proof-module-UnadjustedLangevin-LSI-EntropyMoment-lean}


A finite initial KL divergence already controls the squared Gibbs score when \(\nabla U\) is Lipschitz. Gaussian averaging of the directional score MGF gives \(\mathbb E_\nu e^{\|\nabla U\|^2/(4L)}\le e^{d/2}\). The entropy inequality then proves both integrability and \(G_\mu\le4L H(\mu\mid\nu)+2dL\). This direct argument closes the moment input to ULA without assuming an entropy-cost coupling or a transport inequality.



The Gibbs directional score estimate implies \(\int\exp(\|\nabla U\|^2/(4L))d\nu\le2^{d/2}\le e^{d/2}\), by averaging against an auxiliary Gaussian and completing its square. The entropy variational bound \(\int Fd\mu\le\operatorname{KL}(\mu\mid\nu)+\log\int e^Fd\nu\) then yields \(G_0\le4L H_\nu(\mu)+2dL\). The underlying pointwise convex-duality inequality first proves integrability. If \(\alpha\le L\), the coefficient may be enlarged to \(4L^2/\alpha\).



Earlier results used here: \cite[Section~3.27]{LangevinLean2026ULASC}.
\subsection{\texorpdfstring{Convex duality for the entropy integrand}{Convex duality for the entropy integrand}}
\label{report-step-UnadjustedLangevin-LSI-EntropyMoment-lean-1}


Let \(\Phi(a)=a\log a-a+1\), with its continuous value \(\Phi(0)=1\). For \(a>0\), the elementary inequality \(e^u\ge1+u\), applied to \(u=b-\log a\), gives \(e^b\ge a+ab-a\log a\). Rearranging proves \(ab\le\Phi(a)+e^b-1\). The zero-density case is handled separately: its left side is zero, and its right side is \(e^b\ge0\). Thus this entropy/\allowbreak{}exponential Young inequality remains valid at zero Radon-Nikodym density.



For \(a\ge0\), convex duality for \(a\log a-a+1\) gives \(ab\le a\log a-a+e^b\). The nonnegative entropy correction controls the negative part at \(a=0\) and supports integrability before optimization.

\subsection{\texorpdfstring{Integrability from finite KL and an exponential moment}{Integrability from finite KL and an exponential moment}}
\label{report-step-UnadjustedLangevin-LSI-EntropyMoment-lean-2}


Write \(a=d\mu/d\nu\). Finite KL supplies absolute continuity \(\mu\ll\nu\) and integrability of the actual log likelihood; equivalently \(\Phi(a)\) is integrable under \(\nu\). For a measurable nonnegative observable \(f\), the Young inequality gives \(0\le af\le\Phi(a)+e^f\). Both terms on the right are integrable by the finite-KL and exponential-moment hypotheses.



Domination proves \(af\in L^1(\nu)\). The Radon-Nikodym weighted-integrability equivalence then gives \(f\in L^1(\mu)\). In particular, the moment conclusion is derived before any ordinary expectation is manipulated; it is not an extra assumption hidden in the entropy inequality.



Apply the pointwise inequality with the Radon–Nikodym density and the nonnegative observable. Finite KL and an integrable exponential dominate their product, proving the observable is integrable under the changed law.

\subsection{\texorpdfstring{The entropy variational inequality}{The entropy variational inequality}}
\label{report-step-UnadjustedLangevin-LSI-EntropyMoment-lean-3}


Assume \(\int e^f d\nu\le e^c\). First derive \(f\in L^1(\mu)\) with the preceding theorem. Apply the Young inequality to \(a=d\mu/d\nu\) and \(f-c\), and integrate: \(\int(f-c)d\mu\le\int\Phi(a)d\nu+\int e^{f-c}d\nu-1\).



Probability normalization identifies \(\int\Phi(a)d\nu\) with the genuine finite KL divergence and \(\int(f-c)d\mu=\int f d\mu-c\). The MGF hypothesis makes \(\int e^{f-c}d\nu\le1\), canceling the final two terms. Thus \(\int f d\mu\le H(\mu\mid\nu)+c\), together with the integrability witness used in its derivation. The proof explicitly justifies each weighted integral and subtraction.



Subtract a scalar normalization from the observable before applying the inequality. Choosing that normalization as the logarithm of its exponential moment produces the variational entropy bound with the sharp log-moment term.

\subsection{\texorpdfstring{Integrability of the auxiliary radial Gaussian}{Integrability of the auxiliary radial Gaussian}}
\label{report-step-UnadjustedLangevin-LSI-EntropyMoment-lean-4}


For \(b>0\), the real radial Gaussian \(e^{-b\|v\|^2}\) is integrable on every finite-dimensional real Hilbert space. The proof specializes the complex Gaussian Fourier integrability theorem to real positive precision and takes its norm; the complex norm is exactly the displayed real exponential. This supplies the finite integral needed in the Gaussian linearization argument.



A positive Gaussian precision gives integrable radial decay in finite dimension. This supplies the auxiliary Gaussian averaging measure used to turn directional score estimates into a squared-score estimate.

\subsection{\texorpdfstring{The nonnegative radial Gaussian integral}{The nonnegative radial Gaussian integral}}
\label{report-step-UnadjustedLangevin-LSI-EntropyMoment-lean-5}


Let \(b>0\) and let \(d\) be finite. The radial Gaussian \(v\mapsto e^{-b\|v\|^2}\) is nonnegative and integrable. Its nonnegative extended integral therefore equals its ordinary integral, and Gaussian integration gives

\[\int_{\mathbb R^d}e^{-b\|v\|^2}dv
=\left(\frac\pi b\right)^{d/2}<\infty.\]


The expression on the right is strictly positive because \(b>0\). It supplies the finite nonzero normalizing constant for the auxiliary Gaussian average. In dimension zero the empty product is one, so the same formula remains valid. Nonnegativity also permits Tonelli to exchange this Gaussian averaging with the target integral before their final numerical bound is evaluated.

\subsection{\texorpdfstring{Completing the square in an exponentially tilted Gaussian}{Completing the square in an exponentially tilted Gaussian}}
\label{report-step-UnadjustedLangevin-LSI-EntropyMoment-lean-6}


Complete the square in \(-b\|v\|^2+\langle g,v\rangle\): it equals \(\|g\|^2/(4b)-b\|v-g/(2b)\|^2\). Exponentiating separates the constant factor \(e^{\|g\|^2/(4b)}\) from a translated radial Gaussian. Translation invariance of Lebesgue measure and the radial normalization give \(\int^- e^{-b\|v\|^2}e^{\langle g,v\rangle}dv=e^{\|g\|^2/(4b)}(\pi/b)^{d/2}\). The identity is proved for nonnegative integrals before any integrability of a random score is assumed.



Completing the square in the auxiliary variable evaluates its exponentially tilted Gaussian integral. The remaining factor is an exponential of the squared norm of the vector paired with it.

\subsection{\texorpdfstring{From directional moment bounds to a squared-\allowbreak{}norm moment}{From directional moment bounds to a squared-norm moment}}
\label{report-step-UnadjustedLangevin-LSI-EntropyMoment-lean-7}


Suppose \(\mathbb E e^{\langle g,v\rangle}\le e^{L\|v\|^2/2}\) for every direction \(v\). Multiply by \(e^{-L\|v\|^2}\) and integrate in \(v\). Tonelli applies to the nonnegative jointly measurable integrand, so the order of integration can be exchanged without assuming the desired exponential moment in advance.



Gaussian linearization identifies the left side with \((\pi/L)^{d/2}\mathbb E e^{\|g\|^2/(4L)}\). The directional MGF bound controls the right side by \(\int e^{-(L/2)\|v\|^2}dv=(2\pi/L)^{d/2}\). Cancel the positive finite factor \((\pi/L)^{d/2}\) in the extended nonnegative reals to obtain \(\mathbb E e^{\|g\|^2/(4L)}\le2^{d/2}\).



Average the directional exponential bound against the auxiliary Gaussian. Tonelli applies to the positive integrand, and square completion converts its left side into an exponential squared-norm moment.

\subsection{\texorpdfstring{The exponential Gibbs score moment}{The exponential Gibbs score moment}}
\label{report-step-UnadjustedLangevin-LSI-EntropyMoment-lean-8}


The smooth Gibbs score has the previously proved directional MGF bound \(\mathbb E_\nu e^{\langle\nabla U,v\rangle}\le e^{L\|v\|^2/2}\), obtained from translation invariance and first-order smoothness. Apply Gaussian averaging with \(g=\nabla U\). This gives a finite nonnegative exponential integral bounded by \(2^{d/2}\); measurability and this finiteness prove ordinary integrability.



Convert the nonnegative integral back to its real expectation. Since \(\log2\le1\), \(2^{d/2}\le e^{d/2}\), giving the convenient bound \(\mathbb E_\nu e^{\|\nabla U\|^2/(4L)}\le e^{d/2}\). Only differentiability, Lipschitz gradient, positive \(L\), and finite Gibbs normalization are used; neither convexity nor LSI enters this exponential-moment theorem.



Choosing the auxiliary variance to match \(4L\) gives the factor \(2^{d/2}\). Bounding \(\log2\le1\) turns its logarithm into at most \(d/2\), which is the dimensional constant needed in the score bound.



Earlier results used here: \cite[Section~3.26.2]{LangevinLean2026ULASC}; \cite[Section~3.26.3]{LangevinLean2026ULASC}; \cite[Section~3.26.9]{LangevinLean2026ULASC}; \cite[Section~3.27.2]{LangevinLean2026ULASC}.
\subsection{\texorpdfstring{The gradient moment supplied by finite entropy}{The gradient moment supplied by finite entropy}}
\label{report-step-UnadjustedLangevin-LSI-EntropyMoment-lean-9}


Apply the entropy inequality with \(f(x)=\|\nabla U(x)\|^2/(4L)\) and \(c=d/2\). The Gibbs exponential-moment theorem provides both exponential integrability and the MGF bound; finite initial KL provides the absolute-continuity and entropy inputs. Therefore \(\mathbb E_\mu[\|\nabla U\|^2/(4L)]\le H(\mu\mid\nu)+d/2\), with this observable integrable under \(\mu\).



Multiply by \(4L>0\), proving \(G_\mu\le4L H(\mu\mid\nu)+2dL\) and actual square integrability of the score. The proof also derives finite Gibbs partition from its assumed probability normalization. This replaces the previously missing LSI-to-transport route by a stronger direct entropy/\allowbreak{}score estimate.



Use the observable \(\|\nabla U\|^2/(4L)\) in the entropy variational inequality, then multiply by \(4L>0\). This yields integrability and \(G_0\le4LH_\nu(\mu)+2dL\) simultaneously.



Earlier results used here: \cite[Section~3.26.2]{LangevinLean2026ULASC}; \cite[Section~3.26.3]{LangevinLean2026ULASC}.
\subsection{\texorpdfstring{The gradient moment in the compatible LSI normalization}{The gradient moment in the compatible LSI normalization}}
\label{report-step-UnadjustedLangevin-LSI-EntropyMoment-lean-10}


If \(0<\alpha\le L\), then \(4L\le4L^2/\alpha\). Finite KL is nonnegative, so replacing the direct estimate's coefficient \(4L\) by \(4L^2/\alpha\) preserves the inequality. Thus \(G_\mu\le(4L^2/\alpha)H(\mu\mid\nu)+2dL\), the coefficient used in the ULA one-step calculation. Integrability is inherited from the stronger direct theorem. Compatibility \(\alpha\le L\) is proved from LSI in a separate module.



Because KL is nonnegative and \(\alpha\le L\), the coefficient \(4L\) may be enlarged to \(4L^2/\alpha\). The dimensional term stays \(2dL\), giving exactly the form used in one-step forcing absorption.



Earlier results used here: \cite[Section~3.26.3]{LangevinLean2026ULASC}.
\section{\texorpdfstring{Actual Gaussian ULA laws and their frozen-\allowbreak{}drift interpolation}{Actual Gaussian ULA laws and their frozen-drift interpolation}}
\label{proof-module-UnadjustedLangevin-LSI-ULALaw-lean}


Define the frozen-drift law \(\rho_{k,t}=\operatorname{Law}(X_k-t\nabla U(X_k)+\sqrt{2t}Z)\), with independent standard Gaussian \(Z\), and \(\rho_{k+1}=\rho_{k,h}\). The endpoint and probability identities connect the entropy argument to the actual algorithm.



The interpolation from a law \(\mu\) is the law of \(X-t\nabla U(X)+\sqrt{2t}Z\), with \((X,Z)\) distributed according to the product of \(\mu\) and standard Gaussian measure. Thus the drift remains frozen at the starting state throughout the interpolation. At \(t=0\) the law is exactly \(\mu\), and at \(t=h\) it is its next ULA law. Repeating this transition defines \(\rho_k\); all endpoint entropy identifications use these exact law equalities.



Earlier results used here: \cite[Section~3.36]{LangevinLean2026ULASC}.
\subsection{\texorpdfstring{Entropy of the iterated ULA law}{Entropy of the iterated ULA law}}
\label{report-step-UnadjustedLangevin-LSI-ULALaw-lean-1}


Write \(H_k=\operatorname{KL}(\rho_k\mid\nu)\) for the entropy of the actual law after \(k\) Gaussian transitions. This quantity may initially be extended-valued. The finite-entropy induction below proves \(H_k<\infty\) under the theorem assumptions, and only then treats \((H_k)\) as an ordinary real recurrence.



The entropy belongs to the law after \(k\) actual Gaussian transitions. It can be used as an ordinary sequence only when every term is finite. The later induction establishes precisely that property; the initial definition records no hidden finiteness conclusion.

\subsection{\texorpdfstring{The initial entropy}{The initial entropy}}
\label{report-step-UnadjustedLangevin-LSI-ULALaw-lean-2}


Substitute the defining identity \(\rho_0=\mu\) into the entropy definition. This gives \(H_0=\operatorname{KL}(\mu\mid\nu)\), the initial term appearing in the final exponential bound.



At zero no transition has been applied, so \(\rho_0=\mu\). Substitution into KL preserves the full initial divergence, which is the quantity contracted by the exponential in the final theorem.

\subsection{\texorpdfstring{The frozen-\allowbreak{}drift interpolation law}{The frozen-drift interpolation law}}
\label{report-step-UnadjustedLangevin-LSI-ULALaw-lean-3}


Define \(\rho_t=(T_t)_*(\mu\otimes\gamma)\), where \(T_t(x,z)=x-t\nabla U(x)+\sqrt{2t}z\). The product encodes independence. This is an alias of the actual the transition \(\mu\mapsto T_{h\#}(\mu\otimes\gamma)\); its frozen-drift Gaussian interpolation interpretation applies to \(t\ge0\).



The product law means the Gaussian is independent of the initial state. Conditioning on that state gives a Gaussian centered at \(x-t\nabla U(x)\) with covariance \(2tI_d\). Mixture and drift-error calculations later use this exact conditional law.



Earlier results used here: \cite[Section~3.36.7]{LangevinLean2026ULASC}.
\subsection{\texorpdfstring{Iteration of the Gaussian transition}{Iteration of the Gaussian transition}}
\label{report-step-UnadjustedLangevin-LSI-ULALaw-lean-4}


Iterate the actual independent-Gaussian transition at fixed step \(h\): \(\rho_k=P_h^k\mu\). The definition is on measures, while the separate chain-law theorems identify independent iid constructions with these same laws.



Every application of the transition integrates a fresh independent Gaussian against the current law. Thus this measure recursion represents the marginal laws of any independent-noise realization, even when those realizations use different probability spaces.



Earlier results used here: \cite[Section~3.36.7]{LangevinLean2026ULASC}.
\subsection{\texorpdfstring{The interpolation endpoint is the next iterate}{The interpolation endpoint is the next iterate}}
\label{report-step-UnadjustedLangevin-LSI-ULALaw-lean-5}


At elapsed time \(h\), the interpolation from \(\rho_k\) is \(P_h\rho_k\). The preceding successor identity therefore proves \(\rho_{k,h}=\rho_{k+1}\). This is the law equality used to turn a continuous interpolation bound into the discrete entropy step.



The elapsed-time endpoint is the transition at exactly the fixed step \(h\). Therefore a bound on interpolation entropy at that endpoint is a bound on \(H_{k+1}\), rather than on an auxiliary continuous-time process with a different law.

\subsection{\texorpdfstring{Probability preservation at every iterate}{Probability preservation at every iterate}}
\label{report-step-UnadjustedLangevin-LSI-ULALaw-lean-6}


Induct from the initial probability \(\mu\). At the successor rewrite the iterate as one interpolation step at time \(h\), then apply probability preservation. Every algorithm law is therefore a probability before applying KL and LSI identities.



The induction starts from a probability law and uses preservation of mass by the measurable update pushforward. This gives total mass one at every index before probability-specific KL and LSI formulas are used.

\subsection{\texorpdfstring{The interpolation at time zero}{The interpolation at time zero}}
\label{report-step-UnadjustedLangevin-LSI-ULALaw-lean-7}


At \(t=0\), both the drift increment and Gaussian increment vanish, so the update map is first-coordinate projection. Integrating out the standard Gaussian probability returns \(\mu\). This endpoint identity needs no potential regularity.



The update reduces to the first projection because both time coefficients vanish. The Gaussian factor integrates to one, leaving the starting measure exactly, with no limiting argument or density regularity at zero.



Earlier results used here: \cite[Section~3.34.1]{LangevinLean2026ULASC}; \cite[Section~3.36.7]{LangevinLean2026ULASC}.
\subsection{\texorpdfstring{The zeroth iterated law}{The zeroth iterated law}}
\label{report-step-UnadjustedLangevin-LSI-ULALaw-lean-8}


The zeroth iterate of a function is the identity, hence \(\rho_0=\mu\). The recursive definition makes the initial-entropy rewrite exact rather than an asymptotic endpoint argument.



The zeroth iterate of the transition map leaves every measure unchanged. This identity fixes the initial law independently of the step size and of any later smoothing by a positive Gaussian time.

\subsection{\texorpdfstring{The successor law}{The successor law}}
\label{report-step-UnadjustedLangevin-LSI-ULALaw-lean-9}


The successor-iterate identity gives \(\rho_{k+1}=P_h(\rho_k)\). Apply the identity for one additional application of the transition, whose order matches applying one more transition to the current law.



The current law is passed through the same transition once more. This order is consistent with the pathwise recursion, where the next fresh innovation is combined with the state already reached at the current index.



Earlier results used here: \cite[Section~3.36.7]{LangevinLean2026ULASC}.
\subsection{\texorpdfstring{Probability preservation along the interpolation}{Probability preservation along the interpolation}}
\label{report-step-UnadjustedLangevin-LSI-ULALaw-lean-10}


Lipschitzness makes the gradient continuous and the joint update measurable. The product \(\mu\otimes\gamma\) is a probability, and its pushforward has the same total mass. Evaluating the mapped measure on the whole space proves normalization of \(\rho_t\).



The joint update is measurable because the gradient is continuous. A pushforward preserves total mass, and the independent product source has mass one. These two facts establish probability normalization at every interpolation time under the stated assumptions.



Earlier results used here: \cite[Section~3.36.5]{LangevinLean2026ULASC}; \cite[Section~3.36.7]{LangevinLean2026ULASC}.
\section{\texorpdfstring{Construct the positive Gaussian mixture and Gibbs-\allowbreak{}relative density}{Construct the positive Gaussian mixture and Gibbs-relative density}}
\label{proof-module-UnadjustedLangevin-LSI-InterpolationDensity-lean}


For \(t>0\), the transition kernel is \(K_t(x,y)=(4\pi t)^{-d/2}e^{-\|y-x+t\nabla U(x)\|^2/(4t)}\), and the mixture is \(p_t(y)=\int K_t(x,y)d\mu(x)\). Each kernel is positive and bounded by the common peak \((4\pi t)^{-d/2}\). Consequently the mixture is finite and positive everywhere, not merely almost everywhere. With Gibbs density \(J=Z^{-1}e^{-U}\), the relative density is \(r_t=p_t/J\), and \(r_t\nu=p_tdy\) is the actual proposed law.



Earlier results used here: \cite[Section~3.37]{LangevinLean2026ULASC}.
\subsection{\texorpdfstring{Define the explicit Gaussian mixture before taking real values}{Define the explicit Gaussian mixture before taking real values}}
\label{report-step-UnadjustedLangevin-LSI-InterpolationDensity-lean-1}


For the state-dependent center \(c_t(x)=x-t\nabla U(x)\), let

\[K_t(x,y)=C_t\exp\!\left(-\frac{\|y-c_t(x)\|^2}{4t}\right),\qquad q_t^{\mathrm{ext}}(y)=\int K_t(x,y)\,d\mu(x).\]

The kernel \(K_t(x,y)\) is the normalized Gaussian density with mean \(x-t\nabla U(x)\) and covariance \(2tI_d\). Integrating it against \(\mu(dx)\) forms the observation density. Starting with a nonnegative integral permits Tonelli before any finiteness assertion; the following peak bound proves that the mixture is an ordinary finite function at positive time.



The Gaussian density formula is used for \(t>0\). Keeping the integral in \(\mathbb R_{\ge0}\cup\{\infty\}\) makes nonnegativity immediate and preserves a genuine finiteness obligation before taking a real density. The later normalizer bound discharges that obligation at every proposed state.



The mixture is initially a nonnegative integral, so Tonelli applies without an integrability hypothesis. Only after a uniform finite upper bound is proved is it identified with an ordinary real-valued function.



Earlier results used here: \cite[Section~3.37.13]{LangevinLean2026ULASC}.
\subsection{\texorpdfstring{Apply measurable parameter integration to the mixture}{Apply measurable parameter integration to the mixture}}
\label{report-step-UnadjustedLangevin-LSI-InterpolationDensity-lean-2}


If \(\nabla U\) is Lipschitz, its state-dependent center \(c_t\) is continuous and measurable. The joint Gaussian integrand \((x,y)\mapsto K_t(x,y)\) is therefore measurable: subtraction, scalar multiplication, squared norm and exponential preserve the required measurability.



The kernel is jointly measurable in its starting and observation variables. Integrating the starting variable against the probability law therefore gives a measurable function of the observation.



Earlier results used here: \cite[Section~3.37.13]{LangevinLean2026ULASC}.
\subsection{\texorpdfstring{Prove the Gaussian mixture is the density of the actual update}{Prove the Gaussian mixture is the density of the actual update}}
\label{report-step-UnadjustedLangevin-LSI-InterpolationDensity-lean-3}


For every \(t>0\),

\[\rho_t=(Y_t)_\#(\mu\otimes\mathcal N(0,I))=q_t^{\mathrm{ext}}\,dy.\]

To prove equality of measures, fix a measurable set \(A\). The pushforward formula and Tonelli express the actual proposal probability as

\[\rho_t(A)=\int \mathbb P\bigl(c_t(x)+\sqrt{2t}Z\in A\bigr)\,d\mu(x).\]\[\rho_t(A)=\int_A\left(\int K_t(x,y)\,d\mu(x)\right)dy=\int_Aq_t^{\mathrm{ext}}(y)\,dy.\]

The proof mirrors these three steps with the defining preimage formula for a measurable pushforward, the section-integral formula for a product measure, the fixed-state density identity, and Tonelli under the restricted volume measure on \(A\). It assumes no density for \(\mu\).



For a nonnegative test, Fubini computes the pushforward expectation by first integrating the conditional Gaussian. The resulting density is the stated mixture, identifying it with the law of the actual update.



Earlier results used here: \cite[Section~3.34.1]{LangevinLean2026ULASC}; \cite[Section~3.36.5]{LangevinLean2026ULASC}; \cite[Section~3.36.7]{LangevinLean2026ULASC}; \cite[Section~3.37.13]{LangevinLean2026ULASC}; \cite[Section~3.37.16]{LangevinLean2026ULASC}.
\subsection{\texorpdfstring{Use positivity of every Gaussian component}{Use positivity of every Gaussian component}}
\label{report-step-UnadjustedLangevin-LSI-InterpolationDensity-lean-4}


At every positive time and every point \(y\), the mixture density is strictly positive:

\[q_t^{\mathrm{ext}}(y)>0.\]

For every fixed observation, the kernel is strictly positive at every starting point. Integrating this positive function against a probability measure gives a strictly positive mixture value.



Earlier results used here: \cite[Section~3.37.13]{LangevinLean2026ULASC}; \cite[Section~3.37.14]{LangevinLean2026ULASC}.
\subsection{\texorpdfstring{Bound every mixture value by the common Gaussian peak}{Bound every mixture value by the common Gaussian peak}}
\label{report-step-UnadjustedLangevin-LSI-InterpolationDensity-lean-5}


For \(t>0\), the Gaussian exponent is nonpositive, hence

\[0\le K_t(x,y)\le C_t,\qquad q_t^{\mathrm{ext}}(y)\le C_t.\]

The proof establishes the exponential bound with the fact that an exponential with nonpositive exponent is at most one: the numerator \(-\|y-c_t(x)\|^2\) is nonpositive and the denominator \(4t\) is positive. Multiplication by the nonnegative normalizer preserves the inequality, as does a nonnegative scalar.



monotonicity of the nonnegative integral integrates this pointwise bound. The integral of the constant \(C_t\) equals \(C_t\mu(E)=C_t\) because \(\mu\) is a probability measure. This theorem requires no regularity of the gradient: its bound concerns the nonnegative integrand at every point, and the upper nonnegative integral remains monotone even before measurability is established.



The exponential factor is at most one, independently of the Gaussian center. Hence averaging over any starting probability retains the same peak bound, without requiring an initial moment.



Earlier results used here: \cite[Section~3.37.12]{LangevinLean2026ULASC}; \cite[Section~3.37.13]{LangevinLean2026ULASC}; \cite[Section~3.37.14]{LangevinLean2026ULASC}.
\subsection{\texorpdfstring{Derive pointwise finiteness before real conversion}{Derive pointwise finiteness before real conversion}}
\label{report-step-UnadjustedLangevin-LSI-InterpolationDensity-lean-6}


Fix \(t>0\) and let \(C_t=(4\pi t)^{-d/2}<\infty\). Every Gaussian component satisfies \(0\le K_t(x,y)\le C_t\), independently of its center. Because \(\mu\) is a probability, integration in the starting point gives

\[0\le\int K_t(x,y)d\mu(x)\le\int C_t d\mu=C_t<\infty.\]


Thus the initially nonnegative-integral definition of the mixture is finite at every observation \(y\). It has an ordinary real value pointwise, with the same peak bound, even if \(\mu\) has no density or position moments. The positive-time condition is essential: at zero the Gaussian peak need not have a finite limiting bound.

\subsection{\texorpdfstring{Choose the finite real representative of the mixture}{Choose the finite real representative of the mixture}}
\label{report-step-UnadjustedLangevin-LSI-InterpolationDensity-lean-7}


For \(t>0\), the preceding peak bound makes the mixture finite at every observation. We therefore use its ordinary real-valued representative

\[p_t(y)=\int (4\pi t)^{-d/2}\exp\!\left(-\frac{\|y-x+t\nabla U(x)\|^2}{4t}\right)d\mu(x).\]

This nonnegative integral is at most \((4\pi t)^{-d/2}\) and is strictly positive because every Gaussian component is positive and \(\mu\) is a probability. Its ordinary value and nonnegative extended value agree everywhere. The positive-time law identity, regularity statements, and logarithmic derivatives below all use this same finite representative.

\subsection{\texorpdfstring{Identify the actual law with the finite real density}{Identify the actual law with the finite real density}}
\label{report-step-UnadjustedLangevin-LSI-InterpolationDensity-lean-8}


For \(t>0\), write \(q_t(y)=\int K_t(x,y)d\mu(x)\), which is finite everywhere by the preceding peak bound. For every measurable set \(A\), the actual proposal law satisfies

\[\begin{aligned}
\rho_t(A)&=\int\left(\int_AK_t(x,y)dy\right)d\mu(x)\\
&=\int_A\left(\int K_t(x,y)d\mu(x)\right)dy
=\int_Aq_t(y)dy.
\end{aligned}\]


The first identity conditions on the initial point and uses the actual affine Gaussian proposal density. The second is Tonelli for the nonnegative kernel. Its finite pointwise mixture is precisely the ordinary density defined above, so no change of law accompanies passage to real-valued density notation. This identification remains valid for singular initial laws.

\subsection{\texorpdfstring{Deduce absolute continuity for every positive time}{Deduce absolute continuity for every positive time}}
\label{report-step-UnadjustedLangevin-LSI-InterpolationDensity-lean-9}


The law identity \(\rho_t=q_t^{\mathrm{ext}}dy\) immediately implies \(\rho_t\ll dy\): every Lebesgue-null set has zero integral against the density and thus zero proposal probability.



The actual interpolation law has density \(q_t\) with respect to Lebesgue measure. Hence a Lebesgue-null set has zero interpolation probability, proving absolute continuity at positive time. This holds even for singular initial laws. Finite KL with respect to Gibbs and entropy behavior at zero require the separate score and endpoint arguments.



The actual proposed law has a Lebesgue density, so every Lebesgue-null set has proposed mass zero. This absolute continuity holds for arbitrary initial laws once the elapsed Gaussian time is positive.

\subsection{\texorpdfstring{Convert the nonnegative mixture to its ordinary real integral}{Convert the nonnegative mixture to its ordinary real integral}}
\label{report-step-UnadjustedLangevin-LSI-InterpolationDensity-lean-10}


The real density has the expected ordinary integral formula:

\[q_t(y)=\int C_t\exp\!\left(-\frac{\|y-c_t(x)\|^2}{4t}\right)d\mu(x).\]

The nonnegative Gaussian kernel is measurable and bounded by the integrable constant \(C_t\) under the mixing probability. Its ordinary integral therefore exists and is nonnegative. The finite nonnegative mixture integral equals this same ordinary integral, giving the claimed explicit density formula.



The component kernel is nonnegative and bounded by a finite constant under the mixing probability. Its ordinary integral therefore exists and equals the finite value of the initial nonnegative mixture integral.



Earlier results used here: \cite[Section~3.37.12]{LangevinLean2026ULASC}; \cite[Section~3.37.13]{LangevinLean2026ULASC}; \cite[Section~3.37.14]{LangevinLean2026ULASC}.
\subsection{\texorpdfstring{Pass spatial continuity through the probability mixture}{Pass spatial continuity through the probability mixture}}
\label{report-step-UnadjustedLangevin-LSI-InterpolationDensity-lean-11}


At a fixed positive time, \(q_t:E\to\mathbb R\) is continuous. For every fixed initial state \(x\), the Gaussian kernel \(y\mapsto K_t(x,y)\) is continuous, and its absolute value is bounded by the same constant \(C_t\), independent of \(x\) and \(y\).



The proof first replaces the explicit Gaussian-mixture Lebesgue density by the ordinary integral formula from the previous theorem. It applies continuity under an integral with a common integrable majorant with constant majorant \(x\mapsto C_t\). Measurability in \(x\) follows from gradient continuity. The nonpositive Gaussian exponent proves \(\|K_t(x,y)\|=K_t(x,y)\le C_t\); probability normalization makes this majorant integrable. Pointwise continuity in \(y\) is then assembled from subtraction, norm, square and exponential.



At a fixed positive time the component kernel is continuous in the observation and bounded by one common peak. Dominated convergence passes that continuity through the mixture integral.



Earlier results used here: \cite[Section~3.37.12]{LangevinLean2026ULASC}; \cite[Section~3.37.14]{LangevinLean2026ULASC}.
\subsection{\texorpdfstring{Derive finite partition from target normalization}{Derive finite partition from target normalization}}
\label{report-step-UnadjustedLangevin-LSI-InterpolationDensity-lean-12}


Recall \(\nu=Z^{-1}(e^{-U}dy)\) with the extended partition function \(Z=\int e^{-U}dy\). If this measure is a probability measure, then \(Z\ne\infty\).



No strong convexity or LSI estimate is used. This matters because LSI already includes probability normalization of its target, so its finite partition requirement can be recovered from that normalization rather than imported through a stronger convexity hypothesis.



If the normalized Gibbs measure is a probability, its partition must be finite and nonzero under the displayed assumptions. This permits its density to be used as an ordinary positive denominator.



Earlier results used here: \cite[Section~3.26.2]{LangevinLean2026ULASC}; \cite[Section~3.26.3]{LangevinLean2026ULASC}.
\subsection{\texorpdfstring{Write the explicit density relative to Gibbs}{Write the explicit density relative to Gibbs}}
\label{report-step-UnadjustedLangevin-LSI-InterpolationDensity-lean-13}


The normalized Gibbs target has Lebesgue density \(Z^{-1}e^{-U(y)}\). Therefore the relative density corresponding to the mixture \(q_t(y)dy\) is

\[r_t(y)=Z e^{U(y)}q_t(y).\]

The relative density is the ordinary product \(r_t(y)=Ze^{U(y)}q_t(y)\). The positive finite partition and mixture peak bounds make all factors finite at positive time. Multiplication by the Gibbs density cancels \(Ze^{U(y)}\), producing the actual Lebesgue density \(q_t\).



The definition makes the coordinate conversion explicit: the ratio to the target is obtained by multiplying the Lebesgue density by the reciprocal target density. It does not package an entropy derivative, convergence bound or weak PDE as part of the meaning of a density.



The Gibbs Lebesgue density is \(J(y)=Z^{-1}e^{-U(y)}\), so the relative density is \(r_t(y)=Z e^{U(y)}p_t(y)\). Multiplying by \(J\) recovers \(p_t\) pointwise.



Earlier results used here: \cite[Section~3.26.2]{LangevinLean2026ULASC}.
\subsection{\texorpdfstring{Multiply measurable factors to obtain relative-\allowbreak{}density measurability}{Multiply measurable factors to obtain relative-density measurability}}
\label{report-step-UnadjustedLangevin-LSI-InterpolationDensity-lean-14}


Write the relative density as \(r_t(y)=Ze^{U(y)}p_t(y)\). The mixture \(p_t\) is measurable by parameter integration; differentiability of \(U\) gives measurability of \(e^U\); and the partition \(Z\) is a fixed finite scalar. Products of these measurable factors are measurable. Thus the explicit relative density is measurable, as required to define its weighted target measure. This statement is a measurability conclusion; spatial differentiation and finite entropy are established separately.

\subsection{\texorpdfstring{Obtain relative-\allowbreak{}density continuity at positive time}{Obtain relative-density continuity at positive time}}
\label{report-step-UnadjustedLangevin-LSI-InterpolationDensity-lean-15}


For \(t>0\), spatial continuity of \(q_t\) and continuity of \(e^U\) give continuity of \(r_t(y)=Z e^{U(y)}q_t(y)\).



Continuity of \(p_t\), the potential, and the exponential yields continuity of \(r_t\). The Gibbs density never vanishes, so division introduces no discontinuity.

\subsection{\texorpdfstring{Retain all finiteness guards when proving strict positivity}{Retain all finiteness guards when proving strict positivity}}
\label{report-step-UnadjustedLangevin-LSI-InterpolationDensity-lean-16}


When the target Gibbs measure is a probability measure and \(t>0\), all three factors of \(r_t(y)=Z e^{U(y)}q_t(y)\) are strictly positive. The partition is positive by positivity of the Gibbs partition and finite by finiteness of the partition forced by target probability normalization. The mixture is positive by positivity of every Gaussian-mixture value and finite by the common finite Gaussian peak bound.



The mixture is positive and finite, the partition is positive and finite, and \(e^{U(y)}\) is positive. These facts jointly establish strict positivity of the actual real relative density everywhere.



Earlier results used here: \cite[Section~3.26.3]{LangevinLean2026ULASC}; \cite[Section~3.26.9]{LangevinLean2026ULASC}.
\subsection{\texorpdfstring{Cancel the target density and identify the exact relative law}{Cancel the target density and identify the exact relative law}}
\label{report-step-UnadjustedLangevin-LSI-InterpolationDensity-lean-17}


Fix \(t>0\) and assume the Gibbs partition satisfies \(0<Z<\infty\), as proved under the target normalization hypotheses. Let

\[J(y)=Z^{-1}e^{-U(y)},\qquad r_t(y)=Ze^{U(y)}q_t(y).\]


The mixture \(q_t\) is finite and measurable, and the Gibbs density is strictly positive. Pointwise multiplication therefore gives the ordinary identity

\[r_t(y)J(y)=\bigl(Ze^{U(y)}q_t(y)\bigr)\bigl(Z^{-1}e^{-U(y)}\bigr)=q_t(y).\]


Integrating over any measurable set identifies the measures:

\[r_t\nu=q_tdy=\rho_t.\]


The last equality is the already proved actual Gaussian-mixture law identity. Hence the constructed relative density represents the actual interpolation, rather than an auxiliary law. Positivity, spatial regularity, time differentiation, and finite entropy remain separate properties established in the subsequent analysis.



Earlier results used here: \cite[Section~3.26.1]{LangevinLean2026ULASC}; \cite[Section~3.26.2]{LangevinLean2026ULASC}; \cite[Section~3.26.3]{LangevinLean2026ULASC}; \cite[Section~3.26.9]{LangevinLean2026ULASC}.
\subsection{\texorpdfstring{Transfer normalization from the actual proposal law}{Transfer normalization from the actual proposal law}}
\label{report-step-UnadjustedLangevin-LSI-InterpolationDensity-lean-18}


The measure induced by the constructed relative density is a probability measure:

\[\int r_t\,d\nu=1\qquad(t>0).\]

Integrating the identity \(r_t\nu=\rho_t\) over the entire state space gives \(\int r_t\,d\nu=1\), because the Gaussian update is a probability law. Thus the explicitly constructed density is normalized; no separate assumption about its mass is required.



Integrating the law identity on the whole space gives \(\int r_t\,d\nu=1\). Normalization is therefore inherited from the Gaussian transition rather than proved by a separate algebraic manipulation of a possibly infinite density.



Earlier results used here: \cite[Section~3.26.3]{LangevinLean2026ULASC}.
\section{\texorpdfstring{Gaussian displacement and frozen-\allowbreak{}gradient error energy}{Gaussian displacement and frozen-gradient error energy}}
\label{proof-module-UnadjustedLangevin-LSI-GaussianDrift-lean}


For \(Y_t=X-t\nabla U(X)+\sqrt{2t}Z\), independent Gaussian noise makes the cross term vanish, giving \(\mathbb E\|Y_t-X\|^2=t^2\mathbb E\|\nabla U(X)\|^2+2td\). Gradient Lipschitzness bounds the corresponding frozen-gradient error by \(L^2\) times this moment.



Put \(G_0=\mathbb E\|\nabla U(X)\|^2\) and let \(Z\) be an independent standard Gaussian. The displacement is \(-t\nabla U(X)+\sqrt{2t}Z\), whose squared mean is \(t^2G_0+2td\): the mixed term vanishes by centering and independence. Gradient Lipschitzness multiplies this energy by at most \(L^2\). This controls drift error without assuming \(\mathbb E\|X\|^2<\infty\), since the absolute initial position cancels from the displacement.



Earlier results used here: \cite[Section~3.36]{LangevinLean2026ULASC}.
\subsection{\texorpdfstring{The standard Gaussian second moment}{The standard Gaussian second moment}}
\label{report-step-UnadjustedLangevin-LSI-GaussianDrift-lean-1}


Choose a finite orthonormal basis and use \(\|Z\|^2=\sum_i\langle e_i,Z\rangle^2\). Each standard Gaussian coordinate has second moment one, proved using its dual variance formula and zero mean. Gaussian \(L^2\) membership justifies integrating the finite sum, giving \(\mathbb E\|Z\|^2=d\), also when \(d=0\).



The Gaussian covariance is the identity, so every orthonormal coordinate contributes one to the energy. Summing the finite coordinate squares gives \(d\), without any trace normalization factor. This is the source of the dimension term in the drift-error estimate.

\subsection{\texorpdfstring{The energy of an affine Gaussian displacement}{The energy of an affine Gaussian displacement}}
\label{report-step-UnadjustedLangevin-LSI-GaussianDrift-lean-2}


Expand \(\|ag(X)+bZ\|^2=a^2\|g(X)\|^2+2ab\langle g(X),Z\rangle+b^2\|Z\|^2\). The \(L^2\) hypothesis on \(g\) and Gaussian \(L^2\) membership prove all three terms integrable on the independent product.



The product bilinear-integral theorem factors the cross expectation as \(\langle\mathbb Eg(X),\mathbb EZ\rangle=0\). The product marginals identify the first energy, and the second moment \(\mathbb E\|Z\|^2=d\) identifies the noise energy. Hence the integral is \(a^2\mathbb E\|g\|^2+b^2d\).



The mixed term is integrable by Cauchy–Schwarz and independence factors its expectation. Since the Gaussian mean is zero, that expectation vanishes even when the drift has nonzero mean. The drift and Gaussian square terms then contribute their exact separate energies.

\subsection{\texorpdfstring{Subtracting the initial state from the proposal}{Subtracting the initial state from the proposal}}
\label{report-step-UnadjustedLangevin-LSI-GaussianDrift-lean-3}


Subtract \(x\) from \(T_t(x,z)=x-t\nabla U(x)+\sqrt{2t}z\). Additive cancellation leaves \(-t\nabla U(x)+\sqrt{2t}z\). This algebraic identity shows why a state second moment is unnecessary when only displacement energy is estimated.



The starting position cancels algebraically before integration. Thus its absolute second moment is irrelevant to this displacement identity; the only initial moment required is that of its gradient. This distinction permits finite-KL initial laws without an added position-moment assumption.



Earlier results used here: \cite[Section~3.34.1]{LangevinLean2026ULASC}.
\subsection{\texorpdfstring{The exact frozen-\allowbreak{}drift displacement energy}{The exact frozen-drift displacement energy}}
\label{report-step-UnadjustedLangevin-LSI-GaussianDrift-lean-4}


Apply the affine Gaussian energy formula with \(a=-t\), \(b=\sqrt{2t}\), and \(g=\nabla U\). For \(t\ge0\), \(b^2=2t\), yielding \(\mathbb E\|T_t-X\|^2=t^2\mathbb E\|\nabla U(X)\|^2+2td\). Independence and centering have already discharged the cross term.



With \(a=-t\) and \(b=\sqrt{2t}\), the squared coefficients are \(t^2\) and \(2t\). The exact energy identity therefore has no extra factor from bounding the cross term, unlike a triangle-inequality estimate.



Earlier results used here: \cite[Section~3.34.1]{LangevinLean2026ULASC}.
\subsection{\texorpdfstring{The frozen-\allowbreak{}gradient error energy}{The frozen-gradient error energy}}
\label{report-step-UnadjustedLangevin-LSI-GaussianDrift-lean-5}


Lipschitzness gives \(\|\nabla U(T_t)-\nabla U(X)\|^2\le L^2\|T_t-X\|^2\) pointwise. Convert integrability of the squared gradient to \(L^2\) membership, so the displacement has an integrable squared norm.



Continuity of the proposal and gradient proves measurability of the squared discrepancy. Domination by the integrable displacement energy proves its integrability; monotonicity of the integral and the exact displacement moment then give \(L^2(t^2G_0+2td)\). This is a joint-law estimate; a state-space conditional drift is constructed separately.



Apply gradient Lipschitzness pointwise to the actual state displacement, then square and integrate. The resulting energy is \(L^2(t^2G_0+2td)\), retaining the sharp independent-Gaussian constant in the dimensional term.



Earlier results used here: \cite[Section~3.34.1]{LangevinLean2026ULASC}.
\subsection{\texorpdfstring{Square-\allowbreak{}integrability of the displacement}{Square-integrability of the displacement}}
\label{report-step-UnadjustedLangevin-LSI-GaussianDrift-lean-6}


Pull the square-integrable initial gradient to the first product coordinate and the square-integrable Gaussian identity to the second. Their scalar multiples and sum lie in \(L^2\). The displacement identity rewrites this sum as \(T_t(X,Z)-X\).



The drift term is square-integrable by the initial score hypothesis and the Gaussian term by its second moment. Their finite linear combination is consequently square-integrable, which validates the preceding exact norm-square calculation.



Earlier results used here: \cite[Section~3.34.1]{LangevinLean2026ULASC}.
\section{\texorpdfstring{Construct the conditional Gaussian drift and transfer its energy}{Construct the conditional Gaussian drift and transfer its energy}}
\label{proof-module-UnadjustedLangevin-LSI-ConditionalDrift-lean}


The actual frozen Gaussian interpolation is \(Y_t=X-t\nabla U(X)+\sqrt{2t}Z\), where \((X,Z)\) has the independent law \(P=\mu\otimes\mathcal N(0,I)\). Write \(D_t=\nabla U(Y_t)-\nabla U(X)\). This module constructs a measurable field \(e_t\) satisfying \(e_t(Y_t)=\mathbb E[D_t\mid Y_t]\), proves its flux identities, and derives its squared-energy bound by conditional Jensen. Thus the state-dependent error is a consequence of the joint Gaussian construction. The entropy weak equation is a separate analytic obligation.



Let \((X,Y_t)\) have its actual frozen-Gaussian joint law and set \(e(X,Y_t)=\nabla U(Y_t)-\nabla U(X)\). The conditional drift discrepancy \(\bar e_t(y)=\mathbb E[e(X,Y_t)\mid Y_t=y]\) is defined up to the law of \(Y_t\). Its defining integral identity transfers every admissible observed test from the joint space to the state space. Conditional Jensen gives \(\int\|\bar e_t\|^2d\rho_t\le\mathbb E\|e(X,Y_t)\|^2\); hence conditioning cannot increase the error energy.

\subsection{\texorpdfstring{Construct a conditional mean on the observed state space}{Construct a conditional mean on the observed state space}}
\label{report-step-UnadjustedLangevin-LSI-ConditionalDrift-lean-1}


Let \(Y:\Omega\to S\) be measurable, let \(P\) be a probability measure, and let the Banach-valued vector \(D\) satisfy \(\mathbb E_P\|D\|^2<\infty\). The theorem constructs a strongly measurable \(e:S\to F\) such that

\[e(Y)=\mathbb E_P[D\mid\sigma(Y)]\quad P\text{-almost surely},\qquad \int_S\|e(y)\|^2\,d(Y_\#P)(y)\le\int_\Omega\|D\|^2\,dP.\]

First take the square-integrable conditional expectation \(\mathbb E[D\mid\sigma(Y)]\). Doob–Dynkin factorization writes this observed random vector as a strongly measurable vector field \(e\) evaluated at \(Y\). Conditional Jensen gives \(\int\|e(Y)\|^2dP\le\int\|D\|^2dP\), and the pushforward integral formula identifies its left side with \(\int\|e\|^2d(Y_\#P)\). The field is determined only up to observed-law null sets, which leave its energy and all subsequent test identities unchanged.

\subsection{\texorpdfstring{Identify the vector flux on every measurable set}{Identify the vector flux on every measurable set}}
\label{report-step-UnadjustedLangevin-LSI-ConditionalDrift-lean-2}


Suppose a strongly measurable field \(e\) satisfies the conditional-mean identity \(e(Y)=\mathbb E[D\mid\sigma(Y)]\), with \(D\) integrable. For every measurable \(A\subseteq S\),

\[\int_A e(y)\,d(Y_\#P)(y)=\int_{Y^{-1}(A)}D(\omega)\,dP(\omega).\]

The first step, change of variables for an integral over an observed measurable set, pulls the left integral back to \(\int_{Y^{-1}(A)}e(Y)dP\). The almost-sure conditional-mean identity changes its integrand to the conditional expectation. Because \(Y^{-1}(A)\) belongs to the pulled-back sigma algebra, the defining set-integral identity for conditional expectation replaces that expectation by \(D\) in the set integral. The probability assumption guarantees the finite-measure hypotheses for conditional expectation. This proves a vector-valued identification of the flux, extending beyond the energy estimate.



For any measurable observation set \(A\), the conditional mean satisfies \(\int_A\bar e(y)d\rho(y)=\mathbb E[1_{\{Y\in A\}}e]\). This is the exact setwise identity from which test-function identities are extended.

\subsection{\texorpdfstring{Transfer integrable inner-\allowbreak{}product tests through the conditional mean}{Transfer integrable inner-product tests through the conditional mean}}
\label{report-step-UnadjustedLangevin-LSI-ConditionalDrift-lean-3}


Let \(\varphi:S\to F\) be strongly measurable and suppose \(\langle D,\varphi(Y)\rangle\) is integrable. Then

\[\int_S\langle e(y),\varphi(y)\rangle\,d(Y_\#P)(y)=\int_\Omega\langle D,\varphi(Y)\rangle\,dP.\]\[\mathbb E[\langle D,\varphi(Y)\rangle\mid Y]=\langle\mathbb E[D\mid Y],\varphi(Y)\rangle.\]

Begin with simple observed tests, where the identity follows from the setwise formula. Approximation and the stated integrable pairing bounds then extend it to the admissible vector tests, including the later score pairing.

\subsection{\texorpdfstring{Specify the discrepancy on the actual independent joint law}{Specify the discrepancy on the actual independent joint law}}
\label{report-step-UnadjustedLangevin-LSI-ConditionalDrift-lean-4}


For \(p=(x,z)\), define

\[D_t(p)=\nabla U\!\left(x-t\nabla U(x)+\sqrt{2t}\,z\right)-\nabla U(x).\]

This is exactly \(\nabla U(Y_t)-\nabla U(X)\): mala proposal supplies the interpolated state and the two gradient terms are subtracted in the displayed order. Its sign agrees with the additional relative-density drift \(e_t(y)=\nabla U(y)-\mathbb E[\nabla U(X)\mid Y_t=y]\). No conditional expectation, density, or error bound is part of this definition. For \(t\ge0\), the later energy theorem uses the actual independent Gaussian product law.



The discrepancy compares the frozen initial gradient with the gradient at the actual proposed state. Defining it on the original independent pair space preserves the sharp displacement energy before any conditioning is applied.



Earlier results used here: \cite[Section~3.34.1]{LangevinLean2026ULASC}.
\subsection{\texorpdfstring{Record precisely what is observed at interpolation time}{Record precisely what is observed at interpolation time}}
\label{report-step-UnadjustedLangevin-LSI-ConditionalDrift-lean-5}


The observed variable is \(Y_t(x,z)=x-t\nabla U(x)+\sqrt{2t}z\). Its sigma-algebra \(\sigma(Y_t)\) is generated by inverse images of measurable state-space sets. Conditioning on this sigma-algebra averages out distinctions between initial-noise pairs that yield the same observed state.



The abbreviation \(\sigma(Y_t)\) makes the conditional expectation statements refer to this exact observation. Conditioning retains the interpolated position and averages out information about the starting point and Gaussian innovation that the position alone does not determine. Lipschitzness of the gradient later proves that this pulled-back sigma algebra lies inside the ambient measurable structure.



The observation is the proposed endpoint \(Y_t\), so the conditional quantity may depend only on that state. Information about the particular initial component is averaged out rather than retained in the state-space drift.



Earlier results used here: \cite[Section~3.34.1]{LangevinLean2026ULASC}.
\subsection{\texorpdfstring{Construct the actual interpolation error and prove its energy bound}{Construct the actual interpolation error and prove its energy bound}}
\label{report-step-UnadjustedLangevin-LSI-ConditionalDrift-lean-6}


Assume \(\nabla U\) is \(L\)-Lipschitz and \(G_0=\int\|\nabla U\|^2d\mu<\infty\). For each \(t\ge0\), the theorem constructs a strongly measurable \(e_t:E\to E\) with

\[e_t(Y_t)=\mathbb E[D_t\mid Y_t],\qquad \int\|e_t(y)\|^2\,d\rho_t(y)\le t^2L^2G_0+2tnL^2,\]\[\rho_t=(Y_t)_\#(\mu\otimes\mathcal N(0,I)),\qquad n=\dim E.\]

The general conditional-mean theorem now constructs err. Its pushforward measure is by definition the actual the frozen-drift interpolation law, so the field is integrable on the required marginal. Conditional Jensen followed by the joint Gaussian estimate gives the result; the final ring step expands \(L^2(t^2G_0+2tn)\) into the exact coefficient order used by the entropy interface. Neither the error field nor its numerical energy estimate is a hypothesis.



Conditional Jensen gives \(\|\mathbb E[e\mid Y_t]\|^2\le\mathbb E[\|e\|^2\mid Y_t]\). Integrating and applying the sharp Gaussian displacement identity yields the same \(L^2(t^2G_0+2td)\) upper bound for the observed field.



Earlier results used here: \cite[Section~3.34.1]{LangevinLean2026ULASC}; \cite[Section~3.36.5]{LangevinLean2026ULASC}.
\subsection{\texorpdfstring{Express the constructed error energy in relative-\allowbreak{}density coordinates}{Express the constructed error energy in relative-density coordinates}}
\label{report-step-UnadjustedLangevin-LSI-ConditionalDrift-lean-7}


Suppose a measurable nonnegative \(r\) identifies the actual interpolation law as \(\rho_t=r\nu\). The preceding construction gives a field \(e_t\) satisfying both the conditional-mean identity and

\[\int r(y)\|e_t(y)\|^2\,d\nu(y)\le t^2L^2G_0+2tnL^2.\]

After constructing the state-space error field, use the law identity \(\rho_t=r_t\nu\). Changing measure gives \(\int\|e_t\|^2d\rho_t=\int r_t\|e_t\|^2d\nu\), and the same formula for absolute integrals transfers integrability. Thus the conditional-Jensen energy bound is available in the weighted reference measure required by entropy production.



Since \(\rho_t=r_t\nu\), changing density rewrites the state-space error energy as \(\int r_t\|\bar e_t\|^2d\nu\). This is exactly the forcing energy required by the relative-density weak equation.



Earlier results used here: \cite[Section~3.34.1]{LangevinLean2026ULASC}.
\section{\texorpdfstring{Construct spatial smoothness, finite Fisher information and finite actual entropy}{Construct spatial smoothness, finite Fisher information and finite actual entropy}}
\label{proof-module-UnadjustedLangevin-LSI-GaussianEntropy-lean}


For \(t>0\), the frozen Gaussian update \(Y_t=X-t\nabla U(X)+\sqrt{2t}Z\) has its explicit positive mixture density \(p_t\). Differentiating the Gaussian kernel under the mixing integral proves spatial \(C^1\) regularity. The posterior component-score formula and Jensen give \(\int\|\nabla\log p_t\|^2d\rho_t\le d/(2t)\). If the initial gradient energy is finite, the propagated gradient and this Lebesgue score are square-integrable, making the relative score \(\nabla\log p_t+\nabla U\) square-integrable as well. LSI then gives finite KL of the actual output law. No initial density or position second moment is assumed.



At positive time, Gaussian polynomial decay uniformly bounds the component derivatives, so spatial differentiation passes through the mixture even without an initial position moment. The posterior mixture formula gives \(\nabla\log p_t(y)=\mathbb E[-(Y_t-X+t\nabla U(X))/(2t)\mid Y_t=y]\). Jensen and the exact Gaussian second moment yield \(\int\|\nabla\log p_t\|^2d\rho_t\le d/(2t)\). Adding the potential score gives finite relative Fisher information, and the target LSI then gives finite actual KL.

\subsection{\texorpdfstring{Bound the quadratic Gaussian envelope uniformly}{Bound the quadratic Gaussian envelope uniformly}}
\label{report-step-UnadjustedLangevin-LSI-GaussianEntropy-lean-1}


For \(a>0\), put \(u=a\|z\|^2\). The scalar inequality \(u e^{-u}\le e^{-1}\le1\) gives

\[\|z\|^2e^{-a\|z\|^2}\le a^{-1}.\]

For \(a>0\), the function \(s\mapsto s^2e^{-as^2}\) is bounded on \([0,\infty)\). This bounds quadratic factors times Gaussian decay independently of the component's center, enabling mixture differentiation without a position moment.

\subsection{\texorpdfstring{Bound the first Gaussian envelope without an initial moment}{Bound the first Gaussian envelope without an initial moment}}
\label{report-step-UnadjustedLangevin-LSI-GaussianEntropy-lean-2}


The elementary estimate \(r\le1+r^2\), together with \(e^{-ar^2}\le1\) and the quadratic envelope bound, gives

\[\|z\|e^{-a\|z\|^2}\le1+a^{-1}.\]

The first derivative introduces one displacement factor. Reserving part of the Gaussian exponential bounds that factor uniformly, leaving a constant integrable against any mixing probability.

\subsection{\texorpdfstring{Specify the real Gaussian kernel and its derivative map}{Specify the real Gaussian kernel and its derivative map}}
\label{report-step-UnadjustedLangevin-LSI-GaussianEntropy-lean-3}


With precision \(a\), center \(c\), and the already normalized Gaussian constant \(N_a\), define

\[K_a(c,y)=N_a e^{-a\|y-c\|^2},\qquad D_yK_a(c,y)[v]=-2aK_a(c,y)\langle y-c,v\rangle.\]

the ordinary Gaussian component kernel is the ordinary real density. the spatial derivative of the Gaussian kernel is its proposed Fréchet derivative as a continuous linear functional: the scalar the Gaussian kernel times the coefficient \(-2a\) multiplies the functional \(v\mapsto\langle y-c,v\rangle\). The precision parameter keeps this helper reusable; the ULA specialization is \(a=1/(4t)\). The next theorem proves that this proposed derivative is the actual derivative.



The derivative map is a continuous linear functional in the observation direction, obtained by pairing with the centered displacement. Defining this functional explicitly fixes the object later integrated over the initial law.



Earlier results used here: \cite[Section~3.37.2]{LangevinLean2026ULASC}.
\subsection{\texorpdfstring{Differentiate the kernel by the squared-\allowbreak{}norm and exponential chain rules}{Differentiate the kernel by the squared-norm and exponential chain rules}}
\label{report-step-UnadjustedLangevin-LSI-GaussianEntropy-lean-4}


The map \(y\mapsto y-c\) has derivative the identity. The squared-norm derivative is \(v\mapsto2\langle y-c,v\rangle\). Multiplication by \(-a\), exponentiation, and multiplication by \(N_a\) therefore produce exactly

\[D_yK_a(c,y)[v]=-2aK_a(c,y)\langle y-c,v\rangle.\]

The squared-norm derivative contributes twice the displacement pairing, and division by \(4t\) leaves the coefficient \(-1/(2t)\). Multiplication by the unchanged kernel gives the explicit derivative formula.



Earlier results used here: \cite[Section~3.37.2]{LangevinLean2026ULASC}.
\subsection{\texorpdfstring{Find a common integrable bound for all component derivatives}{Find a common integrable bound for all component derivatives}}
\label{report-step-UnadjustedLangevin-LSI-GaussianEntropy-lean-5}


The norm of the inner-product functional \(v\mapsto\langle y-c,v\rangle\) equals \(\|y-c\|\). Hence

\[\|D_yK_a(c,y)\|=2aN_a\|y-c\|e^{-a\|y-c\|^2}\le2aN_a(1+a^{-1}).\]

The construction uses homogeneity of the norm and the operator norm of the inner-product functional, removes absolute values using positivity of \(N_a\) and the exponential, and applies the first Gaussian envelope bound. The resulting constant is independent of both \(y\) and \(c\). A probability mixture therefore integrates this bound without requiring any moment of its center map.



The common envelope depends on positive time and dimension but not on the initial state. Since the initial law has mass one, that constant envelope is integrable for every mixing distribution.



Earlier results used here: \cite[Section~3.37.2]{LangevinLean2026ULASC}; \cite[Section~3.37.4]{LangevinLean2026ULASC}.
\subsection{\texorpdfstring{Define the actual proposed mixture derivative}{Define the actual proposed mixture derivative}}
\label{report-step-UnadjustedLangevin-LSI-GaussianEntropy-lean-6}


For \(c_t(x)=x-t\nabla U(x)\), define the continuous-linear-map-valued integral

\[P'_t(y)=\int D_yK_{1/(4t)}(c_t(x),y)\,d\mu(x).\]

This is the integrated spatial kernel derivative. The symbol here denotes a spatial Fréchet derivative, not the time derivative of the density. The following domination lemmas establish its integrability, continuity, and equality with the actual spatial derivative of the explicit mixture. Keeping the definition as an integral of linear maps makes each directional evaluation correspond directly to the component kernel derivative.



Integrate the component derivative as a finite-dimensional linear map. Its norm is dominated by the common Gaussian derivative envelope, so this candidate is an actual Bochner integral rather than a purely symbolic expression.

\subsection{\texorpdfstring{Integrate each Gaussian component over the mixing probability law}{Integrate each Gaussian component over the mixing probability law}}
\label{report-step-UnadjustedLangevin-LSI-GaussianEntropy-lean-7}


For a measurable center map \(c\), the function \(x\mapsto K_a(c(x),y)\) is measurable and satisfies \(0\le K_a(c(x),y)\le N_a\). Since \(\mu\) is a probability measure, the constant \(N_a\) is integrable, and so is the kernel.



The calculation starts from integrability of a constant under the mixing probability and uses integrable domination. Its first obligation is measurability, assembled from the center's measurability by composition of the displayed continuous and measurable maps. Its second obligation is the norm bound; positivity removes the absolute value, and the fact that an exponential with nonpositive exponent is at most one supplies the Gaussian peak bound.



At a fixed observation the kernel is bounded by its common peak. Hence it is integrable over the starting probability and its mixture value is finite, as required before differentiating that mixture.



Earlier results used here: \cite[Section~3.37.2]{LangevinLean2026ULASC}; \cite[Section~3.37.4]{LangevinLean2026ULASC}.
\subsection{\texorpdfstring{Integrate the continuous-\allowbreak{}linear-\allowbreak{}map-\allowbreak{}valued component derivative}{Integrate the continuous-linear-map-valued component derivative}}
\label{report-step-UnadjustedLangevin-LSI-GaussianEntropy-lean-8}


The common derivative bound is a finite constant \(2aN_a(1+a^{-1})\). Therefore

\[\int\|D_yK_a(c(x),y)\|\,d\mu(x)<\infty.\]

The proof again applies integrable domination to an integrable constant, now with values in the Banach space of continuous linear functionals. Measurability follows from the explicit derivative formula; the second obligation is exactly the uniform Gaussian derivative envelope. This supplies the Bochner-integrability premise needed to identify the derivative of the mixture.



The linear-map-valued derivative has norm bounded by an integrable scalar envelope. This verifies both its Bochner integrability and the domination needed for the differentiation-under-integral theorem.



Earlier results used here: \cite[Section~3.37.2]{LangevinLean2026ULASC}.
\subsection{\texorpdfstring{Differentiate the mixture under the integral sign}{Differentiate the mixture under the integral sign}}
\label{report-step-UnadjustedLangevin-LSI-GaussianEntropy-lean-9}


For each proposed state \(y\),

\[D_y\!\left(\int K_a(c(x),y)\,d\mu(x)\right)=\int D_yK_a(c(x),y)\,d\mu(x).\]

Dominated differentiation applies because the component derivative is explicit, its norm has a common integrable Gaussian bound, and both kernel and derivative are measurable in the mixing variable. It gives the derivative of the mixture as the integral of the component derivative. Thus mixture differentiability is proved from the kernel, not assumed.



The component derivative formula and common integrable bound satisfy the dominated spatial differentiation hypotheses. Their integral is therefore the genuine derivative of the proposed mixture density.



Earlier results used here: \cite[Section~3.37.2]{LangevinLean2026ULASC}.
\subsection{\texorpdfstring{Prove continuity of the mixture derivative with the same domination}{Prove continuity of the mixture derivative with the same domination}}
\label{report-step-UnadjustedLangevin-LSI-GaussianEntropy-lean-10}


For every fixed initial point \(x\), the derivative kernel \(y\mapsto D_yK_a(c(x),y)\) is continuous. Its operator norm is bounded by the same integrable constant for every \(y\). Dominated continuity therefore proves continuity of

\[y\longmapsto\int D_yK_a(c(x),y)\,d\mu(x).\]

The proof theorem continuity under an integral with a common integrable majorant requires measurability of every section, a common integrable bound, and almost-everywhere continuity of the component map. The preceding lemmas provide the first three facts, while the explicit formula and composition of the displayed continuous and measurable maps prove continuity. This continuity is what strengthens mere differentiability to \(C^1\).



At each starting state the component derivative is continuous in the observation. The same bound dominates it near every observation, so dominated convergence proves continuity of the integrated derivative.



Earlier results used here: \cite[Section~3.37.2]{LangevinLean2026ULASC}.
\subsection{\texorpdfstring{Combine the derivative and its continuity into a C\ensuremath{{}^{1}} theorem}{Combine the derivative and its continuity into a C¹ theorem}}
\label{report-step-UnadjustedLangevin-LSI-GaussianEntropy-lean-11}


A function is \(C^1\) exactly when it has a continuous Fréchet derivative. Apply this criterion to the derivative field \(y\mapsto\int D_yK_a(c(x),y)\,d\mu(x)\).



The integrated derivative field is continuous by the common Gaussian domination. The preceding differentiation theorem identifies it as the derivative at every observation. Together these two facts give \(C^1\) regularity of the mixture for every probability mixing measure, including singular measures.



Differentiability of the mixture together with continuity of its derivative is precisely \(C^1\) regularity. This is stronger than merely continuity of the density and is what the LSI test class requires.

\subsection{\texorpdfstring{Apply Gaussian smoothing to the actual frozen-\allowbreak{}drift law}{Apply Gaussian smoothing to the actual frozen-drift law}}
\label{report-step-UnadjustedLangevin-LSI-GaussianEntropy-lean-12}


The explicit real density of the update satisfies

\[p_t(y)=\int K_{1/(4t)}(x-t\nabla U(x),y)\,d\mu(x).\]

The actual frozen-drift law is already identified with this Gaussian mixture. Its density therefore inherits the established \(C^1\) regularity, regardless of whether the initial law itself has a density.



Earlier results used here: \cite[Section~3.37.12]{LangevinLean2026ULASC}.
\subsection{\texorpdfstring{Identify the exact spatial derivative of the actual density}{Identify the exact spatial derivative of the actual density}}
\label{report-step-UnadjustedLangevin-LSI-GaussianEntropy-lean-13}


The derivative field defined earlier is now certified:

\[D_yp_t(y)=P'_t(y)=\int D_yK_{1/(4t)}(x-t\nabla U(x),y)\,d\mu(x).\]

The actual mixture has measurable frozen center \(c_t(x)=x-t\nabla U(x)\). Applying dominated Gaussian differentiation identifies its derivative with the integral of the displayed component derivative. This exact derivative field can then be chosen jointly measurably while retaining agreement with the actual spatial density derivative.



The spatial derivative of the actual interpolation density equals the integrated component derivative. This identification lets its score be related to Gaussian component scores in the subsequent posterior calculation.



Earlier results used here: \cite[Section~3.37.12]{LangevinLean2026ULASC}.
\subsection{\texorpdfstring{Recover C\ensuremath{{}^{1}} regularity of the potential from its Lipschitz gradient}{Recover C¹ regularity of the potential from its Lipschitz gradient}}
\label{report-step-UnadjustedLangevin-LSI-GaussianEntropy-lean-14}


The assumptions are differentiability of \(U\) and Lipschitz continuity of \(\nabla U\). The Riesz isometry identifies

\[DU(x)=\operatorname{Riesz}(\nabla U(x)).\]

The derivative of the differentiable potential is the inner-product functional represented by \(\nabla U\). A Lipschitz gradient is continuous, and the identification of vectors with these linear functionals is continuous. Therefore the potential has a continuous derivative and belongs to \(C^1\), with no Hessian assumption.



A globally Lipschitz gradient is continuous, and it represents the derivative of the differentiable potential. Thus the potential is \(C^1\), which supplies the regularity of its positive Gibbs density and reciprocal.

\subsection{\texorpdfstring{Establish spatial admissibility of the actual Gibbs-\allowbreak{}relative density}{Establish spatial admissibility of the actual Gibbs-relative density}}
\label{report-step-UnadjustedLangevin-LSI-GaussianEntropy-lean-15}


The relative density is \(r_t(y)=Z e^{U(y)}p_t(y)\). The partition factor \(Z\) is constant in space, \(e^U\) is \(C^1\) because \(U\) is \(C^1\), and the actual Gaussian mixture \(p_t\) is \(C^1\). Therefore \(r_t\) is \(C^1\) at every positive time.



The ordinary product rule combines the constant partition, the continuously differentiable exponential of the potential, and the continuously differentiable Gaussian mixture. The resulting relative density is positive and normalized by the established law identity, so it belongs to the target LSI test class.



The relative density is the product of the mixture density with the positive reciprocal Gibbs density. The product rule gives \(C^1\) regularity, while pointwise positivity and normalization are inherited from the actual law.

\subsection{\texorpdfstring{Apply the target LSI directly to the constructed interpolation law}{Apply the target LSI directly to the constructed interpolation law}}
\label{report-step-UnadjustedLangevin-LSI-GaussianEntropy-lean-16}


Use the actual positive normalized \(C^1\) relative density as the LSI test. This gives

\[2\alpha\,\mathrm{KL}(\rho_t\mid\nu)\le I(\rho_t\mid\nu).\]

This theorem retains extended nonnegative values, so neither finiteness is silently assumed. The proof obtains the Gibbs probability hypothesis from 1; 2 is applied with the already proved positivity, \(C^1\) regularity, and relative probability normalization. Finally tilted measure the Gibbs-relative density \(r_t=Ze^Up_t\) replaces the induced tilted measure with the actual Gaussian interpolation law. The only functional-inequality hypothesis is the target's LSI.



The constructed relative density belongs to the original positive \(C^1\) normalized LSI class. Therefore the inequality applies directly to its actual law, without assuming a new extension of LSI to arbitrary densities.



Earlier results used here: \cite[Section~3.26.3]{LangevinLean2026ULASC}.
\subsection{\texorpdfstring{Record the square-\allowbreak{}norm Jensen inequality for probability measures}{Record the square-norm Jensen inequality for probability measures}}
\label{report-step-UnadjustedLangevin-LSI-GaussianEntropy-lean-17}


For a square-integrable random vector \(F\) under a probability measure, convexity of squared norm gives

\[\left\|\int F\,d\mu\right\|^2\le\int\|F\|^2\,d\mu.\]

The proof obtains convexity by applying pow to convexity and nonnegativity of the norm, with exponent two. Jensen’s inequality for squared norm then applies Jensen on the whole vector space. square-integrability supplies both integrability of \(F\) and of \(\|F\|^2\). This general lemma is used under the normalized Gaussian posterior at each proposed state.



For a probability measure, expand \(\int\|v-\int v\|^2\ge0\). This gives \(\|\int v\|^2\le\int\|v\|^2\), the vector Jensen inequality used for the posterior mean score.

\subsection{\texorpdfstring{Define Bayes posterior on the initial mixture component}{Define Bayes posterior on the initial mixture component}}
\label{report-step-UnadjustedLangevin-LSI-GaussianEntropy-lean-18}


At a fixed proposed point \(y\), let \(p(y)=\int K_a(c(x),y)\,d\mu(x)\). The posterior law of the initial component is

\[\lambda_y(dx)=\frac{K_a(c(x),y)}{p(y)}\,\mu(dx).\]

Weight the initial probability by \(K_a(c(x),y)/p(y)\). Positive precision makes the kernel positive and finite, so its mixture denominator is positive. The next calculation shows that this weight integrates to one, establishing the posterior probability before Jensen is applied.



At observation \(y\), weight the initial law by \(K_t(x,y)/p_t(y)\). The denominator is the actual mixture value, making this a probability describing which initial component contributed to that observation.

\subsection{\texorpdfstring{Prove the real mixture denominator is everywhere positive}{Prove the real mixture denominator is everywhere positive}}
\label{report-step-UnadjustedLangevin-LSI-GaussianEntropy-lean-19}


Every component \(K_a(c(x),y)\) is positive because its normalizer and exponential are positive. Its support as a function of \(x\) is therefore the whole initial-state space, which has probability one.



Every Gaussian component is positive and the initial measure has mass one, so \(p_t(y)>0\) for every observation. This allows the posterior denominator and mixture log score to be formed pointwise.



Earlier results used here: \cite[Section~3.37.4]{LangevinLean2026ULASC}.
\subsection{\texorpdfstring{Check posterior normalization by cancellation of the actual mixture integral}{Check posterior normalization by cancellation of the actual mixture integral}}
\label{report-step-UnadjustedLangevin-LSI-GaussianEntropy-lean-20}


The total posterior mass equals

\[\int\frac{K_a(c(x),y)}{p(y)}\,d\mu(x)=\frac{p(y)}{p(y)}=1.\]

Integrating \(K_t(x,y)/p_t(y)\) in \(x\) gives \(p_t(y)/p_t(y)=1\). Strict positivity and finiteness justify the cancellation, establishing posterior normalization.



Earlier results used here: \cite[Section~3.37.4]{LangevinLean2026ULASC}.
\subsection{\texorpdfstring{Prove posterior displacement has finite conditional second moment}{Prove posterior displacement has finite conditional second moment}}
\label{report-step-UnadjustedLangevin-LSI-GaussianEntropy-lean-21}


The posterior may arise from an initial law without moments. Nevertheless, at each fixed \(y\),

\[\int\|y-c(x)\|^2\,d\lambda_y(x)=\frac{1}{p(y)}\int K_a(c(x),y)\|y-c(x)\|^2\,d\mu(x)<\infty.\]

The Gaussian displacement-square factor times its exponential is uniformly bounded as a function of the center. Dividing its finite initial-law integral by positive \(p_t(y)\) gives a finite conditional second moment.



Earlier results used here: \cite[Section~3.37.2]{LangevinLean2026ULASC}; \cite[Section~3.37.4]{LangevinLean2026ULASC}.
\subsection{\texorpdfstring{Identify the mixture score as the posterior component-\allowbreak{}score mean}{Identify the mixture score as the posterior component-score mean}}
\label{report-step-UnadjustedLangevin-LSI-GaussianEntropy-lean-22}


Dividing the mixture derivative by its positive value gives

\[\nabla\log p(y)=-2a\int(y-c(x))\,d\lambda_y(x).\]

Divide the integrated gradient identity by \(p_t(y)\). The result is the posterior average of \(-(y-x+t\nabla U(x))/(2t)\), identifying the actual mixture score exactly.



Earlier results used here: \cite[Section~3.37.4]{LangevinLean2026ULASC}.
\subsection{\texorpdfstring{Use posterior Jensen to bound the pointwise mixture score energy}{Use posterior Jensen to bound the pointwise mixture score energy}}
\label{report-step-UnadjustedLangevin-LSI-GaussianEntropy-lean-23}


Applying square-norm Jensen to the posterior score identity yields

\[p(y)\|\nabla\log p(y)\|^2\le4a^2\int K_a(c(x),y)\|y-c(x)\|^2\,d\mu(x).\]

Apply vector Jensen under the normalized posterior to the component score. Its squared posterior mean is bounded by the posterior mean of its squared norm, with no extra factor two.



Earlier results used here: \cite[Section~3.37.4]{LangevinLean2026ULASC}.
\subsection{\texorpdfstring{Compute the exact Gaussian component displacement energy}{Compute the exact Gaussian component displacement energy}}
\label{report-step-UnadjustedLangevin-LSI-GaussianEntropy-lean-24}


At precision \(1/(4t)\), the component density is the law of \(c+\sqrt{2t}Z\). Therefore

\[\int K_{1/(4t)}(c,y)\|y-c\|^2\,dy=\mathbb E\|\sqrt{2t}Z\|^2=2td.\]

Conditional on the initial point, the centered displacement is \(\sqrt{2t}Z\). Dividing its squared moment \(2td\) by \(4t^2\) gives component-score energy \(d/(2t)\).



Earlier results used here: \cite[Section~3.37.17]{LangevinLean2026ULASC}; \cite[Section~3.37.2]{LangevinLean2026ULASC}; \cite[Section~3.37.4]{LangevinLean2026ULASC}.
\subsection{\texorpdfstring{Integrate posterior Jensen to obtain the universal smoothing Fisher bound}{Integrate posterior Jensen to obtain the universal smoothing Fisher bound}}
\label{report-step-UnadjustedLangevin-LSI-GaussianEntropy-lean-25}


Integrate the pointwise score inequality, use Fubini, and insert the component displacement moment:

\[\int p(y)\|\nabla\log p(y)\|^2\,dy\le4\left(\frac1{4t}\right)^2\int2td\,d\mu=\frac d{2t}.\]

Multiply the posterior Jensen inequality by \(p_t(y)\) and integrate in \(y\). The posterior denominator cancels, and Tonelli reduces the upper bound to the exact Gaussian component-score moment.



Earlier results used here: \cite[Section~3.37.4]{LangevinLean2026ULASC}.
\subsection{\texorpdfstring{Propagate finite gradient energy without a state second moment}{Propagate finite gradient energy without a state second moment}}
\label{report-step-UnadjustedLangevin-LSI-GaussianEntropy-lean-26}


Write

\[\nabla U(Y_t)=\bigl(\nabla U(Y_t)-\nabla U(X)\bigr)+\nabla U(X).\]

The first term is in \(L^2\) because Lipschitzness bounds it by \(L\|Y_t-X\|\), and the actual frozen displacement has finite second moment whenever \(\nabla U(X)\in L^2\). The second term is in \(L^2\) by assumption. The construction uses the Lipschitz-gradient displacement-energy bound and the \(L^2\) characterization for the discrepancy, adds the pulled-back initial gradient, and transfers the sum to the output measure with the pushforward integral identity for square energies. The initial position cancels from the displacement, which is why no \(\mathbb E\|X\|^2\) premise appears.



The gradient difference is at most \(L\) times displacement, so \(\|\nabla U(Y_t)\|^2\le2\|\nabla U(X)\|^2+2L^2\|Y_t-X\|^2\). This uses only a score moment of the initial law.



Earlier results used here: \cite[Section~3.34.1]{LangevinLean2026ULASC}; \cite[Section~3.36.5]{LangevinLean2026ULASC}.
\subsection{\texorpdfstring{Transfer the smoothing Fisher estimate to the actual update law}{Transfer the smoothing Fisher estimate to the actual update law}}
\label{report-step-UnadjustedLangevin-LSI-GaussianEntropy-lean-27}


The generic mixture Fisher theorem applies with \(c_t(x)=x-t\nabla U(x)\). Rewriting its integral through \(\rho_t=p_t\,dy\) gives

\[\nabla\log p_t\in L^2(\rho_t),\qquad\int\|\nabla\log p_t\|^2\,d\rho_t\le\frac d{2t}.\]

The actual proposal pushforward and mixture density describe the same law. Therefore their Lebesgue-score square integrals agree, transferring the universal smoothing Fisher bound to the law used in entropy production.



Earlier results used here: \cite[Section~3.37.12]{LangevinLean2026ULASC}.
\subsection{\texorpdfstring{Retain everywhere positivity for the real density}{Retain everywhere positivity for the real density}}
\label{report-step-UnadjustedLangevin-LSI-GaussianEntropy-lean-28}


The extended mixture density is positive and finite at every point for \(t>0\). Its real conversion is consequently positive:

\[p_t(y)>0.\]

Every Gaussian component is positive, so its mixture against a probability is positive. The common Gaussian peak also makes it finite. Hence \(0<p_t(y)<\infty\) at every observation, justifying the logarithm and reciprocal used in its score.



Pointwise finiteness and positivity are retained when choosing the real representative of the mixture. Hence its logarithm and reciprocal may be differentiated at every observation, rather than only on an unspecified almost-everywhere set.

\subsection{\texorpdfstring{Split the relative score into potential and Lebesgue scores}{Split the relative score into potential and Lebesgue scores}}
\label{report-step-UnadjustedLangevin-LSI-GaussianEntropy-lean-29}


Since \(r_t=Ze^Up_t\), positivity and finite normalization give

\[\log r_t=\log Z+U+\log p_t,\qquad\nabla\log r_t=\nabla U+\nabla\log p_t.\]

Since \(\log r_t=\log p_t+U+\log Z\), spatial differentiation gives \(\nabla\log r_t=\nabla\log p_t+\nabla U\). The constant normalization contributes no gradient.



Earlier results used here: \cite[Section~3.26.2]{LangevinLean2026ULASC}; \cite[Section~3.26.3]{LangevinLean2026ULASC}; \cite[Section~3.26.9]{LangevinLean2026ULASC}.
\subsection{\texorpdfstring{Construct finite relative-\allowbreak{}score energy from two proved L\ensuremath{{}^{2}} fields}{Construct finite relative-score energy from two proved L² fields}}
\label{report-step-UnadjustedLangevin-LSI-GaussianEntropy-lean-30}


The potential gradient remains in \(L^2(\rho_t)\), while the Gaussian smoothing bound gives \(\nabla\log p_t\in L^2(\rho_t)\). Their sum is therefore in \(L^2\), and the score decomposition identifies it with the actual relative score:

\[\nabla\log r_t\in L^2(\rho_t).\]

The proof obtains continuity of the Lebesgue score from \(C^1\) regularity of its positive density, converts the finite Fisher integral to square-integrability, and adds it to the propagated finite gradient energy. The final function extensionality uses the score decomposition \(\nabla\log r_t=\nabla\log p_t+\nabla U\) pointwise. This is the finite-energy fact used for the bounded logarithmic entropy tests.



The two square-integrable fields are the Lebesgue score and the propagated potential gradient. Their sum has energy at most twice each separate energy, providing finite relative-score energy.



Earlier results used here: \cite[Section~3.26.3]{LangevinLean2026ULASC}.
\subsection{\texorpdfstring{Prove finite relative Fisher information before taking real values}{Prove finite relative Fisher information before taking real values}}
\label{report-step-UnadjustedLangevin-LSI-GaussianEntropy-lean-31}


Square integrability of the actual relative score implies

\[I(\rho_t\mid\nu)=\int\|\nabla\log r_t\|^2\,d\rho_t<\infty.\]

The squared relative score is integrable under the actual interpolation law. Its nonnegative integral is therefore finite and equals its ordinary energy integral. This identifies actual relative Fisher information as a finite number before LSI or entropy production is combined with it.



Integrating the nonnegative squared relative score against the actual proposed law gives the same finite ordinary energy already constructed. Thus the extended relative Fisher information is finite before any real-value use.



Earlier results used here: \cite[Section~3.26.3]{LangevinLean2026ULASC}.
\subsection{\texorpdfstring{Derive finite actual KL from LSI and the constructed finite Fisher energy}{Derive finite actual KL from LSI and the constructed finite Fisher energy}}
\label{report-step-UnadjustedLangevin-LSI-GaussianEntropy-lean-32}


At positive time, the constructed density is an admissible LSI test, so \(2\alpha\,\mathrm{KL}(\rho_t\mid\nu)\le I(\rho_t\mid\nu)\). With \(\alpha>0\) and finite Fisher information, the KL divergence must be finite.



The construction uses the target's LSI probability hypothesis, the finite relative-score energy bound, and the target LSI applied to the constructed positive normalized density. If KL were infinity, its product with the positive nonzero factor a nonnegative scalar (2 * \ensuremath{\alpha}) would be infinity. The inequality would then force Fisher information to be infinity, contradicting the proved finiteness. No entropy continuity or initial density smoothness is assumed to obtain this positive-time finiteness.



LSI bounds nonnegative actual KL by the finite relative Fisher information divided by \(2\alpha>0\). Consequently positive-time actual entropy is finite, completing the finiteness chain needed for its later derivative identity.



Earlier results used here: \cite[Section~3.26.3]{LangevinLean2026ULASC}.
\section{\texorpdfstring{Time derivatives and spacetime integrability under the Gaussian mixture}{Time derivatives and spacetime integrability under the Gaussian mixture}}
\label{proof-module-UnadjustedLangevin-LSI-GaussianTimeDerivative-lean}


Let \(g_x=\nabla U(x)\), \(c_t(x)=x-tg_x\), \(d=\dim E\), and \(K_t(x,y)=(4\pi t)^{-d/2}\exp(-\|y-c_t(x)\|^2/(4t))\). This module differentiates the actual normalized proposal, constructs \(\dot p_t(y)=\int\partial_tK_t(x,y)\,\mu(dx)\), and proves its local spacetime integrability. The only initial moment used is \(G_0=\int\|g_x\|^2d\mu<\infty\); this is derived from finite initial KL elsewhere.



The explicit kernel time derivative is \(\dot K_t=K_t[-d/(2t)+\|y-x\|^2/(4t^2)-\|\nabla U(x)\|^2/4]\). On a compact positive time interval, reciprocal powers of time are bounded and Gaussian polynomial factors have integrable envelopes. The finite initial score energy controls the remaining starting-point term. These estimates justify both differentiating the mixture in time and integrating its derivative over space and time; they are established before the entropy chain rule uses that derivative.

\subsection{\texorpdfstring{The Gaussian normalizer as an explicit time power}{The Gaussian normalizer as an explicit time power}}
\label{report-step-UnadjustedLangevin-LSI-GaussianTimeDerivative-lean-1}


The proposal normalizer is defined as the reciprocal of an actual Gaussian integral, not as a free scalar parameter. Gaussian Fourier integration evaluates \(\int e^{-a\|z\|^2}dz=(\pi/a)^{d/2}\) for \(a>0\). Substitute \(a=1/(4t)\), simplify \(\pi/a=4\pi t\), and invert the real power to obtain \(N_t=(4\pi t)^{-d/2}\). Positivity of \(t\) justifies the Gaussian formula and the real-power inversion.



The Gaussian normalizer has time dependence \((4\pi t)^{-d/2}\). Expressing it as a real power separates the positive constant from the time power and permits ordinary scalar differentiation.



Earlier results used here: \cite[Section~3.37.1]{LangevinLean2026ULASC}; \cite[Section~3.37.12]{LangevinLean2026ULASC}; \cite[Section~3.37.2]{LangevinLean2026ULASC}.
\subsection{\texorpdfstring{The time derivative of the Gaussian normalizer}{The time derivative of the Gaussian normalizer}}
\label{report-step-UnadjustedLangevin-LSI-GaussianTimeDerivative-lean-2}


Differentiating \(t^{-d/2}\) gives \(-(d/(2t))t^{-d/2}\). Positive time supplies the reciprocal and leaves the normalizer multiplied by its logarithmic derivative.



Earlier results used here: \cite[Section~3.37.12]{LangevinLean2026ULASC}.
\subsection{\texorpdfstring{The frozen-\allowbreak{}drift Gaussian kernel}{The frozen-drift Gaussian kernel}}
\label{report-step-UnadjustedLangevin-LSI-GaussianTimeDerivative-lean-3}


The real component kernel is \(K_t(x,y)=N_t\exp(-\|y-x+tg_x\|^2/(4t))\). The gradient \(g_x\) is evaluated at the initial state \(x\), and is constant when differentiating in \(t\) or \(y\). Its integral over \(y\) is one at positive time by the previously proved Gaussian proposal normalization.



The frozen kernel is the actual Gaussian density conditional on the initial state. Its center moves linearly with time through the fixed vector \(\nabla U(x)\), while its covariance grows as \(2tI_d\).



Earlier results used here: \cite[Section~3.37.12]{LangevinLean2026ULASC}.
\subsection{\texorpdfstring{Its explicit candidate time derivative}{Its explicit candidate time derivative}}
\label{report-step-UnadjustedLangevin-LSI-GaussianTimeDerivative-lean-4}


Define the explicit candidate derivative \(\dot K_t=K_t[-d/(2t)+\|y-x\|^2/(4t^2)-\|g_x\|^2/4]\). The dimension term comes from the normalizer; the remaining terms come from differentiating the exponent. In the expanded exponent, the mixed inner-product term is independent of time, which explains its absence from this expression. The following theorem verifies that this formula is the derivative.



The derivative field is defined by the explicit normalizer and exponent formula. Naming this function before differentiation makes later measurability and integrability statements refer to one fixed spacetime integrand.

\subsection{\texorpdfstring{Expanding the frozen Gaussian exponent}{Expanding the frozen Gaussian exponent}}
\label{report-step-UnadjustedLangevin-LSI-GaussianTimeDerivative-lean-5}


Set \(z=y-x\). The squared-norm identity gives \(\|z+sg_x\|^2=\|z\|^2+2s\langle z,g_x\rangle+s^2\|g_x\|^2\). For \(s\ne0\), division by \(-4s\) yields \(-\|z+sg_x\|^2/(4s)=-\|z\|^2/(4s)-\langle z,g_x\rangle/2-s\|g_x\|^2/4\). The vector identity and scalar denominator cancellation are proved separately in the proof calculation.



Expanding \(\|y-x+t g\|^2\) and dividing by \(4t\) yields \(\|y-x\|^2/(4t)+\langle y-x,g\rangle/2+t\|g\|^2/4\). The middle term has zero time derivative, simplifying the formula.

\subsection{\texorpdfstring{Differentiating the frozen Gaussian kernel}{Differentiating the frozen Gaussian kernel}}
\label{report-step-UnadjustedLangevin-LSI-GaussianTimeDerivative-lean-6}


Write the expanded exponent as \(-A/s-B-sC\), where \(A=\|y-x\|^2/4\), \(B=\langle y-x,g_x\rangle/2\), and \(C=\|g_x\|^2/4\). Its derivative at positive \(t\) is \(A/t^2-C\). The expansion agrees with the original exponent on a neighborhood avoiding zero, so this is also the derivative of the original expression.



Apply the exponential chain rule and multiply by the actual normalizer. Its derivative contributes \(-dK_t/(2t)\); the exponential derivative contributes \(K_t(A/t^2-C)\). Collecting terms gives exactly its explicit time derivative. No derivative of \(g_x\), and therefore no Hessian, is taken.



Combine the normalizer derivative with the derivative of the expanded exponent. The cross term contributes no time derivative, leaving exactly the stated squared-position and squared-gradient terms.

\subsection{\texorpdfstring{An absolute bound on the kernel time derivative}{An absolute bound on the kernel time derivative}}
\label{report-step-UnadjustedLangevin-LSI-GaussianTimeDerivative-lean-7}


Put \(r=y-x+tg_x\) and \(e=e^{-\|r\|^2/(4t)}\). The triangle inequality and \((a+b)^2\le2a^2+2b^2\) give \(\|y-x\|^2\le2\|r\|^2+2t^2\|g_x\|^2\). The Gaussian envelope gives \(\|r\|^2e\le4t\), while \(e\le1\). Consequently \(\|y-x\|^2e\le8t+2t^2\|g_x\|^2\).



Bound the absolute value of each derivative-coefficient term by their nonnegative sum, multiply by \(e\), and insert these bounds. The resulting convenient envelope is \(|\dot K_t(x,y)|\le N_t[(d+4)/t+\|g_x\|^2]\). It is uniform in \(y\), and is integrable in the initial variable whenever \(G_0<\infty\). Its purpose is time differentiation under the initial-state integral, rather than output-space integration.



Take absolute values term by term in the explicit derivative. Nonnegativity of the kernel allows its factor to be retained, producing a sum of Gaussian-weighted nonnegative energy terms.



Earlier results used here: \cite[Section~3.37.12]{LangevinLean2026ULASC}; \cite[Section~3.37.14]{LangevinLean2026ULASC}.
\subsection{\texorpdfstring{Measurability of the normalizer}{Measurability of the normalizer}}
\label{report-step-UnadjustedLangevin-LSI-GaussianTimeDerivative-lean-8}


At positive time the Gaussian normalizer has the continuous explicit form \((4\pi t)^{-d/2}\), hence is measurable. Assigning an arbitrary value at zero gives a measurable extension if a whole-time parameter field is desired. Only the positive-time values are used in the Gaussian density and derivative identities.



At positive time the normalizer is continuous, hence measurable. Its reciprocal-power expression also permits a measurable extension when the proof organizes all times in one ambient domain.



Earlier results used here: \cite[Section~3.37.12]{LangevinLean2026ULASC}.
\subsection{\texorpdfstring{Joint measurability of the kernel time derivative}{Joint measurability of the kernel time derivative}}
\label{report-step-UnadjustedLangevin-LSI-GaussianTimeDerivative-lean-9}


Lipschitzness makes \(x\mapsto g_x\) continuous and measurable. Combine this with the measurable normalizer and the explicit algebraic formula for \(\dot K_t\). Products, norms, scalar inverses, and exponentials preserve measurability, giving joint measurability in \((t,y,x)\). This is the input needed to integrate out the initial variable while retaining measurable time-space dependence.



The explicit derivative uses only measurable gradient, norms, scalar operations, and the Gaussian exponential. Their composition is jointly measurable in time, initial state, and observation.

\subsection{\texorpdfstring{The candidate time derivative of the mixture}{The candidate time derivative of the mixture}}
\label{report-step-UnadjustedLangevin-LSI-GaussianTimeDerivative-lean-10}


Define \(\dot p_t(y)=\int\dot K_t(x,y)\,d\mu(x)\). At this point it is an ordinary parameter integral. The subsequent differentiation theorem proves it is the actual derivative of \(p_t(y)\), and the energy estimates prove its integrability, so no analytic assertion is built into this definition.



Integrating the explicit kernel derivative over the initial probability defines the candidate mixture time derivative. Its later identification as a genuine derivative requires the uniform envelope proved separately.

\subsection{\texorpdfstring{Joint measurability of the mixture derivative}{Joint measurability of the mixture derivative}}
\label{report-step-UnadjustedLangevin-LSI-GaussianTimeDerivative-lean-11}


The component derivative is jointly measurable in \((t,y,x)\). Its strongly measurable parameter integral over the probability measure \(\mu\) is therefore measurable in \((t,y)\). This gives a single measurable version of \(\dot p\) for spacetime Fubini and the entropy chain rule.



Parameter integration of the jointly measurable derivative field gives a measurable candidate in time and observation. Absolute integrability, when invoked later, ensures it is also the ordinary integral required by the chain rule.

\subsection{\texorpdfstring{The genuine time derivative of the mixture}{The genuine time derivative of the mixture}}
\label{report-step-UnadjustedLangevin-LSI-GaussianTimeDerivative-lean-12}


Fix \(t>0\). Continuity of \(N_t\) supplies a common neighborhood \(S\) on which \(s>t/2\) and \(N_s<N_t+1\). The pointwise derivative envelope then gives \(|\dot K_s(x,y)|\le(N_t+1)[2(d+4)/t+\|g_x\|^2]\) for every \(s\in S\). This bound is integrable in \(x\) by the initial gradient-energy hypothesis, and the component kernel at \(t\) is integrable because it is bounded by \(N_t\).



A positive-time neighborhood bounds reciprocal time factors. An integrable kernel derivative envelope then permits time differentiation under the initial-law integral, proving that the candidate is the actual density derivative.



Earlier results used here: \cite[Section~3.37.12]{LangevinLean2026ULASC}; \cite[Section~3.37.14]{LangevinLean2026ULASC}.
\subsection{\texorpdfstring{Equality of the two Gaussian kernel formulas}{Equality of the two Gaussian kernel formulas}}
\label{report-step-UnadjustedLangevin-LSI-GaussianTimeDerivative-lean-13}


The spatial Fisher calculation uses precision notation \(q_a(c,y)=Z_a^{-1}e^{-a\|y-c\|^2}\). Substitute \(a=1/(4t)\) and \(c=x-tg_x\). The normalizers agree by definition, and \(-a\|y-c\|^2=-\|y-c\|^2/(4t)\). This pointwise identity makes the time and spatial calculations use precisely the same normalized Gaussian kernel.



Both kernels have the same Gaussian center and covariance. Substituting their normalizers and exponents proves equality pointwise, allowing spatial and time results developed in the two notations to be combined.



Earlier results used here: \cite[Section~3.37.12]{LangevinLean2026ULASC}.
\subsection{\texorpdfstring{A weighted derivative bound}{A weighted derivative bound}}
\label{report-step-UnadjustedLangevin-LSI-GaussianTimeDerivative-lean-14}


Retain the Gaussian factor instead of replacing it by its supremum. With \(r=y-x+tg_x\), the bound \(\|y-x\|^2\le2\|r\|^2+2t^2\|g_x\|^2\) implies that the absolute derivative coefficient is at most \(d/(2t)+\|r\|^2/(2t^2)+\|g_x\|^2\). Multiplication by \(K_t\ge0\) yields \(|\dot K_t|\le K_t[d/(2t)+\|r\|^2/(2t^2)+\|g_x\|^2]\). Unlike the initial-variable envelope, this bound remains integrable in the output variable.



The weighted envelope retains the kernel's Gaussian decay and separates the initial-gradient square from displacement-dependent terms. This is the form suitable for integrating first in the observation.



Earlier results used here: \cite[Section~3.37.14]{LangevinLean2026ULASC}.
\subsection{\texorpdfstring{The Gaussian energy envelope}{The Gaussian energy envelope}}
\label{report-step-UnadjustedLangevin-LSI-GaussianTimeDerivative-lean-15}


Use the common nonnegative envelope \(\mathcal E_t(x,y)=K_t(x,y)(1+\|g_x\|^2+\|y-c_t(x)\|^2)\). Its three terms correspond to Gaussian mass, initial gradient energy, and Gaussian displacement energy. This envelope can dominate the derivative with a time-dependent scalar coefficient and can also be reused in spatial integration by parts.



The energy combines one, displacement square, and initial-gradient square under the Gaussian kernel. These are chosen to dominate every polynomial term appearing in the explicit time derivative.

\subsection{\texorpdfstring{Nonnegativity of the energy envelope}{Nonnegativity of the energy envelope}}
\label{report-step-UnadjustedLangevin-LSI-GaussianTimeDerivative-lean-16}


At positive time the actual normalizer and the Gaussian exponential are positive, so \(K_t\ge0\). Every term in \(1+\|g_x\|^2+\|y-c_t(x)\|^2\) is nonnegative. Their product is nonnegative; this identifies its norm with itself in the Fubini integrability proof.



Every component is a nonnegative square or constant multiplied by a nonnegative kernel. Their sum is therefore nonnegative, permitting nonnegative Fubini before finiteness is established.



Earlier results used here: \cite[Section~3.37.14]{LangevinLean2026ULASC}.
\subsection{\texorpdfstring{The component energy integral}{The component energy integral}}
\label{report-step-UnadjustedLangevin-LSI-GaussianTimeDerivative-lean-17}


For fixed \(x\), Gaussian normalization gives \(\int K_t(x,y)dy=1\), and the centered Gaussian covariance gives \(\int K_t(x,y)\|y-c_t(x)\|^2dy=2td\). Both integrands are genuinely integrable. The initial gradient factor is constant in \(y\), so its integral is \(\|g_x\|^2\). Expand the common energy into these three integrable terms to obtain \(\int\mathcal E_t(x,y)dy=1+\|g_x\|^2+2td\).



For each fixed initial state, integrate the kernel to one and its displacement square using the Gaussian moment formula. The resulting expression depends only on time, dimension, and the initial-gradient norm.



Earlier results used here: \cite[Section~3.37.6]{LangevinLean2026ULASC}; \cite[Section~3.37.7]{LangevinLean2026ULASC}.
\subsection{\texorpdfstring{Integrability of the pair energy envelope}{Integrability of the pair energy envelope}}
\label{report-step-UnadjustedLangevin-LSI-GaussianTimeDerivative-lean-18}


The common energy is jointly measurable. For every initial \(x\), its output section is integrable by the component energy theorem. Since the energy is nonnegative, the integral of its norm is exactly \(1+\|g_x\|^2+2td\), an integrable function of \(x\) under the probability law \(\mu\). The product integrability criterion therefore proves \(\mathcal E_t\in L^1(\mu\otimes dy)\).



The initial-gradient square is integrable by the score-energy hypothesis, and all remaining terms are finite constants. The component integral calculation therefore proves integrability of the full pair energy.

\subsection{\texorpdfstring{The integrated pair energy}{The integrated pair energy}}
\label{report-step-UnadjustedLangevin-LSI-GaussianTimeDerivative-lean-19}


Joint integrability allows Fubini, giving \(\int\mathcal E_t\,d(\mu\otimes dy)=\int(1+\|g_x\|^2+2td)d\mu(x)\). Probability normalization and the definition of \(G_0\) reduce this to \(1+G_0+2td\). The displayed equality supplies a finite quantitative bound, rather than relying on an unspecified integrability witness.



Integrating the component formula over the starting law replaces the squared gradient by \(G_0\). This gives an explicit finite energy expression, retaining its time dependence for uniform compact-interval estimates.

\subsection{\texorpdfstring{A uniform coefficient on a positive time interval}{A uniform coefficient on a positive time interval}}
\label{report-step-UnadjustedLangevin-LSI-GaussianTimeDerivative-lean-20}


Set \(C_t=d/(2t)+1/(2t^2)+1\). It is positive for \(t>0\), and is bounded above on any interval \([a,b]\) with \(a>0\) by \(C_a\). The coefficient is deliberately coarse: only local spacetime integrability, rather than a sharp density-derivative estimate, is needed for the entropy chain rule.



On a positive compact interval, \(t\) is bounded above and \(1/t\) is bounded above by the reciprocal of the lower endpoint. One coefficient can therefore dominate all scalar time factors uniformly.

\subsection{\texorpdfstring{Derivative domination by the energy envelope}{Derivative domination by the energy envelope}}
\label{report-step-UnadjustedLangevin-LSI-GaussianTimeDerivative-lean-21}


In the weighted derivative bound, denote the three nonnegative coefficients by \(A=d/(2t)\), \(B=1/(2t^2)\), and \(1\). Each is at most \(C_t=A+B+1\). Therefore \(A+B\|r\|^2+\|g_x\|^2\le C_t(1+\|g_x\|^2+\|r\|^2)\). Multiply by \(K_t\ge0\) to prove \(|\dot K_t|\le C_t\mathcal E_t\).



Multiply that coefficient by the nonnegative Gaussian energy envelope. The pointwise derivative bound then holds throughout the interval with one common integrable function, the key hypothesis for spacetime domination.



Earlier results used here: \cite[Section~3.37.14]{LangevinLean2026ULASC}.
\subsection{\texorpdfstring{Integrability of the kernel time derivative}{Integrability of the kernel time derivative}}
\label{report-step-UnadjustedLangevin-LSI-GaussianTimeDerivative-lean-22}


The component time derivative is jointly measurable, and its norm is bounded by \(C_t\mathcal E_t\). The common energy is integrable on \(\mu\otimes dy\), and \(C_t\) is a finite scalar at positive time. Domination therefore proves joint integrability of \(\dot K_t\), which licenses the changes of integration order used in the time weak equation.



The energy envelope has finite pair integral and the derivative is measurable. Absolute-value domination consequently proves integrability of the actual kernel time derivative under the Gaussian pair measure.

\subsection{\texorpdfstring{Spatial integrability of the mixture derivative}{Spatial integrability of the mixture derivative}}
\label{report-step-UnadjustedLangevin-LSI-GaussianTimeDerivative-lean-23}


Fubini applied to the jointly integrable component derivative first gives \(\dot p_t\in L^1(dy)\). The norm inequality for Bochner integration gives \(|\dot p_t(y)|\le\int|\dot K_t(x,y)|d\mu(x)\). Integrate in \(y\), exchange the two integrals, and dominate by \(C_t\mathcal E_t\). The exact envelope integral yields \(\int|\dot p_t(y)|dy\le C_t(1+G_0+2td)\).



The norm of the mixture derivative is at most the mixture of the kernel derivative norms. Fubini and the finite pair bound therefore give spatial integrability after the initial state is averaged out.

\subsection{\texorpdfstring{Spacetime integrability of the mixture derivative}{Spacetime integrability of the mixture derivative}}
\label{report-step-UnadjustedLangevin-LSI-GaussianTimeDerivative-lean-24}


Joint measurability of \(\dot p\) permits the product integrability criterion on \([a,b]\times E\), where \(a>0\). Every time section is in \(L^1(dy)\). For times in the interval, \(C_t\le C_a\) and \(1+G_0+2td\le1+G_0+2bd\), giving the common finite scalar bound \(\int|\dot p_t|dy\le C_a(1+G_0+2bd)\). The time interval has finite measure, so the section norms are integrable. Hence \(\dot p\in L^1(dt\,dy)\) locally away from zero. If \(a>b\), the restricted time measure is zero and the same argument remains valid.



The uniform compact-time coefficient and the energy integral provide an integrable time majorant. Integrating that majorant over the finite interval proves the full spacetime \(L^1\) estimate used by bounded entropy calculus.

\section{\texorpdfstring{Jointly measurable density and relative score}{Jointly measurable density and relative score}}
\label{proof-module-UnadjustedLangevin-LSI-GaussianMeasurability-lean}


Spacetime integration requires joint measurability in time and state, rather than measurability separately at each fixed time. This module constructs a single measurable extension of the actual positive-time relative score and proves that it agrees with the gradient used in Fisher information.



All positive-time kernels and their explicit derivatives are jointly measurable in time, starting point, and observation. Parameter integration therefore yields jointly measurable mixture fields. Dividing by the positive mixture gives one measurable score field for every spacetime integral. At each positive time this field agrees pointwise with the spatial derivative expression already obtained by Gaussian smoothing; assigning an arbitrary value at time zero does not affect integrals over positive time.

\subsection{\texorpdfstring{Integrate a jointly measurable Gaussian kernel}{Integrate a jointly measurable Gaussian kernel}}
\label{report-step-UnadjustedLangevin-LSI-GaussianMeasurability-lean-1}


The density is a nonnegative parameter integral \(p_t(y)=\int K_t(x,y)d\mu(x)\). The normalizer, affine center, norm, and exponential are measurable in \((t,y,x)\). The parameterized nonnegative integral is therefore measurable. Passing from the extended integral to its real value preserves measurability. Multiplication by \(Z e^{U(y)}\) gives the joint measurability of \(r_t(y)\).



The explicit positive-time formulas give one jointly measurable family on the domain \((0,\infty)\times E\). If a value at zero is needed to define an ambient spacetime field, it may be assigned separately on that measurable time slice. No choice at zero affects any positive-time density identity or spacetime integral.



The Gaussian kernel is a composition of measurable algebraic operations with the continuous exponential. Parameter integration preserves joint measurability in the remaining time and observation variables.



Earlier results used here: \cite[Section~3.37.13]{LangevinLean2026ULASC}.
\subsection{\texorpdfstring{Measurability of the spatial derivative}{Measurability of the spatial derivative}}
\label{report-step-UnadjustedLangevin-LSI-GaussianMeasurability-lean-2}


The derivative kernel is a continuous linear functional proportional to the inner product with the Gaussian residual:

\[D_yK_t(x,y)[v]=-\frac{K_t(x,y)}{2t}\langle y-x+t\nabla U(x),v\rangle.\]

The inner-product-to-dual map is continuous, so this is jointly strongly measurable as a continuous-linear-map-valued function. Its Bochner integral in \(x\), denoted \(A_t(y)\), remains strongly measurable. At positive times the earlier differentiation theorem identifies \(A_t=D p_t\).



The explicit spatial derivative is the kernel times a linear displacement functional. Its joint measurability and finite-dimensional coordinate integration yield a measurable mixture derivative field.



Earlier results used here: \cite[Section~3.37.2]{LangevinLean2026ULASC}.
\subsection{\texorpdfstring{One relative-\allowbreak{}score field for all spacetime integrals}{One relative-score field for all spacetime integrals}}
\label{report-step-UnadjustedLangevin-LSI-GaussianMeasurability-lean-3}


Define

\[\mathcal S(t,y)=\nabla U(y)+p_t(y)^{-1}\mathcal R^{-1}(A_t(y)),\]

where \(\mathcal R^{-1}\) is the inverse Riesz identification of the continuous dual with the Hilbert space. Continuous scalar multiplication, inversion as a measurable real function, and the previous measurability results make \(\mathcal S\) jointly strongly measurable.



For \(t>0\), \(p_t>0\), so the log chain rule and \(r_t=Z e^U p_t\) give

\[\mathcal S(t,y)=\nabla U(y)+\nabla\log p_t(y)=\nabla\log r_t(y).\]

Thus the measurable field is the actual relative score wherever it is used. Its auxiliary extension at nonpositive time introduces no analytic premise into the proof.



Dividing the mixture gradient by its strictly positive density defines one jointly measurable Lebesgue score. Adding the continuous potential gradient gives the relative score field used uniformly throughout positive-time spacetime integrals.



Earlier results used here: \cite[Section~3.26.3]{LangevinLean2026ULASC}.
\section{\texorpdfstring{Data processing and a finite-\allowbreak{}entropy change of reference}{Data processing and a finite-entropy change of reference}}
\label{proof-module-UnadjustedLangevin-LSI-EntropyEndpoint-lean}


A Gaussian ULA interpolation applies the same transition to the input law and the Gibbs law. Data processing controls their mutual entropy. A one-sided upper bound on the updated Gibbs density then changes the reference back to the original target, proving finiteness as well as an upper bound. Propagated gradient moments make the logarithmic density correction integrable, without any initial position-moment assumption.



Data processing under the common Gaussian transition gives \(\operatorname{KL}(\mu P_t\mid\nu P_t)\le\operatorname{KL}(\mu\mid\nu)\). The likelihood-ratio decomposition from \(\mu P_t\) to \(\nu\) adds \(\log(d(\nu P_t)/d\nu)\). The explicit density comparison bounds this added term by a constant plus a multiple of \(\|\nabla U\|^2\), whose expectation is finite after propagation. Integrability of both signs of the log likelihood is verified before summing these entropy terms.

\subsection{\texorpdfstring{KL is unchanged by adjoining the same independent law}{KL is unchanged by adjoining the same independent law}}
\label{report-step-UnadjustedLangevin-LSI-EntropyEndpoint-lean-1}


Adjoin the same independent probability innovation \(\eta\) to two finite measures. The log likelihood of \(\mu\otimes\eta\) relative to \(\nu\otimes\eta\) depends only on the first coordinate, so KL is unchanged. The proof uses the standard conditional-product KL identity with the constant Markov kernel \(x\mapsto\eta\); its conditional products are exactly the ordinary product measures. This equality supplies the independent-noise step in data processing.



Taking a product with the same independent probability does not change the likelihood ratio: it depends only on the first coordinate. Integrating out the extra factor therefore preserves the relative entropy.

\subsection{\texorpdfstring{Adding log likelihoods through an intermediate law}{Adding log likelihoods through an intermediate law}}
\label{report-step-UnadjustedLangevin-LSI-EntropyEndpoint-lean-2}


For \(\mu\ll\rho\ll\nu\), the Radon-Nikodym chain rule gives \(d\mu/d\nu=(d\mu/d\rho)(d\rho/d\nu)\) almost everywhere under the integrated law \(\mu\). Each factor is positive and finite there: positivity is transported along absolute continuity, and finiteness uses the probability measures' finite Radon-Nikodym derivatives.



Convert the product of extended densities to real numbers, then apply the logarithm product rule. This proves \(\log(d\mu/d\nu)=\log(d\mu/d\rho)+\log(d\rho/d\nu)\), as an almost-everywhere identity under \(\mu\). The positivity and finiteness checks prevent using logarithm algebra at a zero or infinite density.



For three mutually compatible absolutely continuous laws, their density ratios multiply. Taking logarithms gives the sum of the two log likelihoods on the common full-measure set, before any integration is attempted.

\subsection{\texorpdfstring{An exponential bound on the negative log likelihood}{An exponential bound on the negative log likelihood}}
\label{report-step-UnadjustedLangevin-LSI-EntropyEndpoint-lean-3}


For \(\mu\ll\nu\), the reciprocal likelihood \(e^{-\log(d\mu/d\nu)}\) agrees \(\mu\)-almost everywhere with the real reverse Radon-Nikodym derivative \((d\nu/d\mu)\). This derivative is integrable under \(\mu\), even when \(\nu\) has a singular component relative to \(\mu\). Thus the negative exponential of the log likelihood is always integrable; finite KL has not been assumed at this stage.



The exponential of the negative log likelihood is a reciprocal density ratio. Its integral under the first law is bounded by the reference mass, giving control of the negative log-likelihood tail.

\subsection{\texorpdfstring{Log-\allowbreak{}likelihood integrability from an upper bound}{Log-likelihood integrability from an upper bound}}
\label{report-step-UnadjustedLangevin-LSI-EntropyEndpoint-lean-4}


Suppose the log likelihood \(\ell=\log(d\mu/d\nu)\) has an integrable upper bound \(b\). Where \(\ell\ge0\), \(|\ell|\le b\le|b|\). Where \(\ell<0\), the elementary exponential bound gives \(|\ell|=-\ell\le e^{-\ell}\). Hence globally \(|\ell|\le|b|+e^{-\ell}\). The first term is integrable by hypothesis, and the second by reciprocal-likelihood integrability. Measurable domination proves \(\ell\in L^1(\mu)\), controlling its negative tail without a lower density bound.



An integrable upper bound controls the positive part, while the exponential estimate controls the negative part. Together they establish absolute integrability of the log likelihood needed for the entropy chain rule.

\subsection{\texorpdfstring{Entropy comparison through a reference density bound}{Entropy comparison through a reference density bound}}
\label{report-step-UnadjustedLangevin-LSI-EntropyEndpoint-lean-5}


Assume finite \(H(\mu\mid\rho)\), \(\rho\ll\nu\), and \(d\rho/d\nu\le e^b\), with \(b\in L^1(\mu)\). Finite KL gives \(\mu\ll\rho\), hence \(\mu\ll\nu\). The density bound transfers to \(\mu\)-almost every point, where positivity of the intermediate density permits logarithms. Thus \(\log(d\rho/d\nu)\le b\).



The log-likelihood chain rule gives \(\log(d\mu/d\nu)\le\log(d\mu/d\rho)+b\). Its right side is integrable. The preceding one-sided integrability theorem therefore proves the new log likelihood integrable and hence \(H(\mu\mid\nu)<\infty\). Only after proving finiteness do we integrate the inequality, yielding \(H(\mu\mid\nu)\le H(\mu\mid\rho)+\mathbb E_\mu b\). Finite entropy against the new reference is a conclusion, rather than an additional premise.



Decompose the log likelihood through the intermediate reference law. Its first entropy is controlled by data processing, and the density-ratio upper bound controls the added term by an integrable observable.

\subsection{\texorpdfstring{Data processing for the Gaussian transition}{Data processing for the Gaussian transition}}
\label{report-step-UnadjustedLangevin-LSI-EntropyEndpoint-lean-6}


The interpolation is the measurable image of the initial law times standard Gaussian noise under \((x,z)\mapsto x-t\nabla U(x)+\sqrt{2t}z\). Apply KL data processing to this same map for both initial laws. The independent-noise product identity removes the common Gaussian factor, giving \(H(\rho_t^\mu\mid\rho_t^\nu)\le H(\mu\mid\nu)\). This argument compares two evolved laws and does not assert that either law is invariant under the discrete step.



Both inputs are passed through the same measurable Gaussian update with the same independent noise law. Relative-entropy data processing therefore bounds the two output laws by the divergence of their inputs.



Earlier results used here: \cite[Section~3.36.5]{LangevinLean2026ULASC}; \cite[Section~3.36.7]{LangevinLean2026ULASC}.
\subsection{\texorpdfstring{Propagation of the gradient moment}{Propagation of the gradient moment}}
\label{report-step-UnadjustedLangevin-LSI-EntropyEndpoint-lean-7}


Write \(\nabla U(Y_t)=[\nabla U(Y_t)-\nabla U(X)]+\nabla U(X)\). The initial gradient is \(L^2\), and the frozen-error energy theorem makes the bracketed discrepancy \(L^2\) on the independent initial/\allowbreak{}noise space. Their sum is \(L^2\); the map-measure equivalence transfers this to the actual interpolation law.



Pointwise, the squared norm of the sum is at most twice each squared norm. Integrate and insert \(\mathbb E\|\nabla U(Y_t)-\nabla U(X)\|^2\le L^2(t^2G_0+2td)\). The initial-coordinate marginal gives the remaining \(G_0\), so \(G_t\le2(1+L^2t^2)G_0+4tdL^2\). This proves both integrability and the quantitative estimate without a position moment.



Lipschitzness bounds \(\|\nabla U(Y_t)\|^2\) by \(2\|\nabla U(X)\|^2+2L^2\|Y_t-X\|^2\). The exact displacement moment then yields the displayed propagated energy.



Earlier results used here: \cite[Section~3.34.1]{LangevinLean2026ULASC}; \cite[Section~3.36.5]{LangevinLean2026ULASC}; \cite[Section~3.36.7]{LangevinLean2026ULASC}.
\subsection{\texorpdfstring{The interpolation entropy upper bound}{The interpolation entropy upper bound}}
\label{report-step-UnadjustedLangevin-LSI-EntropyEndpoint-lean-8}


Let \(\rho_t\) be the evolved input and \(\sigma_t\) the evolved Gibbs target. Data processing gives \(H(\rho_t\mid\sigma_t)\le H_0<\infty\). Suppose the actual reference law satisfies \(d\sigma_t/d\nu\le\exp(b+c\|\nabla U\|^2)\), with its absolute continuity retained explicitly.



Finite initial KL first supplies finite \(G_0\) through the direct entropy/\allowbreak{}score estimate. The propagated gradient-moment theorem then proves \(G_t<\infty\). Consequently the logarithmic reference correction \(b+c\|\nabla U\|^2\) is integrable under \(\rho_t\). Apply the finite-entropy change-of-reference theorem to obtain \(H(\rho_t\mid\nu)<\infty\) and \(H(\rho_t\mid\nu)\le H_0+b+cG_t\). The pointwise density comparison itself is established in the Gibbs interpolation module.



Insert the explicit log density-ratio bound and the propagated gradient moment into the entropy decomposition. This gives a finite upper bound on actual interpolation entropy, suitable for sending elapsed time to zero.



Earlier results used here: \cite[Section~3.26.3]{LangevinLean2026ULASC}.
\section{\texorpdfstring{A Gaussian density comparison at the initial endpoint}{A Gaussian density comparison at the initial endpoint}}
\label{proof-module-UnadjustedLangevin-LSI-GibbsInterpolationBound-lean}


The entropy limit at zero is obtained by comparing the updated Gibbs density with the original Gibbs density. This calculation uses the differentiable potential and its Lipschitz gradient directly, with no initial density smoothness or change-of-variables theorem.



For \(0<t<1/(4L)\), set \(\delta=1-4tL>0\) and \(s=t/\delta\). The smoothness inequalities compare the frozen Gaussian exponent starting at \(x\) with a reversed Gaussian exponent centered using \(\nabla U(y)\). Completing the square leaves the correction \(s\|\nabla U(y)\|^2/4\). Integrating the reversed Gaussian contributes the normalizer ratio \((s/t)^{d/2}=\delta^{-d/2}\). The resulting upper bound on \(d(\nu P_t)/d\nu\) is the key initial-endpoint estimate.



Earlier results used here: \cite[Section~3.46]{LangevinLean2026ULASC}.
\subsection{\texorpdfstring{The Gaussian normalizer ratio}{The Gaussian normalizer ratio}}
\label{report-step-UnadjustedLangevin-LSI-GibbsInterpolationBound-lean-1}


For \(t>0\), integration of the radial Gaussian gives \(N_t=(4\pi t)^{-d/2}\). Therefore for positive \(s,t\),

\[N_t/N_s=(s/t)^{d/2}.\]

The proof obtains the individual integrals from the finite-dimensional Gaussian integral, uses positivity to invert them, and applies the multiplication and division identities for real powers.



The normalizer scales as \(t^{-d/2}\), so replacing \(t\) by \(s=t/\delta\) contributes \((s/t)^{d/2}=\delta^{-d/2}\) when the reversed Gaussian is integrated. The direction of this ratio matters for the entropy correction's sign.



Earlier results used here: \cite[Section~3.37.1]{LangevinLean2026ULASC}; \cite[Section~3.37.12]{LangevinLean2026ULASC}; \cite[Section~3.37.2]{LangevinLean2026ULASC}.
\subsection{\texorpdfstring{Complete the square using only first derivatives}{Complete the square using only first derivatives}}
\label{report-step-UnadjustedLangevin-LSI-GibbsInterpolationBound-lean-2}


Write \(g_x=\nabla U(x)\), \(g_y=\nabla U(y)\), and choose positive \(s,t\) with \(s(1-4tL)=t\). The smooth upper model at \(y\) bounds \(U(y)-U(x)\) by \(\langle g_y,y-x\rangle+(L/2)\|y-x\|^2\). Lipschitzness gives \(\|g_x-g_y\|\le L\|x-y\|\). Expanding squares and bounding their cross term yields

\[U(y)-U(x)-\frac{\|y-x+tg_x\|^2}{4t}\le-\frac{\|x-y+sg_y\|^2}{4s}+\frac s4\|g_y\|^2.\]

The relation between \(s\) and \(t\) absorbs the Lipschitz error in the negative quadratic. Exponentiation preserves order. Multiplying by the positive Gaussian normalizer rewrites this as an upper bound on the Gibbs-weighted transition kernel.



Only the first-order smooth upper model and Lipschitz gradient comparison are used. Their quadratic terms are collected into the reversed Gaussian square, leaving the explicit positive correction in the observation-gradient norm.



Earlier results used here: \cite[Section~3.37.12]{LangevinLean2026ULASC}; \cite[Section~3.37.14]{LangevinLean2026ULASC}; \cite[Section~3.9.1]{LangevinLean2026ULASC}.
\subsection{\texorpdfstring{Integrate the comparison against the Gibbs start}{Integrate the comparison against the Gibbs start}}
\label{report-step-UnadjustedLangevin-LSI-GibbsInterpolationBound-lean-3}


The relative density of the updated Gibbs law is

\[r_{\nu,t}(y)=\int N_t\exp\!\left(U(y)-U(x)-\frac{\|y-x+t\nabla U(x)\|^2}{4t}\right)dx.\]

Insert the pointwise comparison. The remaining \(x\)-dependent expression is a normalized Gaussian with time \(s\) centered at \(y-s\nabla U(y)\), so its integral is one. The factors independent of \(x\) give

\[r_{\nu,t}(y)\le(s/t)^{d/2}\exp\!\left(\frac s4\|\nabla U(y)\|^2\right).\]

Normalization of the original Gibbs probability cancels its partition function. Taking \(s=t/(1-4tL)\) later supplies an explicit reference-density bound whose logarithmic cost tends to zero.



After pointwise comparison, integrate the reversed Gaussian in the starting variable. Its total mass is one after its own normalizer is inserted, yielding the density-ratio upper bound needed at the initial entropy endpoint.



Earlier results used here: \cite[Section~3.26.2]{LangevinLean2026ULASC}; \cite[Section~3.26.3]{LangevinLean2026ULASC}; \cite[Section~3.26.9]{LangevinLean2026ULASC}; \cite[Section~3.37.12]{LangevinLean2026ULASC}; \cite[Section~3.37.14]{LangevinLean2026ULASC}; \cite[Section~3.37.6]{LangevinLean2026ULASC}; \cite[Section~3.37.7]{LangevinLean2026ULASC}; \cite[Section~3.46.3]{LangevinLean2026ULASC}.
\section{\texorpdfstring{The initial entropy endpoint from an explicit upper envelope}{The initial entropy endpoint from an explicit upper envelope}}
\label{proof-module-UnadjustedLangevin-LSI-InitialEntropy-lean}


The scalar entropy comparison needs an upper boundary at time zero, even if the initial density is irregular. This module constructs a continuous explicit function \(B(t)\) with \(B(0)=H_0\) and actual interpolation entropy \(H_t\le B(t)\) for sufficiently small positive \(t\). The proof uses the updated Gibbs density bound and data processing; no initial Fisher information, state second moment, or assumed entropy continuity is required.



With \(\delta(t)=1-4tL\) and \(G_0=\int\|\nabla U\|^2d\mu\), the endpoint majorant is \(M(t)=H(0)-(d/2)\log\delta(t)+t[2(1+L^2t^2)G_0+4tdL^2]/(4\delta(t))\). It is defined in a positive neighborhood where \(\delta(t)>0\), satisfies \(M(0)=H(0)\), and is continuous there. The bound \(H(t)\le M(t)\) therefore supplies the upper endpoint limit needed by scalar comparison, without requiring differentiability or even a prescribed entropy density at zero.

\subsection{\texorpdfstring{The explicit initial-\allowbreak{}endpoint majorant}{The explicit initial-endpoint majorant}}
\label{report-step-UnadjustedLangevin-LSI-InitialEntropy-lean-1}


Put \(\delta(t)=1-4tL\) and \(G_0=\mathbb E_\mu\|\nabla U\|^2\). Define the explicit envelope \(B(t)=H_0-\frac d2\log\delta(t)+\frac{t}{4\delta(t)}[2(1+L^2t^2)G_0+4tdL^2]\). The logarithmic term is the normalizer correction from the updated Gibbs density; the final term uses the propagated gradient moment. The definition is a scalar function at all times, while its entropy interpretation below requires \(t>0\) and \(\delta(t)>0\).



The majorant combines the data-processing entropy with the Gaussian normalizer correction and a gradient-energy correction. Its explicit denominator is positive for all sufficiently small positive times, so the expression is an ordinary finite scalar there.



Earlier results used here: \cite[Section~3.26.3]{LangevinLean2026ULASC}.
\subsection{\texorpdfstring{The majorant at zero}{The majorant at zero}}
\label{report-step-UnadjustedLangevin-LSI-InitialEntropy-lean-2}


At \(t=0\), the logarithm has argument \(\delta(0)=1\), hence is zero. The coefficient \(t/(4\delta(t))\) is zero as well. Both correction terms vanish exactly, proving \(B(0)=H_0\). This is a direct algebraic evaluation of the constructed envelope, not an assumed limit of the actual entropy path.



At zero the logarithmic term is \(\log1=0\), and the entire energy correction has a factor \(t\). Hence the value is precisely the initial entropy, even when the initial law has no smooth density.



Earlier results used here: \cite[Section~3.26.3]{LangevinLean2026ULASC}.
\subsection{\texorpdfstring{Continuity of the majorant at zero}{Continuity of the majorant at zero}}
\label{report-step-UnadjustedLangevin-LSI-InitialEntropy-lean-3}


The denominator \(\delta(t)\) has value one at zero, so it stays nonzero locally. The logarithm is continuous at that same positive value. All other terms in \(B\) are polynomial or constant scalar operations. Their continuity proves \(B(t)\to B(0)=H_0\) as \(t\to0\). The initial gradient integral is an already defined scalar; its genuine finiteness in the entropy application is established by the finite-KL moment theorem.



The denominator tends to one and is nonzero near zero. Continuity of the logarithm at one, polynomial factors, and reciprocal then proves continuity of the full majorant at the endpoint.

\subsection{\texorpdfstring{The actual entropy lies below the majorant}{The actual entropy lies below the majorant}}
\label{report-step-UnadjustedLangevin-LSI-InitialEntropy-lean-4}


For positive \(t\) with \(\delta=1-4tL>0\), choose the comparison Gaussian time \(s=t/\delta\). The updated Gibbs law has an explicit relative density bounded by \((s/t)^{d/2}\exp(s\|\nabla U(y)\|^2/4)\). Since \(s/t=\delta^{-1}\), this is \(\exp[-(d/2)\log\delta+(s/4)\|\nabla U(y)\|^2]\).



The actual relative density identifies the Radon-Nikodym derivative and proves absolute continuity of the updated Gibbs reference law. The change-of-reference entropy estimate then gives finite \(H_t\) and \(H_t\le H_0-(d/2)\log\delta+(s/4)G_t\). The coefficient \(s/4\) is nonnegative, so inserting \(G_t\le2(1+L^2t^2)G_0+4tdL^2\) proves \(H_t\le B(t)\).



The density-ratio comparison bounds actual output entropy by the majorant whenever \(1-4tL>0\). The finite score moment supplies all integrability required for the likelihood decomposition.



Earlier results used here: \cite[Section~3.26.3]{LangevinLean2026ULASC}.
\subsection{\texorpdfstring{The entropy majorant on a positive neighborhood}{The entropy majorant on a positive neighborhood}}
\label{report-step-UnadjustedLangevin-LSI-InitialEntropy-lean-5}


For all sufficiently small positive \(t\), specifically \(t<1/(8L)\), one has \(1-4tL>0\). The previous positive-time entropy bound therefore applies eventually in the right neighborhood of zero, giving \(H_t\le B(t)\). Combined with continuity of \(B\) and \(B(0)=H_0\), this supplies the upper initial boundary used in scalar comparison. It does not insert an assumption that \(H_t\) is differentiable or continuous at zero.



A sufficiently small positive neighborhood satisfies the denominator condition. Thus the entropy upper bound is available along every approach to zero through positive times, exactly the endpoint control used in scalar comparison.



Earlier results used here: \cite[Section~3.26.3]{LangevinLean2026ULASC}.
\section{\texorpdfstring{The conditional one-\allowbreak{}step KL contraction with bias 6}{The conditional one-step KL contraction with bias 6}}
\label{proof-module-UnadjustedLangevin-LSI-OneStep-lean}


Assume the displayed LSI, gradient-moment and interpolation differential bounds. For \(0<h\le\alpha/(4L^2)\), a scalar integrating-factor argument proves \(H_{k+1}\le e^{-\alpha h}H_k+6h^2dL^2\).



Write \(x=\alpha h\). The step restriction and compatibility \(\alpha\le L\) imply \(0<x\le1/4\), \(hL\le1/4\), and \(h^2L^2(4L^2/\alpha)\le\alpha/4\). The scalar differential inequality therefore has uniform forcing \(\alpha H(0)/4+(5/2)hdL^2\). Its initial-entropy part is absorbed by the quantitative gap between \(e^{-3x/2}\) and \(e^{-x}\), while the remaining forcing is of order \(h^2dL^2\).

\subsection{\texorpdfstring{Absorb a small initial-\allowbreak{}entropy forcing into the exponential gap}{Absorb a small initial-entropy forcing into the exponential gap}}
\label{report-step-UnadjustedLangevin-LSI-OneStep-lean-1}


For \(0\le x\le1/4\), the tangent bound for the exponential gives \(e^{-3x/2}\ge1/2\). Multiply \(e^{x/2}\ge1+x/2\) by \(e^{-3x/2}\); then \(e^{-x}\ge e^{-3x/2}+(x/2)e^{-3x/2}\ge e^{-3x/2}+x/4\). This quantitative rate gap absorbs the initial-entropy forcing.



The comparison retains a deliberate gap between the stronger raw decay \(e^{-3x/2}\) and the final factor \(e^{-x}\). The lower bound \(x/4\) on this gap is precisely large enough to absorb the initial-entropy-dependent forcing generated by the score moment.

\subsection{\texorpdfstring{Uniform forcing, affine comparison and the coefficient 6}{Uniform forcing, affine comparison and the coefficient 6}}
\label{report-step-UnadjustedLangevin-LSI-OneStep-lean-2}


Combine \(2\alpha H\le I\) and \(H'\le-3I/4+t^2L^2G_0+2tdL^2\). The gradient bound \(G_0\le(4L^2/\alpha)H(0)+2dL\) and \(h\le\alpha/(4L^2)\) give a uniform inequality \(H'\le-(3\alpha/2)H+b\), with \(b=(\alpha/4)H(0)+(5/2)hdL^2\). The proof establishes \(hL\le1/4\), \(\alpha h\le1/4\), and the coefficient bound \(h^2L^2(4L^2/\alpha)\le\alpha/4\) before this substitution.



Apply the scalar affine differential-inequality comparison and bound its forcing integral by \(bh\), using \(1-e^{-ah}\le ah\). The exponential-gap lemma absorbs \((\alpha h/4)H(0)\), changing the decay factor from \(e^{-3\alpha h/2}\) to \(e^{-\alpha h}\). Nonnegativity of \(H(0)\) justifies multiplying that scalar bound. Finally enlarge the remaining nonnegative \((5/2)h^2dL^2\) term to \(6h^2dL^2\), obtaining the stated one-step estimate.



The dimensional forcing integrates over one time interval and therefore gains another factor \(h\). After absorbing the initial-entropy contribution, the remaining bound is \((5/2)h^2dL^2\), which is enlarged to the uniform coefficient six used in the recurrence. All simplifications use the positive step and curvature assumptions.

\section{\texorpdfstring{Integrating entropy dissipation with almost-\allowbreak{}everywhere derivatives}{Integrating entropy dissipation with almost-everywhere derivatives}}
\label{proof-module-UnadjustedLangevin-LSI-IntegratedEntropy-lean}


The actual entropy chain rule yields local absolute continuity and a derivative inequality almost everywhere. These scalar lemmas use precisely that regularity, together with a continuous one-sided upper envelope at time zero. They do not require pointwise derivatives at every time.



On each \([a,b]\subset(0,h]\), local absolute continuity and an almost-everywhere derivative inequality suffice for the integrating-factor argument. The resulting bound uses the entropy at \(a\), not an assumed derivative at zero. An explicit upper majorant continuous at zero then allows \(a\downarrow0\). Thus finite initial entropy can enter the comparison even when the initial law has no smooth positive density; all spatial smoothing is used only after a positive elapsed time.

\subsection{\texorpdfstring{Integrate the affine differential inequality}{Integrate the affine differential inequality}}
\label{report-step-UnadjustedLangevin-LSI-IntegratedEntropy-lean-1}


Suppose \(H\) is absolutely continuous on \([s,t]\) and \(H^{\prime}\le-aH+b\) almost everywhere, with \(a>0\). Set \(F(u)=e^{au}(H(u)-b/a)\). Products of absolutely continuous functions on a compact interval are absolutely continuous. The product rule holds almost everywhere and gives \(F^{\prime}(u)=e^{au}(a(H(u)-b/a)+H^{\prime}(u))\le0\).



The fundamental theorem of calculus for absolutely continuous functions gives \(F(t)-F(s)=\int_s^tF^{\prime}\le0\). Multiply by \(e^{-at}>0\), cancel the exponentials, and rearrange:

\[H(t)\le e^{-a(t-s)}H(s)+\frac ba(1-e^{-a(t-s)}).\]

The proof removes no exceptional set by assuming a derivative there; both differentiability and the inequality are used only almost everywhere.



Local absolute continuity gives an integrable derivative and the fundamental theorem on each positive compact interval. Multiplying by the smooth integrating factor preserves absolute continuity, so integrating its almost-everywhere derivative bound is legitimate.

\subsection{\texorpdfstring{Send the positive initial time to zero}{Send the positive initial time to zero}}
\label{report-step-UnadjustedLangevin-LSI-IntegratedEntropy-lean-2}


Suppose \(H\) is absolutely continuous on every \([s,T]\), \(0<s<T\), and \(H(s)\le M(s)\) for all sufficiently small positive \(s\). If \(M\) is continuous at zero, the preceding estimate gives

\[H(T)\le e^{-a(T-s)}M(s)+\frac ba(1-e^{-a(T-s)}).\]

The right side tends to \(e^{-aT}M(0)+(b/a)(1-e^{-aT})\). Order is preserved under limits, giving the desired initial-endpoint bound. Single endpoints have zero Lebesgue measure, so an inequality on the open interval suffices. Continuity or differentiability of \(H\) at zero is not needed.



The upper initial majorant converges to the actual finite initial entropy. Applying it at the lower positive endpoint and sending that endpoint to zero therefore bounds the limiting comparison term from above, which is all the scalar inequality requires.

\subsection{\texorpdfstring{Absorb the ULA forcing with its exact constants}{Absorb the ULA forcing with its exact constants}}
\label{report-step-UnadjustedLangevin-LSI-IntegratedEntropy-lean-3}


Assume \(2\alpha H\le I\), \(H^{\prime}\le-3I/4+t^2L^2G_0+2tdL^2\), and \(G_0\le4L^2M(0)/\alpha+2dL\). Let \(0<h\le\alpha/(4L^2)\), \(0<\alpha\le L\), and \(M(0)\ge0\). Then \(hL\le1/4\), \(\alpha h\le1/4\),

\[h^2L^2(4L^2/\alpha)\le\alpha/4,\qquad h^2L^2(2dL)\le hdL^2/2.\]

The derivative is therefore at most \(-3\alpha H/2+B\), where \(B=\alpha M(0)/4+(5/2)hdL^2\). The affine comparison and \(1-e^{-ah}\le ah\) imply \(H(h)\le e^{-3\alpha h/2}M(0)+Bh\). Finally \(e^{-x}\ge e^{-3x/2}+x/4\) for \(0\le x\le1/4\) absorbs the entropy part of \(Bh\). Bounding \(5/2\) by 6 gives

\[H(h)\le e^{-\alpha h}M(0)+6h^2dL^2.\]

The final actual-law theorem proves every analytic premise before invoking this scalar result.



The time-dependent forcing is first bounded uniformly using \(t\le h\). The initial-entropy contribution is absorbed by the exponential gap, while the dimensional forcing integrates to a multiple of \(h^2dL^2\), giving the stated one-step coefficient.

\section{\texorpdfstring{Iterate entropy contraction and obtain the bias constant 8}{Iterate entropy contraction and obtain the bias constant 8}}
\label{proof-module-UnadjustedLangevin-LSI-DiscreteEntropy-lean}


Starting from \(H_{k+1}\le e^{-\alpha h}H_k+6h^2dL^2\), induction retains the exact geometric sum. The small-step denominator estimate \(1-e^{-\alpha h}\ge\frac34\alpha h\) gives \(H_k\le e^{-\alpha hk}H_0+8hdL^2/\alpha\).



After one step the entropy recurrence is \(H_{k+1}\le qH_k+c\), where \(q=e^{-\alpha h}\in(0,1)\) and \(c=6h^2dL^2\). Iteration gives \(q^kH_0+c(1-q^k)/(1-q)\). For \(\alpha h\le1/4\), the lower bound \(1-q\ge3\alpha h/4\) converts the accumulated \(O(h^2)\) local forcing into the global floor \(8hdL^2/\alpha\). Positivity of the denominator is essential when making this comparison.



Earlier results used here: \cite[Section~3.15]{LangevinLean2026ULASC}.
\subsection{\texorpdfstring{The exact finite geometric entropy sum}{The exact finite geometric entropy sum}}
\label{report-step-UnadjustedLangevin-LSI-DiscreteEntropy-lean-1}


Set \(q=e^{-\alpha h}\in(0,1)\), \(b=6h^2dL^2\). The affine recurrence theorem sums \(H_{k+1}\le qH_k+b\) into \(H_k\le q^kH_0+b(1-q^k)/(1-q)\). Rewrite \(q^k=e^{-\alpha hk}\) using the exponential natural-multiple identity. No sign assumption on \(H_0\) is needed.



The exact term \((1-q^k)/(1-q)\) records that the forcing acts only \(k\) times. Writing \(q^k=e^{-\alpha hk}\) identifies physical elapsed time as \(hk\), preserving the correct rate in the final statement.



Earlier results used here: \cite[Section~3.15.2]{LangevinLean2026ULASC}.
\subsection{\texorpdfstring{A lower bound for the geometric denominator}{A lower bound for the geometric denominator}}
\label{report-step-UnadjustedLangevin-LSI-DiscreteEntropy-lean-2}


The tangent bound \(e^x\ge1+x\), multiplied by \(e^{-x}\), gives \(1-e^{-x}\ge x/(1+x)\). For \(0\le x\le1/4\), elementary positive-denominator algebra shows \(x/(1+x)\ge3x/4\). This lower bound controls the geometric-series denominator in the bias floor.



The tangent inequality gives \(e^{-x}\le1/(1+x)\), so its complement is at least \(x/(1+x)\). Since \(x\le1/4\), the denominator is at most \(5/4\); the weaker coefficient \(3/4\) is therefore valid and sufficient.

\subsection{\texorpdfstring{The iterated entropy bound and its bias floor}{The iterated entropy bound and its bias floor}}
\label{report-step-UnadjustedLangevin-LSI-DiscreteEntropy-lean-3}


The nonnegative forcing gives \(H_k\le e^{-\alpha hk}H_0+6h^2dL^2/(1-e^{-\alpha h})\). For \(\alpha h\le1/4\), the denominator is at least \(3\alpha h/4>0\). Divide and cancel to obtain the bias floor \(8hdL^2/\alpha\). The denominator is proved strictly positive before multiplication or division.



Substituting the positive lower bound on \(1-q\) gives \(6h^2dL^2/(3\alpha h/4)=8hdL^2/\alpha\). One power of the step cancels, leaving the expected first-order entropy bias floor.



Earlier results used here: \cite[Section~3.15.3]{LangevinLean2026ULASC}.
\section{\texorpdfstring{An integrated chain rule from spacetime integrability}{An integrated chain rule from spacetime integrability}}
\label{proof-module-UnadjustedLangevin-LSI-EntropyChainRule-lean}


To differentiate entropy, bounded approximations are first integrated in time. A pointwise time derivative together with joint spacetime \(L^1\) suffices for Fubini and the fundamental theorem of calculus; no common spatial majorant for every time derivative is assumed. The resulting spatial integral is absolutely continuous and has the integrated derivative almost everywhere. A dominated-limit theorem then passes truncated integral identities to the actual entropy identity.



For a spacetime-integrable derivative field \(F(t,x)\), Fubini gives an absolutely continuous time primitive \(A(t)=A(a)+\int_a^t\int F(s,x)d\nu(x)ds\). The fundamental theorem of calculus identifies its derivative almost everywhere with the inner integral. If bounded approximations have convergent endpoint values and a common integrable flux majorant, dominated convergence transfers their integrated identities to the limit. This obtains entropy absolute continuity from an integrated identity rather than assuming pointwise differentiability at every time.

\subsection{\texorpdfstring{An integrated identity from a spacetime derivative}{An integrated identity from a spacetime derivative}}
\label{report-step-UnadjustedLangevin-LSI-EntropyChainRule-lean-1}


Assume \(f(a,\cdot)\in L^1(\nu)\), \(\dot f\in L^1((a,b]\times\nu)\), and \(\partial_t f(t,x)=\dot f(t,x)\) for every \(t\in[a,b]\) and almost every \(x\). Fubini gives \(\dot f(\cdot,x)\in L^1(a,b)\) for almost every \(x\). Apply the scalar fundamental theorem to these sections: \(f(b,x)=f(a,x)+\int_a^b\dot f(t,x)dt\).



The time-integrated derivative is integrable in \(x\) by spacetime integrability. The section identity therefore proves terminal integrability of \(f(b,\cdot)\), rather than requiring it separately. Integrate the identity and exchange time and space to obtain \(\int f(b)d\nu=\int f(a)d\nu+\int_a^b\int\dot f(t,x)d\nu(x)dt\).



For almost every spatial point, integrate its time derivative between the endpoints. Spacetime absolute integrability permits Fubini to exchange the two integrations, yielding the integrated identity for the spatial integral.

\subsection{\texorpdfstring{Absolute continuity from spacetime integrability}{Absolute continuity from spacetime integrability}}
\label{report-step-UnadjustedLangevin-LSI-EntropyChainRule-lean-2}


Apply the integrated fundamental theorem on every subinterval \([a,t]\subseteq[a,b]\). Restricting the spacetime measure preserves integrability, giving both integrability of \(f(t,\cdot)\) and \(F(t)=F(a)+\int_a^t q(u)du\), where \(F(t)=\int f(t,x)d\nu(x)\) and \(q(t)=\int\dot f(t,x)d\nu(x)\).



Fubini gives \(q\in L^1(a,b)\). An indefinite integral of an \(L^1\) function is absolutely continuous; adding the constant \(F(a)\) preserves this property. Equality on the interval transfers absolute continuity to \(F\). Thus regularity of the spatial integral is a conclusion of spacetime integrability and pointwise differentiation.



An integral of an \(L^1\) time function is absolutely continuous. The spacetime norm bound makes the spatially integrated derivative such an \(L^1\) function, supplying this regularity without an assumed entropy derivative.

\subsection{\texorpdfstring{The almost-\allowbreak{}everywhere derivative of the integrated quantity}{The almost-everywhere derivative of the integrated quantity}}
\label{report-step-UnadjustedLangevin-LSI-EntropyChainRule-lean-3}


The subinterval argument gives \(F(t)=F(a)+\int_a^t q(u)du\) for every \(t\in[a,b]\), with \(q\in L^1(a,b)\). Lebesgue differentiation supplies derivative \(q(t)\) for the indefinite integral at almost every interior time. For each such \(t\), an open neighborhood stays inside \((a,b)\), where the displayed identity holds. The local equality transfers the derivative to \(F\), yielding \(F'(t)=\int\dot f(t,x)d\nu(x)\) almost everywhere. No pointwise derivative assertion at every time, or at the endpoints, is inferred.



The fundamental theorem for absolutely continuous primitives gives the derivative almost everywhere. The exceptional time set has measure zero, which is sufficient for all later integrating-factor inequalities.

\subsection{\texorpdfstring{Passing an integrated identity through dominated limits}{Passing an integrated identity through dominated limits}}
\label{report-step-UnadjustedLangevin-LSI-EntropyChainRule-lean-4}


Suppose truncated identities have the form \(B_n-A_n=\int q_n d\nu\). Assume the two endpoint sequences converge to their intended values, \(q_n\to q\) almost everywhere, and \(|q_n|\le g\) with \(g\in L^1(\nu)\). Dominated convergence gives convergence of the right-hand integrals, while the endpoint limits give convergence of the left-hand differences. Their equality for every \(n\) and uniqueness of limits prove \(B-A=\int qd\nu\).



In the entropy application, the endpoints are truncated entropy integrals and \(q_n\) are bounded-test fluxes on a fixed time interval. Spatial energy estimates supply the pointwise flux limit and a time-integrable dominating bound. This lemma passes the integrated identities to the limit without assuming in advance that the full entropy has already been differentiated.



Pass the approximating integrated identities to the limit in both endpoints and flux. Endpoint domination and spacetime flux domination are checked separately because they use different reference measures and majorants.

\subsection{\texorpdfstring{Regularity supplied by the limiting integral identity}{Regularity supplied by the limiting integral identity}}
\label{report-step-UnadjustedLangevin-LSI-EntropyChainRule-lean-5}


Assume only \(q\in L^1(a,b)\) and the exact identity \(H(t)=H(a)+\int_a^t q(u)du\) throughout \([a,b]\). The primitive of \(q\) is absolutely continuous; adding \(H(a)\) and using the identity proves that \(H\) is absolutely continuous.



Lebesgue differentiation gives derivative \(q(t)\) for the primitive at almost every interior time. The identity holds in an interior neighborhood of each such time, so the derivative transfers to \(H\). This gives precisely the final entropy comparison inputs: local absolute continuity and the actual flux as derivative almost everywhere.



The limiting identity expresses the function as a constant plus an integral of an \(L^1\) flux. It therefore supplies both local absolute continuity and its almost-everywhere derivative, closing the regularity needed for entropy comparison.

\section{\texorpdfstring{Bounded logarithmic tests and their entropy primitives}{Bounded logarithmic tests and their entropy primitives}}
\label{proof-module-UnadjustedLangevin-LSI-BoundedLog-lean}


The unbounded logarithm in entropy production is replaced by a smooth bounded function. The replacement has derivative between zero and one, converges to the identity, and admits an entropy primitive dominated by the actual KL integrand. These properties separately justify the spatial test, the time chain rule, and the limit in the entropy identity.



For \(M=m+1\), use \(\chi_m(z)=M\arctan(z/M)\), with derivative \(w_m(z)=(1+(z/M)^2)^{-1}\). The function \(\chi_m\) is bounded, \(0\le w_m\le1\), and for fixed \(z\), \(\chi_m(z)\to z\) and \(w_m(z)\to1\). The primitive \(\Phi_m(r)=\int_1^r\chi_m(\log u)du\) increases toward the entropy integrand \(r\log r-r+1\). These bounds separate admissibility of each test from the limiting unbounded entropy identity.

\subsection{\texorpdfstring{Choose the bounded test and its weight}{Choose the bounded test and its weight}}
\label{report-step-UnadjustedLangevin-LSI-BoundedLog-lean-1}


For a natural number \(m\), write \(c=m+1>0\). Define

\[\chi_m(z)=c\arctan(z/c),\qquad w_m(z)=\frac1{1+(z/c)^2},\qquad\Phi_m(r)=\int_1^r\chi_m(\log u)du.\]

The positive offset in \(c\) prevents division by zero even for \(m=0\). The interval integral is oriented: when \(r<1\), its sign reverses. The formulas define these expressions for every real argument; the entropy claims below explicitly require \(r>0\).



The scale \(M=m+1\) is strictly positive for every allowed index, avoiding a zero denominator. The arctangent bounds the test uniformly in its argument, while the derivative weight remains between zero and one.

\subsection{\texorpdfstring{Differentiate and bound the test}{Differentiate and bound the test}}
\label{report-step-UnadjustedLangevin-LSI-BoundedLog-lean-2}


The arctangent chain rule gives \(\chi_m(0)=0\) and \(\chi_m^{\prime}(z)=w_m(z)\). The denominator is at least one, so \(0\le w_m\le1\). The derivative bound implies that \(\chi_m\) is 1-Lipschitz, hence \(|\chi_m(z)|\le|z|\). Independently, \(|\arctan u|<\pi/2\) gives \(|\chi_m(z)|\le c\pi/2\). This last bound controls the time derivative integrand without any log-integrability requirement. Since arctangent has the sign of its argument, \(\chi_m\) does too.



smoothness of all orders records smoothness of the scalar test. the scalar derivative identity gives its exact derivative; the proof multiplies the derivative of arctangent by \(1/c\) and cancels the outer factor \(c\).



Differentiating arctangent gives the rational weight \(1/(1+(z/M)^2)\). Its denominator is at least one, proving nonnegativity and the upper bound used to dominate every truncated score energy.

\subsection{\texorpdfstring{Recover the logarithm and its primitive}{Recover the logarithm and its primitive}}
\label{report-step-UnadjustedLangevin-LSI-BoundedLog-lean-3}


As \(m\to\infty\), \(z/(m+1)\to0\). The derivative of arctangent at zero is one; its difference quotient therefore gives \(\chi_m(z)\to z\). At \(z=0\) the claim is exact at every index. Continuity of inversion at the positive denominator one gives \(w_m(z)\to1\).



The function \(u\mapsto\chi_m(\log u)\) is measurable and bounded, hence integrable on every finite interval, including intervals containing zero. When \(r>0\), continuity in a neighborhood of \(r\) and the fundamental theorem of calculus imply \(\Phi_m^{\prime}(r)=\chi_m(\log r)\).



For fixed \(z\), the ratio \(z/M\) tends to zero, and \(\arctan u/u\to1\). This gives recovery of the logarithmic test and, after integration in the density variable, its entropy primitive.

\subsection{\texorpdfstring{Dominate the entropy endpoints}{Dominate the entropy endpoints}}
\label{report-step-UnadjustedLangevin-LSI-BoundedLog-lean-4}


We have \(\Phi_m(1)=0\). For \(r\ge1\), both \(\log u\) and its test are nonnegative on \([1,r]\), giving \(\Phi_m(r)\ge0\). For \(0<r\le1\), the test is nonpositive on \([r,1]\), and reversing the interval again gives nonnegativity.



On the positive side, \(0\le\chi_m(\log u)\le\log u\). On the negative side, \(\log u\le\chi_m(\log u)\le0\); reversing the interval preserves the corresponding upper bound for the primitive. In both cases,

\[0\le\Phi_m(r)\le\int_1^r\log u\,du=r\log r-r+1.\]

Pointwise convergence of the test and domination by \(|\log u|\), integrable on the finite interval between 1 and \(r\), prove \(\Phi_m(r)\to r\log r-r+1\). These are exactly the hypotheses for dominated convergence against the finite-KL relative density later.



The entropy integrand \(r\log r-r+1\) is nonnegative and dominates the bounded primitives. For a probability density its integral equals KL, so finite actual entropy supplies the endpoint majorant for dominated convergence.

\subsection{\texorpdfstring{Identify the spatial gradient of the test}{Identify the spatial gradient of the test}}
\label{report-step-UnadjustedLangevin-LSI-BoundedLog-lean-5}


For a differentiable scalar field \(f\), the Fréchet chain rule gives \(D(\chi_m\circ f)=w_m(f)Df\). Applying the Hilbert-space identification of linear derivatives with gradients yields

\[\nabla(\chi_m\circ f)(x)=w_m(f(x))\nabla f(x),\qquad\|\nabla(\chi_m\circ f)(x)\|\le\|\nabla f(x)\|.\]

The inequality uses \(0\le w_m\le1\). With \(f=\log r_t\), it supplies square-integrability of the bounded log test gradient from the actual Fisher bound.



Applying the spatial chain rule multiplies the relative score by \(w_m(\log r)\). Since this weight is at most one, the gradient's squared energy is bounded by genuine Fisher information, proving admissibility of every bounded test.

\section{\texorpdfstring{Derive the actual Gaussian weak equation for bounded finite-\allowbreak{}energy tests}{Derive the actual Gaussian weak equation for bounded finite-energy tests}}
\label{proof-module-UnadjustedLangevin-LSI-GaussianWeakEquation-lean}


Write \(g(x)=\nabla U(x)\), \(c_t(x)=x-tg(x)\), and \(K_t(x,y)=N_t e^{-\|y-c_t(x)\|^2/(4t)}\). The spatial flux \(F_{t,x}(y)=K_t(x,y)(g(x)-(y-c_t(x))/(2t))\) has divergence equal to the actual Gaussian kernel time derivative. Global integration by parts therefore gives \(\int\dot p_t\psi=-\int F_{t,x}\cdot\nabla\psi\,d\mu\,dy\) for every bounded \(C^1\) test with \(\nabla\psi\in L^2(\rho_t)\). This module proves all integrability and boundary-product obligations from finite initial gradient energy and Gaussian second moments. It differentiates no Hessian of \(U\) and assumes no weak equation.



For the frozen component define \(F_t(x,y)=\nabla_yK_t(x,y)+K_t(x,y)\nabla U(x)\). Direct differentiation gives \(\operatorname{div}_yF_t=\dot K_t\), with this sign corresponding to \(\partial_tp=\Delta p+\operatorname{div}(p\,\text{frozen gradient})\). Coordinate integration by parts yields \(\int\dot p_t\psi=-\iint\langle F_t,\nabla\psi\rangle d\mu\,dy\). Gaussian energy domination supplies the required flux and derivative integrability for bounded \(C^1\) tests with finite gradient energy.



Earlier results used here: \cite[Section~3.38]{LangevinLean2026ULASC}.
\subsection{\texorpdfstring{Choose the spatial flux with the correct Fokker–Planck sign}{Choose the spatial flux with the correct Fokker–Planck sign}}
\label{report-step-UnadjustedLangevin-LSI-GaussianWeakEquation-lean-1}


The component density evolves under frozen drift \(-g(x)\). Its time derivative is \(\Delta_yK_t+\operatorname{div}_y(K_tg(x))\), so define

\[F_{t,x}(y)=\nabla_yK_t(x,y)+K_t(x,y)g(x)=K_t(x,y)\left(g(x)-\frac{y-c_t(x)}{2t}\right).\]

frozen gaussian flux is the last expression. Its divergence equals the density's time derivative, while the weak equation consequently carries a minus sign after integration by parts. The frozen gradient \(g(x)\) is constant in the spatial variable \(y\), which avoids differentiating the possibly merely Lipschitz gradient field.



The flux uses the frozen starting gradient \(\nabla U(x)\), not the current gradient at the observation. Its divergence equals the kernel time derivative, while the difference from the target drift later appears as entropy forcing.

\subsection{\texorpdfstring{Differentiate each coordinate of the flux explicitly}{Differentiate each coordinate of the flux explicitly}}
\label{report-step-UnadjustedLangevin-LSI-GaussianWeakEquation-lean-2}


Let \(r=y-c_t(x)\). For a direction \(v\), differentiating the kernel and the affine residual gives

\[D_v\langle F_{t,x}(y),v\rangle=K_t(x,y)\left(\frac{\langle r,v\rangle^2}{4t^2}-\frac{\langle r,v\rangle\langle g(x),v\rangle}{2t}-\frac{\|v\|^2}{2t}\right).\]

Differentiate each flux component using the Gaussian gradient and product rules. The resulting diagonal quadratic displacement and drift terms are the coordinate contributions to its divergence.

\subsection{\texorpdfstring{Sum the directional derivatives to obtain a coordinate-\allowbreak{}free divergence}{Sum the directional derivatives to obtain a coordinate-free divergence}}
\label{report-step-UnadjustedLangevin-LSI-GaussianWeakEquation-lean-3}


Sum over an orthonormal basis \((b_i)_{i=1}^d\). Parseval gives \(\sum_i\langle r,b_i\rangle^2=\|r\|^2\), completeness gives \(\sum_i\langle r,b_i\rangle\langle g,b_i\rangle=\langle r,g\rangle\), and \(\sum_i\|b_i\|^2=d\). Thus

\[\operatorname{div}F_{t,x}=K_t\left(-\frac d{2t}+\frac{\|r\|^2}{4t^2}-\frac{\langle r,g\rangle}{2t}\right)=K_t\left(-\frac d{2t}+\frac{\|y-x\|^2}{4t^2}-\frac{\|g\|^2}{4}\right).\]

Summing over an orthonormal coordinate basis collects displacement squares into a Euclidean squared norm and diagonal constants into the dimension. This gives a coordinate-free divergence formula.

\subsection{\texorpdfstring{Transport the actual independent pair to its Gaussian Lebesgue density}{Transport the actual independent pair to its Gaussian Lebesgue density}}
\label{report-step-UnadjustedLangevin-LSI-GaussianWeakEquation-lean-4}


The pair \((X,Y_t)\), constructed from independent initial state and standard noise, satisfies

\[\mathcal L(X,Y_t)(dx,dy)=K_t(x,y)\,\mu(dx)\,dy.\]

The identity for the conditional Gaussian starting-proposal pair-density formula already proves this transport for a measurable center map. Apply it to \(x\mapsto x-t\nabla U(x)\), whose measurability follows from the Lipschitz gradient. It then checks that the real kernel here is exactly the previously normalized proposal density. This identity is essential: it lets the proof establish integrability under a simple independent probability law and transfer the result to the Lebesgue pair integrals needed for integration by parts.



The actual independent state-noise pair is pushed to the starting-proposal pair. Its conditional Gaussian density is the kernel, so nonnegative and integrable pair expectations have the stated Lebesgue representation.



Earlier results used here: \cite[Section~3.34.1]{LangevinLean2026ULASC}; \cite[Section~3.37.12]{LangevinLean2026ULASC}; \cite[Section~3.37.13]{LangevinLean2026ULASC}; \cite[Section~3.38.1]{LangevinLean2026ULASC}.
\subsection{\texorpdfstring{Prove flux pairings integrable by transporting two L\ensuremath{{}^{2}} factors}{Prove flux pairings integrable by transporting two L² factors}}
\label{report-step-UnadjustedLangevin-LSI-GaussianWeakEquation-lean-5}


For a continuous scalar field \(\varphi\in L^2(\rho_t)\), the coordinate flux pairing is integrable:

\[\int K_t(x,y)\left|\left\langle g(x)-\frac{y-c_t(x)}{2t},v\right\rangle\varphi(y)\right|\,d\mu(x)\,dy<\infty.\]

The flux pairing can be transported to a product of square-integrable fields on the independent-noise source. Cauchy–Schwarz proves absolute integrability there and hence in the kernel representation.



Earlier results used here: \cite[Section~3.34.1]{LangevinLean2026ULASC}; \cite[Section~3.36.5]{LangevinLean2026ULASC}; \cite[Section~3.37.4]{LangevinLean2026ULASC}.
\subsection{\texorpdfstring{Admit the gradient pairing of a C\ensuremath{{}^{1}} finite-\allowbreak{}energy test}{Admit the gradient pairing of a C¹ finite-energy test}}
\label{report-step-UnadjustedLangevin-LSI-GaussianWeakEquation-lean-6}


Choose the scalar field \(\varphi(y)=D\psi(y)[v]=\langle\nabla\psi(y),v\rangle\). For \(\psi\in C^1\), this field is continuous. If \(\nabla\psi\in L^2(\rho_t)\), its coordinate pairing is also in \(L^2\).



For a fixed direction \(v\), the derivative of the \(C^1\) test is \(D\psi(y)[v]=\langle\nabla\psi(y),v\rangle\), a continuous function. Its absolute value is bounded by \(\|v\|\|\nabla\psi(y)\|\), so it has finite square energy. The flux-pairing estimate then gives the required integrability of \(F_i\partial_i\psi\).



A bounded \(C^1\) test with finite gradient square energy lies in the intended weak-test class. The preceding flux estimate verifies integrability of its actual gradient pairing before integration by parts.

\subsection{\texorpdfstring{Control the boundary product from boundedness of the test}{Control the boundary product from boundedness of the test}}
\label{report-step-UnadjustedLangevin-LSI-GaussianWeakEquation-lean-7}


A bounded continuous test satisfies \(|\psi|\le M\), hence belongs to \(L^2(\rho_t)\) because \(\rho_t\) is a probability measure. The general flux-pairing result therefore gives integrability of \(\langle F_{t,x}(y),v\rangle\psi(y)\).



A bounded test has finite square energy under the interpolation probability. Pair it with the finite-energy flux by Cauchy–Schwarz to obtain integrability of \(F_i\psi\), the boundary product required for global integration by parts. No separate pointwise decay of the test at infinity is assumed.



The test itself is bounded, so multiplying by an integrable flux component preserves integrability. This supplies the boundary-product condition in the global coordinate integration-by-parts theorem.

\subsection{\texorpdfstring{Match the divergence to the independently differentiated Gaussian kernel}{Match the divergence to the independently differentiated Gaussian kernel}}
\label{report-step-UnadjustedLangevin-LSI-GaussianWeakEquation-lean-8}


Direct differentiation of the normalized Gaussian kernel gives its time derivative:

\[\dot K_t(x,y)=K_t(x,y)\left(-\frac d{2t}+\frac{\|y-x\|^2}{4t^2}-\frac{\|\nabla U(x)\|^2}{4}\right).\]

Summing the coordinate flux derivatives gives the same scalar expression as the kernel time derivative. The identity follows by expanding the centered displacement square and canceling its mixed terms. Thus the weak PDE uses a proved equality of explicit derivatives, rather than an assumed evolution equation.



The divergence and time derivative have been calculated independently from their explicit formulas. Algebraic comparison identifies them pointwise, fixing both the diffusion and drift signs in the weak equation.

\subsection{\texorpdfstring{Dominate each unit-\allowbreak{}coordinate flux derivative by a quadratic Gaussian energy}{Dominate each unit-coordinate flux derivative by a quadratic Gaussian energy}}
\label{report-step-UnadjustedLangevin-LSI-GaussianWeakEquation-lean-9}


For \(\|v\|=1\), Cauchy–Schwarz bounds \(|\langle r,v\rangle|\le\|r\|\) and \(|\langle g,v\rangle|\le\|g\|\). The inequality \(\|r\|\|g\|\le\|r\|^2+\|g\|^2\) then gives

\[\left|D_v\langle F_{t,x},v\rangle\right|\le\left(\frac1{4t^2}+\frac1{2t}\right)K_t(x,y)(1+\|g(x)\|^2+\|r\|^2).\]

A derivative of one flux coordinate is bounded by the kernel times a constant plus quadratic displacement and gradient energies. These are precisely the terms included in the established Gaussian energy envelope.



Earlier results used here: \cite[Section~3.37.4]{LangevinLean2026ULASC}.
\subsection{\texorpdfstring{Integrate the coordinate derivative envelope under the actual pair measure}{Integrate the coordinate derivative envelope under the actual pair measure}}
\label{report-step-UnadjustedLangevin-LSI-GaussianWeakEquation-lean-10}


The common envelope \(K_t(1+\|g(X)\|^2+\|Y-c_t(X)\|^2)\) is jointly integrable: its integral is \(1+G_0+2td\), where \(G_0=\int\|\nabla U\|^2d\mu\). The preceding pointwise bound therefore proves joint integrability of every unit-coordinate flux derivative.



Integrating the quadratic envelope first in the Gaussian observation and then in the initial law gives a finite value from dimension and \(G_0\). This verifies absolute integrability of each coordinate derivative.

\subsection{\texorpdfstring{Apply global integration by parts after proving all three product integrabilities}{Apply global integration by parts after proving all three product integrabilities}}
\label{report-step-UnadjustedLangevin-LSI-GaussianWeakEquation-lean-11}


For one unit coordinate \(v\), global integration by parts gives

\[\int D_v\langle F_{t,x}(y),v\rangle\,\psi(y)\,d\mu(x)\,dy=-\int\langle F_{t,x}(y),v\rangle D_v\psi(y)\,d\mu(x)\,dy.\]

The product, its directional derivative, and the flux-gradient pairing are all integrable. These checks make global integration by parts applicable without assuming compact support of the bounded test.

\subsection{\texorpdfstring{Derive the actual weak PDE for every bounded C\ensuremath{{}^{1}} finite-\allowbreak{}energy test}{Derive the actual weak PDE for every bounded C¹ finite-energy test}}
\label{report-step-UnadjustedLangevin-LSI-GaussianWeakEquation-lean-12}


Sum the coordinate identities over an orthonormal basis. The divergence identity turns the sum on the left into \(\dot K_t\), while completeness of the basis turns the sum on the right into \(\langle F_{t,x},\nabla\psi\rangle\). Fubini then yields

\[\int\dot p_t(y)\psi(y)\,dy=-\int\langle F_{t,x}(y),\nabla\psi(y)\rangle\,d\mu(x)\,dy.\]

Sum the coordinate identities and use divergence equal to the kernel time derivative. Fubini identifies the mixture derivative on the left, producing the actual weak PDE for every admitted bounded finite-energy test.

\section{\texorpdfstring{Relative Fisher energy and the sharp frozen-\allowbreak{}gradient error}{Relative Fisher energy and the sharp frozen-gradient error}}
\label{proof-module-UnadjustedLangevin-LSI-GaussianEnergyIntegrability-lean}


Let \(\rho_t\) be the actual frozen Gaussian interpolation, \(r_t=d\rho_t/d\nu\), \(s_t=\nabla\log r_t\), and \(G_0=\mathbb E_\mu\|\nabla U\|^2\). This module proves \(I_t\le4(1+L^2t^2)G_0+8tdL^2+d/t\), local spacetime integrability of the weighted relative-score energy, and the exact pair-law error bound \(\mathbb E\|\nabla U(Y_t)-\nabla U(X)\|^2\le L^2(t^2G_0+2td)\). These estimates provide domination for the unbounded entropy test and its truncation limits.



The actual relative score is the sum of the Lebesgue score and \(\nabla U\). The inequality \(\|a+b\|^2\le2\|a\|^2+2\|b\|^2\), Gaussian score energy \(d/(2t)\), and propagated gradient energy give a finite bound on every positive compact time interval. The frozen-gradient discrepancy has energy at most \(L^2(t^2G_0+2td)\). These two energy estimates dominate flux pairings in spacetime and justify the later removal of bounded-log weights.



Earlier results used here: \cite[Section~3.38]{LangevinLean2026ULASC}.
\subsection{\texorpdfstring{A square bound for the sum of two vectors}{A square bound for the sum of two vectors}}
\label{report-step-UnadjustedLangevin-LSI-GaussianEnergyIntegrability-lean-1}


The triangle inequality gives \(\|x+y\|\le\|x\|+\|y\|\). Square both nonnegative sides and use \(2\|x\|\|y\|\le\|x\|^2+\|y\|^2\), which follows from \((\|x\|-\|y\|)^2\ge0\). This proves \(\|x+y\|^2\le2\|x\|^2+2\|y\|^2\), the pointwise estimate used to split relative Fisher energy.



Expand \(\|a+b\|^2\), bound the mixed term by \(2\|a\|\|b\|\), and use \(2uv\le u^2+v^2\). This gives the factor two without requiring any relation between the vectors.

\subsection{\texorpdfstring{Splitting the relative-\allowbreak{}score energy}{Splitting the relative-score energy}}
\label{report-step-UnadjustedLangevin-LSI-GaussianEnergyIntegrability-lean-2}


The jointly measurable score field agrees at positive time with the actual relative score. Its decomposition is \(s_t=\nabla U+\nabla\log p_t\). The propagated gradient is square integrable, and Gaussian smoothing gives \(\mathbb E_{\rho_t}\|\nabla\log p_t\|^2\le d/(2t)\). These facts also prove integrability of the relative-score square.



Apply the squared sum bound pointwise and integrate: \(I_t\le2G_t+2\mathbb E_{\rho_t}\|\nabla\log p_t\|^2\le2G_t+d/t\). The proof keeps the integrability witnesses for all terms before using integral monotonicity.



Split the relative score into the Lebesgue score and potential gradient. The preceding square inequality bounds its energy by twice each component energy, which are estimated by Gaussian smoothing and propagated score moments.



Earlier results used here: \cite[Section~3.26.3]{LangevinLean2026ULASC}.
\subsection{\texorpdfstring{An explicit finite relative-\allowbreak{}score energy bound}{An explicit finite relative-score energy bound}}
\label{report-step-UnadjustedLangevin-LSI-GaussianEnergyIntegrability-lean-3}


The interpolation gradient-moment estimate is \(G_t\le2(1+L^2t^2)G_0+4tdL^2\). Substitute it into \(I_t\le2G_t+d/t\) and distribute the factor two. This gives \(I_t\le4(1+L^2t^2)G_0+8tdL^2+d/t\). Every quantity on the right is determined by the initial gradient energy and the original potential; no entropy-production assumption is used.



Insert the Gaussian Fisher bound \(d/(2t)\) and the propagated gradient estimate. The result is finite at each positive time and uniformly integrable on intervals bounded away from zero.



Earlier results used here: \cite[Section~3.26.3]{LangevinLean2026ULASC}.
\subsection{\texorpdfstring{The weighted density representation of score energy}{The weighted density representation of score energy}}
\label{report-step-UnadjustedLangevin-LSI-GaussianEnergyIntegrability-lean-4}


The explicitly constructed relative density satisfies \(r_t\nu=\rho_t\). The weighted-density integrability equivalence therefore transports integrability of \(\|s_t\|^2\) under \(\rho_t\) to integrability of \(r_t\|s_t\|^2\) under the fixed target \(\nu\). The corresponding integral identity is \(\int r_t\|s_t\|^2d\nu=\int\|s_t\|^2d\rho_t\). Strict positivity and measurability of the constructed density justify the ordinary weighted-integral formulas.



Changing from the actual proposed law to its density multiplies the score square by \(p_t\), or by \(r_t\) under Gibbs. These are equivalent representations of the same finite energy.



Earlier results used here: \cite[Section~3.26.3]{LangevinLean2026ULASC}.
\subsection{\texorpdfstring{Spacetime integrability of relative-\allowbreak{}score energy}{Spacetime integrability of relative-score energy}}
\label{report-step-UnadjustedLangevin-LSI-GaussianEnergyIntegrability-lean-5}


The relative density and the chosen relative-score field are jointly measurable in time and space. Their weighted energy is nonnegative. On a positive interval \([a,b]\), the section at every time is integrable by the density-change theorem.



Since \(0<a\le t\le b\), the quantitative estimate is bounded uniformly by \(4(1+L^2b^2)G_0+8bdL^2+d/a\). The integral of the section norm is exactly its nonnegative energy, so this finite scalar bound makes it integrable in time. The product criterion yields \(r_t\|s_t\|^2\in L^1(dt\,d\nu)\) on the interval. This is the spacetime dominating function required to pass from bounded entropy tests to the full relative Fisher term.



The explicit time bound is a sum of powers of time and \(1/t\). On a positive compact interval each is integrable, proving spacetime square integrability of the measurable relative-score field.



Earlier results used here: \cite[Section~3.26.3]{LangevinLean2026ULASC}.
\subsection{\texorpdfstring{The actual frozen-\allowbreak{}gradient discrepancy}{The actual frozen-gradient discrepancy}}
\label{report-step-UnadjustedLangevin-LSI-GaussianEnergyIntegrability-lean-6}


For the fixed independent pair \((X,Z)\sim\mu\otimes\mathcal N(0,I)\), define \(e_t(X,Z)=\nabla U(X-t\nabla U(X)+\sqrt{2t}Z)-\nabla U(X)\). The underlying probability space is independent of time. This avoids assuming a jointly measurable choice of a state-space conditional expectation when integrating the error energy over time.



The discrepancy is evaluated at the actual independent-noise proposal and its initial state. Its measurability follows from the continuous gradient and update map, fixing one field for subsequent pair integrals.



Earlier results used here: \cite[Section~3.34.1]{LangevinLean2026ULASC}.
\subsection{\texorpdfstring{Spacetime integrability of discrepancy energy}{Spacetime integrability of discrepancy energy}}
\label{report-step-UnadjustedLangevin-LSI-GaussianEnergyIntegrability-lean-7}


The proposal and Lipschitz gradient make \((t,X,Z)\mapsto\|e_t(X,Z)\|^2\) jointly measurable. For \(t\ge0\), the exact Gaussian calculation already gives integrability and \(\mathbb E\|e_t\|^2\le L^2(t^2G_0+2td)\). On a compact nonnegative time interval this is at most \(L^2(b^2G_0+2bd)\). The time measure is finite, so Fubini's integrability criterion proves the squared error is integrable on the time/\allowbreak{}initial-state/\allowbreak{}noise product.



The pointwise gradient Lipschitz estimate bounds its square by \(L^2\) times displacement square. The exact displacement formula is polynomial in time with finite coefficients, giving an integrable spacetime energy.



Earlier results used here: \cite[Section~3.34.1]{LangevinLean2026ULASC}.
\subsection{\texorpdfstring{The gradient-\allowbreak{}error energy under the Gaussian pair law}{The gradient-error energy under the Gaussian pair law}}
\label{report-step-UnadjustedLangevin-LSI-GaussianEnergyIntegrability-lean-8}


Map the independent initial/\allowbreak{}noise pair by \(T_t(x,z)=(x,x-t\nabla U(x)+\sqrt{2t}z)\). The squared error as a function of the output pair is \(F(x,y)=\|\nabla U(y)-\nabla U(x)\|^2\). Its composition with \(T_t\) is exactly the frozen-gradient discrepancy from the Gaussian energy calculation. Measurability of \(T_t\) and continuity of \(F\) let the map-measure integrability and integral theorems transport both the integrability witness and the bound \(L^2(t^2G_0+2td)\) to the actual initial/\allowbreak{}output pair law.



The pair-law change of variables transports the discrepancy energy from independent noise to the Gaussian starting-observation pair. Because the integrand is nonnegative, the identity can be established before signed Fubini is needed.



Earlier results used here: \cite[Section~3.34.1]{LangevinLean2026ULASC}; \cite[Section~3.36.5]{LangevinLean2026ULASC}.
\subsection{\texorpdfstring{The kernel representation of gradient-\allowbreak{}error energy}{The kernel representation of gradient-error energy}}
\label{report-step-UnadjustedLangevin-LSI-GaussianEnergyIntegrability-lean-9}


The actual pair law is identified with \(K_t(x,y)\,\mu(dx)dy\) by the previously proved Gaussian proposal pair-density theorem. The proof specializes its measurable center to \(x-t\nabla U(x)\); the precision and step-size versions of the kernel agree pointwise.



Rewrite the pair-law energy theorem using this exact density identity. The weighted-density integrability equivalence gives \(K_t\|\nabla U(y)-\nabla U(x)\|^2\in L^1(\mu\otimes dy)\), and its integral formula gives the sharp bound \(\int K_t\|\nabla U(y)-\nabla U(x)\|^2d\mu\,dy\le L^2(t^2G_0+2td)\). This transport preserves the exact constants needed in the final entropy recurrence; a cruder pointwise Gaussian envelope is not substituted for the error energy.



Rewriting the pair measure using its conditional Gaussian density gives the kernel-weighted energy. Thus the same sharp bound is available in the weak PDE's Lebesgue integral representation.



Earlier results used here: \cite[Section~3.34.1]{LangevinLean2026ULASC}; \cite[Section~3.37.12]{LangevinLean2026ULASC}; \cite[Section~3.37.13]{LangevinLean2026ULASC}; \cite[Section~3.37.4]{LangevinLean2026ULASC}; \cite[Section~3.38.1]{LangevinLean2026ULASC}.
\section{\texorpdfstring{The complete entropy chain rule for the actual Gaussian interpolation}{The complete entropy chain rule for the actual Gaussian interpolation}}
\label{proof-module-UnadjustedLangevin-LSI-GaussianEntropyIdentity-lean}


This module proves the entropy identity and positive-time absolute continuity for the actual frozen-drift Gaussian interpolation. The proof starts with bounded smooth log tests, applies the global Gaussian weak equation and a spacetime \(L^1\) time chain rule, then removes both spatial and time truncations by dominated convergence. All regularity and integrability needed in these steps are constructed from the original Gaussian update and the initial gradient energy.



Write \(\pi=Z^{-1}e^{-U}dx\), \(\rho_t=\operatorname{Law}(X-t\nabla U(X)+\sqrt{2t}Z_0)\), \(r_t=d\rho_t/d\pi\), and \(S_t=\nabla\log r_t\). Here \(Z_0\) is an independent standard Gaussian. The eventual conclusion is \(H(\rho_b\mid\pi)-H(\rho_a\mid\pi)=\int_a^b q(t)dt\), with the explicit raw Gaussian flux \(q\), for every \(0<a\le b\). This identity proves local absolute continuity and the a.e. derivative; they are not assumed.



Use bounded tests \(\chi_m(\log r_t)\) and their entropy primitives to obtain an exact integrated identity on \([a,b]\subset(0,\infty)\). Finite Fisher and Gaussian velocity energies dominate the flux uniformly in \(m\), while the genuine entropy integrand dominates the endpoints. Dominated convergence then removes the bounded-log modification in both places. The limiting identity gives local absolute continuity and the derivative \(H'=-I+\mathbb E\langle\nabla U(Y_t)-\nabla U(X),\nabla\log r_t(Y_t)\rangle\) almost everywhere.



Earlier results used here: \cite[Section~3.46]{LangevinLean2026ULASC}.
\subsection{\texorpdfstring{Removing bounded log weights by dominated convergence}{Removing bounded log weights by dominated convergence}}
\label{report-step-UnadjustedLangevin-LSI-GaussianEntropyIdentity-lean-1}


Let \(w_m(z)=[1+(z/(m+1))^2]^{-1}\). For a measurable scalar field \(z\) and an integrable field \(f\),

\[\int w_m(z(x))f(x)\,d\nu(x)\longrightarrow\int f(x)\,d\nu(x).\]

The pointwise limit is \(w_m(z(x))\to1\), proved in the pointwise convergence of bounded-log derivative weights to one. Each weighted integrand is measurable, because the weight is obtained by division, squaring and inversion of measurable real functions. Furthermore \(0\le w_m\le1\), so \(|w_m(z)f|\le|f|\). The proof invokes dominated convergence with the single integrable majorant norm. This lemma later removes the truncation from the actual Gaussian flux, whose two factors are square integrable.



The derivative weights converge pointwise to one and stay between zero and one. Multiplying by the already integrable unweighted pairing gives a common absolute majorant, so dominated convergence removes the weight without changing the reference measure.

\subsection{\texorpdfstring{Why the square-\allowbreak{}integrable flux pairing is integrable}{Why the square-integrable flux pairing is integrable}}
\label{report-step-UnadjustedLangevin-LSI-GaussianEntropyIdentity-lean-2}


If \(v,s\in L^2(\nu;E)\), then

\[|\langle v,s\rangle|\le\|v\|\|s\|,\qquad \|v\|\|s\|\in L^1(\nu).\]

Cauchy–Schwarz bounds the integral of the product of the two norms by their finite square norms. The absolute inner product is bounded by this norm product, proving absolute integrability of the pairing. This remains true on the stated measure space without assuming its total mass is finite.



For square-integrable fields \(a,b\), use \(|\langle a,b\rangle|\le(\|a\|^2+\|b\|^2)/2\). The right side has finite integral, proving absolute flux integrability before signed Fubini or limiting arguments.

\subsection{\texorpdfstring{The real Gibbs density and its normalization}{The real Gibbs density and its normalization}}
\label{report-step-UnadjustedLangevin-LSI-GaussianEntropyIdentity-lean-3}


Write \(Z=\int e^{-U(x)}\,dx\), and let \(\pi\) be the normalized Gibbs law. Define

\[J(y)=Z^{-1}e^{-U(y)}.\]

The existing probability-mass argument proves \(0<Z<\infty\); differentiability makes \(U\) measurable and gives strict positivity of the partition function. Consequently \(J(y)>0\) at every point. The measure identity is

\[\pi(dy)=J(y)\,dy.\]

The density identity follows from multiplying the unnormalized Gibbs weight by the finite scalar \(Z^{-1}\). Since \(0<Z<\infty\), the resulting measurable function \(J\) is positive and has integral one. The same expression therefore represents the target in ordinary Lebesgue integration; strong convexity is not needed for this change-of-density step.



The Gibbs density is positive and finite pointwise, and its integral is one. Thus \(J\Phi_m(r_t)\) may be read as a Lebesgue representation of the same bounded relative-entropy primitive integrated against Gibbs.



Earlier results used here: \cite[Section~3.26.1]{LangevinLean2026ULASC}; \cite[Section~3.26.2]{LangevinLean2026ULASC}; \cite[Section~3.26.3]{LangevinLean2026ULASC}; \cite[Section~3.26.9]{LangevinLean2026ULASC}.
\subsection{\texorpdfstring{The time derivative of the weighted entropy primitive}{The time derivative of the weighted entropy primitive}}
\label{report-step-UnadjustedLangevin-LSI-GaussianEntropyIdentity-lean-4}


Let \(p_t\) be the actual Lebesgue density of the frozen Gaussian update, and let \(r_t=p_t/J\) be its relative density. For the bounded entropy primitive \(\Phi_m\), one has \(\Phi_m'(r)=\chi_m(\log r)\) for \(r>0\). At every positive time, the actual kernel differentiation theorem gives \(\partial_t p_t=\dot p_t\), hence

\[\partial_t r_t=Ze^{U}\dot p_t,\qquad \partial_t[J\Phi_m(r_t)]=\chi_m(\log r_t)\dot p_t.\]

The proof applies the time chain rule first to the constant multiple defining \(r_t\), then to the derivative formula for the bounded entropy primitive, then to multiplication by \(J\). The remaining coefficient cancels because \(J(y)Ze^{U(y)}=1\), using the inverse of the strictly positive real partition function and \(e^{-U}e^U=1\). This cancellation is what permits the spacetime estimate to use the Lebesgue \(L^1\) bound on \(\dot p_t\).



Because \(J\) is time independent and \(Jr_t=p_t\), differentiating \(J\Phi_m(r_t)\) gives \(\chi_m(\log r_t)\dot p_t\). This cancellation avoids differentiating the Gibbs density in time.



Earlier results used here: \cite[Section~3.26.2]{LangevinLean2026ULASC}; \cite[Section~3.26.3]{LangevinLean2026ULASC}; \cite[Section~3.26.9]{LangevinLean2026ULASC}.
\subsection{\texorpdfstring{Identifying finite entropy through the chosen density}{Identifying finite entropy through the chosen density}}
\label{report-step-UnadjustedLangevin-LSI-GaussianEntropyIdentity-lean-5}


Suppose \(\mu=r\nu\) is a probability measure, \(r\ge0\) is measurable, and \(H(\mu\mid\nu)<\infty\). Then

\[\mathrm{kl}(r)=r\log r+1-r\in L^1(\nu),\qquad \int\mathrm{kl}(r)\,d\nu=H(\mu\mid\nu).\]

The Radon–Nikodym derivative of \(r\nu\) relative to \(\nu\) is \(r\) almost everywhere. Finite KL gives integrability of its logarithm under \(r\nu\); changing measure makes \(r\log r\) integrable under \(\nu\). Since \(r\) and the constant one both have finite integral, so does \(r\log r-r+1\). Their probability masses agree, so integrating the last two terms cancels them and gives the stated entropy identity.



The actual KL is finite at positive time by LSI and the Fisher bound. The density formula therefore identifies it with the ordinary integral of \(r_t\log r_t-r_t+1\), whose mass corrections cancel by normalization.

\subsection{\texorpdfstring{Integrability of every actual bounded entropy}{Integrability of every actual bounded entropy}}
\label{report-step-UnadjustedLangevin-LSI-GaussianEntropyIdentity-lean-6}


At \(t>0\), the actual interpolation has a positive \(C^1\) relative density \(r_t\). Its Fisher information is finite from the initial score moment and Gaussian smoothing; LSI then gives finite actual KL. The bounded primitive satisfies

\[0\le\Phi_m(r_t(y))\le\mathrm{kl}(r_t(y)).\]

The preceding density lemma proves that the right-hand side is integrable. The left-hand side is measurable: continuity of \(r_t\) and the derivative of \(\Phi_m\) at positive arguments give continuity of their composition. The proof integrable domination argument therefore derives integrability for each truncation from a verified KL majorant. Initial density differentiability, initial Fisher information and a second moment of the state are not required.



The bounded primitive lies between zero and the genuine nonnegative entropy integrand. Finite actual KL therefore proves integrability of each approximating entropy before its time identity is formed.



Earlier results used here: \cite[Section~3.26.3]{LangevinLean2026ULASC}.
\subsection{\texorpdfstring{Convergence of bounded entropy to actual KL}{Convergence of bounded entropy to actual KL}}
\label{report-step-UnadjustedLangevin-LSI-GaussianEntropyIdentity-lean-7}


For every fixed positive time,

\[\int\Phi_m(r_t)\,d\pi\longrightarrow H(\rho_t\mid\pi).\]

Pointwise convergence \(\Phi_m(r)\to\mathrm{kl}(r)\) holds for \(r>0\). The domination \(0\le\Phi_m(r_t)\le\mathrm{kl}(r_t)\), together with integrability proved above, gives dominated convergence under the fixed Gibbs measure. Before taking the limit, Identify the integral of \(r_t\log r_t-r_t+1\) with the actual KL through tilted measure the Gibbs-relative density \(r_t=Ze^Up_t\). Thus the limiting quantity is the KL of the actual Gaussian law, rather than an independently specified entropy path.



Pointwise convergence of the primitive to \(r\log r-r+1\), together with the finite entropy majorant, gives convergence of both positive-time endpoint entropies. This limit is independent of any time-flux limit.



Earlier results used here: \cite[Section~3.26.3]{LangevinLean2026ULASC}.
\subsection{\texorpdfstring{The bounded entropy and its time flux}{The bounded entropy and its time flux}}
\label{report-step-UnadjustedLangevin-LSI-GaussianEntropyIdentity-lean-8}


The two functions used in the time integration are

\[H_m(t)=\int\Phi_m(r_t(y))\,d\pi(y),\qquad q_m(t)=\int\chi_m(\log r_t(y))\dot p_t(y)\,dy.\]

the bounded entropy \(H_m(t)\) uses the fixed Gibbs reference measure; the bounded entropy flux \(q_m(t)\) uses Lebesgue measure. Their different measures match the exact cancellation in the weighted chain rule. The definitions are explicit functions of the actual mixture density and its actual kernel time derivative.



The bounded entropy is paired with its explicit time flux on a compact positive interval. Boundedness of the log test lets the spacetime \(L^1\) density-derivative estimate dominate that flux.



Earlier results used here: \cite[Section~3.26.3]{LangevinLean2026ULASC}.
\subsection{\texorpdfstring{Admissibility of the actual bounded log test}{Admissibility of the actual bounded log test}}
\label{report-step-UnadjustedLangevin-LSI-GaussianEntropyIdentity-lean-9}


The spatial weak test is \(\psi_m(y)=\chi_m(\log r_t(y))\). Positivity and \(C^1\) regularity of \(r_t\), combined with smoothness of \(\chi_m\), show that \(\psi_m\in C^1\). Its uniform bound is

\[|\psi_m(y)|\le(m+1)\pi/2.\]

The spatial chain rule gives

\[\nabla\psi_m(y)=w_m(\log r_t(y))\nabla\log r_t(y),\qquad \|\nabla\psi_m(y)\|\le\|\nabla\log r_t(y)\|.\]

The test-gradient norm is at most the relative-score norm, whose square is integrable. Continuity of the test derivative supplies measurability. Thus the bounded logarithmic test is \(C^1\) and has finite gradient energy, exactly the conditions required by the actual Gaussian weak equation.



The test is bounded and \(C^1\), and its gradient norm is at most the relative-score norm. Finite Fisher information therefore places its gradient in the exact finite-energy class admitted by the actual weak equation.



Earlier results used here: \cite[Section~3.26.3]{LangevinLean2026ULASC}.
\subsection{\texorpdfstring{The bounded entropy flux is the actual spatial flux}{The bounded entropy flux is the actual spatial flux}}
\label{report-step-UnadjustedLangevin-LSI-GaussianEntropyIdentity-lean-10}


Let \(K_t(x,y)\) be the Gaussian transition density, \(g=\nabla U\), and \(S_t=\nabla\log r_t\). The proved Gaussian weak equation, applied to the admissible test above, gives

\[q_m(t)=-\iint K_t(x,y)w_m(\log r_t(y))\left\langle g(x)-\frac{y-x+tg(x)}{2t},S_t(y)\right\rangle\,\mu(dx)\,dy.\]

The underlying flux is \(K_t[g(x)-(y-x+tg(x))/(2t)]\). Its divergence equals the actual kernel time derivative; global integration by parts is justified by the previously proved finite-gradient-energy estimates. Substitute the test-gradient chain rule and the equality between the explicit measurable score field and the actual relative score. The remaining rearrangement is bilinearity of the inner product and scalar multiplication.



The weak equation converts the time-derivative pairing into the negative spatial flux pairing. Its sign is fixed by integration by parts, and its integrability has been established for this very bounded test.



Earlier results used here: \cite[Section~3.26.3]{LangevinLean2026ULASC}.
\subsection{\texorpdfstring{The bounded time chain rule without a spatial time-\allowbreak{}majorant}{The bounded time chain rule without a spatial time-majorant}}
\label{report-step-UnadjustedLangevin-LSI-GaussianEntropyIdentity-lean-11}


On any \(0<a\le b\), the weighted primitive has derivative \(\chi_m(\log r_t)\dot p_t\). The kernel theorem proves \(\dot p\in L^1([a,b]\times E)\), while \(\chi_m\) is bounded by \((m+1)\pi/2\). Thus the derivative of the weighted primitive is jointly integrable.



Integrability at time \(a\) follows from the bounded-entropy lemma after rewriting Gibbs measure as \(J(y)dy\). The generic spacetime FTC/\allowbreak{}Fubini theorem then yields

\[H_m(b)-H_m(a)=\int_a^b q_m(t)\,dt.\]

Integrate the pointwise chain rule in time first, using the bounded test to dominate the time derivative. Fubini then yields the entropy endpoint identity without requiring an integrable spatial envelope uniform near each individual observation.



Earlier results used here: \cite[Section~3.26.3]{LangevinLean2026ULASC}; \cite[Section~3.46.3]{LangevinLean2026ULASC}.
\subsection{\texorpdfstring{Expressing the flux on one independent probability space}{Expressing the flux on one independent probability space}}
\label{report-step-UnadjustedLangevin-LSI-GaussianEntropyIdentity-lean-12}


Take \(X\sim\mu\) and \(Z\sim N(0,I)\) independently, and set \(Y_t=X-tg(X)+\sqrt{2t}Z\). Define

\[V_t(X,Z)=g(X)-\frac{\sqrt{2t}}{2t}Z,\qquad \widehat S_t(X,Z)=S_t(Y_t),\qquad q(t)=-\mathbb E\langle V_t,\widehat S_t\rangle.\]

These fields are defined on the same independent initial-state and Gaussian-noise probability for every time. The frozen velocity is the component spatial flux divided by the positive transition density and evaluated at the actual proposal. Using this fixed reference space makes joint measurability and time domination available for the limiting entropy identity.



The Gaussian pair-law identity replaces kernel-weighted starting-observation integrals by expectations over one fixed independent state-noise probability. This puts all energy and pairing estimates on the same reference space.



Earlier results used here: \cite[Section~3.34.1]{LangevinLean2026ULASC}.
\subsection{\texorpdfstring{Square integrability of the two actual Gaussian factors}{Square integrability of the two actual Gaussian factors}}
\label{report-step-UnadjustedLangevin-LSI-GaussianEntropyIdentity-lean-13}


The initial gradient belongs to \(L^2(\mu)\), and the standard Gaussian noise belongs to \(L^2\). Pulling both to the independent product space and taking their scalar combination proves \(V_t\in L^2(\mu\otimes\gamma)\). The actual relative score belongs to \(L^2(\rho_t)\). Since \(\rho_t\) is the pushforward of \(\mu\otimes\gamma\) by \(Y_t\), the measurable-map characterization of \(L^2\) gives

\[\widehat S_t=S_t\circ Y_t\in L^2(\mu\otimes\gamma).\]

the finite velocity energy from the initial-gradient and Gaussian moments performs the product-space linear combination. the finite score energy transported through the proposal map first identifies the explicit score with the actual gradient of \(\log r_t\), then applies the pushforward integral identity for square energies. Consequently their inner product is integrable.



The relative score has finite square energy by the Fisher estimate. The frozen velocity has finite square energy by its Gaussian component and initial gradient, so their flux pairing is absolutely integrable by Cauchy–Schwarz.



Earlier results used here: \cite[Section~3.26.3]{LangevinLean2026ULASC}; \cite[Section~3.36.5]{LangevinLean2026ULASC}.
\subsection{\texorpdfstring{The exact frozen velocity energy}{The exact frozen velocity energy}}
\label{report-step-UnadjustedLangevin-LSI-GaussianEntropyIdentity-lean-14}


Independence and the zero Gaussian mean cancel the drift-noise cross term:

\[\mathbb E\|V_t\|^2=\mathbb E\|g(X)\|^2+\frac{d}{2t},\qquad t>0.\]

The proof specializes the already verified affine-Gaussian moment formula with coefficients \(1\) and \(-\sqrt{2t}/(2t)\). The squared coefficient is \((2t)/(4t^2)=1/(2t)\), using \((\sqrt{2t})^2=2t\). This exact finite energy gives a compact-time bound that is uniform away from zero. No moment of \(X\) itself is used.



The frozen velocity consists of the Gaussian component score plus the initial drift. Their mixed expectation vanishes by centered independent Gaussian noise, giving the exact sum of component-score and initial-gradient energies.

\subsection{\texorpdfstring{Transporting the bounded flux to independent noise}{Transporting the bounded flux to independent noise}}
\label{report-step-UnadjustedLangevin-LSI-GaussianEntropyIdentity-lean-15}


The law of \((X,Y_t)\) has density \(K_t(x,y)\) with respect to \(\mu(dx)dy\). Therefore the joint flux identity becomes

\[q_m(t)=-\mathbb E\!\left[w_m(\log r_t(Y_t))\langle V_t,\widehat S_t\rangle\right].\]

Consider the measurable map t sending \((x,z)\) to \((x,Y_t)\), the kernel k, and the weighted pairing f. the starting-proposal pair-density identity identifies the pushforward law. The integral under its real density becomes \(\iint K F\), and the change-of-variables formula for an integrable pushforward changes that integral back to the independent product space. Finally \(Y_t-(X-tg(X))=\sqrt{2t}Z\) reduces the residual to the specified velocity field. The change of variables concerns the actual Gaussian law and does not assume an invariant Euler step.



Changing variables in the bounded flux retains the same state-space test evaluated at the actual proposal. This lets the derivative weight remain bounded by one while the independent-noise energy majorant is applied.



Earlier results used here: \cite[Section~3.26.3]{LangevinLean2026ULASC}; \cite[Section~3.34.1]{LangevinLean2026ULASC}; \cite[Section~3.36.5]{LangevinLean2026ULASC}; \cite[Section~3.37.4]{LangevinLean2026ULASC}.
\subsection{\texorpdfstring{Removing the truncation at fixed positive time}{Removing the truncation at fixed positive time}}
\label{report-step-UnadjustedLangevin-LSI-GaussianEntropyIdentity-lean-16}


The raw pairing \(\langle V_t,\widehat S_t\rangle\) is integrable by the two \(L^2\) estimates. The scalar map \((x,z)\mapsto\log r_t(Y_t)\) is measurable. Since \(0\le w_m\le1\) and \(w_m\to1\), dominated convergence on the independent probability space yields

\[q_m(t)\longrightarrow q(t),\qquad t>0.\]

The dominating function is \(|\langle V_t,\widehat S_t\rangle|\). Crucially this limit is performed after spatial integration by parts; it does not ask for integrability of \(|\log r_t|\,|\dot p_t|\), which is not supplied by the original assumptions.



At a fixed positive time, all factors except the bounded-log weight remain unchanged. Since that weight tends to one and its absolute value is at most one, the integrable full flux dominates the entire sequence.



Earlier results used here: \cite[Section~3.26.3]{LangevinLean2026ULASC}; \cite[Section~3.34.1]{LangevinLean2026ULASC}; \cite[Section~3.36.5]{LangevinLean2026ULASC}.
\subsection{\texorpdfstring{A fixed-\allowbreak{}reference formula for the relative Fisher energy}{A fixed-reference formula for the relative Fisher energy}}
\label{report-step-UnadjustedLangevin-LSI-GaussianEntropyIdentity-lean-17}


Define

\[I(t)=\int r_t(y)\|S_t(y)\|^2\,d\pi(y).\]

The actual relative-density identity gives \(I(t)=\int\|S_t\|^2\,d\rho_t\). Applying the actual Gaussian pushforward once more yields

\[\mathbb E\|\widehat S_t\|^2=I(t).\]

The first definition fixes the reference measure for time measurability and spacetime integrability. The equality uses the weighted representation of actual relative Fisher information and the measurable-map Bochner integral formula. The proof supplies strong measurability of the score norm square explicitly, fixing time in the joint measurable score field.



Pushing the relative-score square through the actual proposal map expresses \(I_t\) as an expectation on the independent-noise source. This representation is jointly measurable and directly comparable with the discrepancy energy.



Earlier results used here: \cite[Section~3.26.3]{LangevinLean2026ULASC}; \cite[Section~3.36.5]{LangevinLean2026ULASC}.
\subsection{\texorpdfstring{A uniform bound for every truncated flux}{A uniform bound for every truncated flux}}
\label{report-step-UnadjustedLangevin-LSI-GaussianEntropyIdentity-lean-18}


For any vectors \(v,s\) and any \(0\le w\le1\),

\[|w\langle v,s\rangle|\le\|v\|\|s\|\le\tfrac12(\|v\|^2+\|s\|^2).\]

The second inequality is \((\|v\|-\|s\|)^2\ge0\). Integrate on the independent Gaussian space, use the norm inequality for Bochner integrals, and substitute the two exact energies:

\[|q_m(t)|\le\frac12\left(G_0+\frac d{2t}+I(t)\right),\qquad G_0=\int\|g\|^2\,d\mu.\]

The absolute truncated pairing is bounded by half the sum of the two unweighted square energies. This bound contains no truncation index and is therefore suitable for both pointwise-time and time-integrated limits.



Earlier results used here: \cite[Section~3.26.3]{LangevinLean2026ULASC}; \cite[Section~3.34.1]{LangevinLean2026ULASC}; \cite[Section~3.36.5]{LangevinLean2026ULASC}.
\subsection{\texorpdfstring{Joint measurability of both time fluxes}{Joint measurability of both time fluxes}}
\label{report-step-UnadjustedLangevin-LSI-GaussianEntropyIdentity-lean-19}


The bounded flux is a parameter integral of \(\chi_m(\log r_t(y))\dot p_t(y)\). Both factors are jointly measurable in \((t,y)\), so parameter integration makes \(t\mapsto q_m(t)\) measurable. The raw flux uses the fixed independent space: the map \((t,x,z)\mapsto(t,Y_t)\) is measurable, the explicit relative-score field is jointly strongly measurable, and the velocity field is a measurable algebraic combination of \(g(X)\), \(t^{-1}\), \(\sqrt{2t}\), and \(Z\). Their inner product and its parameter integral are measurable. The definitions at nonpositive times serve only to make these globally specified measurable functions; all analytic conclusions are restricted to positive times.



The chosen score and velocity fields are jointly measurable, and the bounded-log weights are continuous functions of measurable densities. Their products and parameter integrals therefore define measurable scalar fluxes in time.



Earlier results used here: \cite[Section~3.34.1]{LangevinLean2026ULASC}.
\subsection{\texorpdfstring{The integrable majorant for the time limit}{The integrable majorant for the time limit}}
\label{report-step-UnadjustedLangevin-LSI-GaussianEntropyIdentity-lean-20}


Fix \(0<a\le b\). The actual Gaussian energy theorem proves \(I\in L^1([a,b])\). Set

\[D(t)=\tfrac12\left(G_0+\frac d{2a}+I(t)\right).\]

The interval has finite Lebesgue measure, so the constant term is integrable; hence \(D\in L^1([a,b])\). Since \(t\ge a>0\), \(d/(2t)\le d/(2a)\). The uniform flux estimate gives \(|q_m(t)|\le D(t)\) for every \(m\). The fixed-time limit \(q_m(t)\to q(t)\), followed by continuity of the absolute value, also gives \(|q(t)|\le D(t)\). The mathematical statement supplies the single majorant together with all three verified integrability/\allowbreak{}domination conclusions, ready for dominated convergence in time.



On a compact positive interval, the score and velocity square energies have a common integrable time bound. Their half-sum dominates every truncated flux and the limiting flux, establishing the spacetime dominated-convergence hypothesis.



Earlier results used here: \cite[Section~3.26.3]{LangevinLean2026ULASC}.
\subsection{\texorpdfstring{The full actual entropy identity}{The full actual entropy identity}}
\label{report-step-UnadjustedLangevin-LSI-GaussianEntropyIdentity-lean-21}


Let \(H(t)=H(\rho_t\mid\pi)\), represented by the actual finite KL. For each truncation,

\[H_m(b)-H_m(a)=\int_a^b q_m(t)\,dt.\]

At both endpoints, the spatial dominated-convergence theorem gives \(H_m\to H\). In time, the common integrable majorant \(D\) and the pointwise convergence \(q_m\to q\) give \(\int_a^b q_m\to\int_a^b q\). Uniqueness of limits therefore yields

\[H(b)-H(a)=\int_a^b q(t)\,dt.\]

the exact integrated entropy identity applies the general dominated convergence at both endpoints and in the integrated flux under Lebesgue measure restricted to \((a,b]\). Its endpoint convergence hypotheses are the already established actual-KL limits, and all time domination comes from the previous theorem. No entropy-path regularity is supplied as a hypothesis.



Take the truncation limit in the two endpoint values and in the integrated flux separately. Their established dominations then yield an exact identity for the genuine actual entropy between any two positive times.



Earlier results used here: \cite[Section~3.26.3]{LangevinLean2026ULASC}.
\subsection{\texorpdfstring{Absolute continuity and the a.\allowbreak{}e.\allowbreak{} derivative of actual entropy}{Absolute continuity and the a.e. derivative of actual entropy}}
\label{report-step-UnadjustedLangevin-LSI-GaussianEntropyIdentity-lean-22}


The raw flux is measurable and satisfies \(|q|\le D\), so it is integrable on every positive-time compact interval. Apply the full entropy identity with terminal time \(t\in[a,b]\):

\[H(t)=H(a)+\int_a^t q(u)\,du.\]

An indefinite integral of an \(L^1\) function is absolutely continuous, and the Lebesgue differentiation theorem identifies its derivative almost everywhere. Thus

\[H\in AC([a,b]),\qquad H'(t)=q(t)\quad\text{for almost every }t\in(a,b).\]

The proof derives interval integrability from the actual raw-flux measurability and the verified majorant, then invokes absolute continuity and the almost-everywhere derivative of an integrable time primitive. This completes the analytic time-regularity bridge from the original potential, LSI and finite initial score-energy assumptions. A later score-pairing identity rewrites the raw Gaussian flux as negative Fisher energy plus the frozen gradient discrepancy, yielding the quantitative entropy-production estimate.



The flux is integrable on each positive compact interval, so the exact integrated identity makes entropy locally absolutely continuous there. Its derivative equals that flux almost everywhere, supplying the precise temporal regularity used in the final scalar proof.



Earlier results used here: \cite[Section~3.26.3]{LangevinLean2026ULASC}.
\section{\texorpdfstring{Averaging the Gaussian component score}{Averaging the Gaussian component score}}
\label{proof-module-UnadjustedLangevin-LSI-GaussianScorePairing-lean}


The Gaussian component score \(S_{t,x}(y)=-(y-x+t\nabla U(x))/(2t)\) averages to the actual mixture score. This module establishes the vector identity \(p_t\nabla\log p_t=\int K_tS_{t,x}d\mu\), proves all joint integrability needed for pairing it with a continuous \(L^2(\rho_t)\) vector field, and transfers the resulting identity to the fixed independent initial/\allowbreak{}noise space used by entropy production.



The score of one Gaussian component is \(s_t(x,y)=-(y-x+t\nabla U(x))/(2t)\). Its conditional mean given \(Y_t=y\) is the mixture score \(\nabla\log p_t(y)\). Multiplying by any admissible state-space vector test and integrating proves the corresponding score-pairing identity without choosing a conditional mean explicitly. Square-integrability of both factors supplies absolute integrability, which is needed before exchanging the starting-point and observation integrals.

\subsection{\texorpdfstring{The score of one Gaussian component}{The score of one Gaussian component}}
\label{report-step-UnadjustedLangevin-LSI-GaussianScorePairing-lean-1}


For a fixed initial state \(x\), the Gaussian proposal has center \(c_t(x)=x-t\nabla U(x)\). Its score in the proposed state is \(S_{t,x}(y)=\nabla_y\log K_t(x,y)=-(y-c_t(x))/(2t)\). This definition is an explicit affine vector field. It is different from the mixture score, which averages this field over the posterior initial-state distribution.



Differentiate the logarithm of one Gaussian component in the observation variable. The normalizer is constant in that variable, leaving the negative centered displacement divided by \(2t\).

\subsection{\texorpdfstring{Square-\allowbreak{}integrability of the component score}{Square-integrability of the component score}}
\label{report-step-UnadjustedLangevin-LSI-GaussianScorePairing-lean-2}


Represent the actual initial/\allowbreak{}output pair as \((X,Y_t)\), with \(Y_t=c_t(X)+\sqrt{2t}Z\) and independent standard Gaussian \(Z\). Substituting into the component score cancels the center: \(S_{t,X}(Y_t)=-(\sqrt{2t}/(2t))Z\). The Gaussian identity random variable is square integrable; scalar multiplication and pullback to the product preserve this property. The map-measure theorem transfers it to the actual pair law. No initial position or gradient moment is needed for this component score.



Under the independent-noise representation the centered component displacement is \(\sqrt{2t}Z\). Its score energy is therefore \(d/(2t)\), finite for every positive time.



Earlier results used here: \cite[Section~3.34.1]{LangevinLean2026ULASC}; \cite[Section~3.36.5]{LangevinLean2026ULASC}.
\subsection{\texorpdfstring{Pulling a finite-\allowbreak{}energy state field to the pair space}{Pulling a finite-energy state field to the pair space}}
\label{report-step-UnadjustedLangevin-LSI-GaussianScorePairing-lean-3}


If \(V\in L^2(\rho_t)\), then \(V(Y_t)\in L^2(\mu\otimes\mathcal N(0,I))\), because \(\rho_t\) is the actual pushforward by the proposal. Map the independent pair to \((X,Y_t)\); a second application of the map-measure equivalence gives \((x,y)\mapsto V(y)\in L^2\) under the actual pair law. Continuity of \(V\) supplies the measurability required in both transfers.



A state-space field with finite square energy under \(\rho_t\) has the same square energy when composed with the proposed coordinate on the joint pair. This is pushforward change of variables.



Earlier results used here: \cite[Section~3.34.1]{LangevinLean2026ULASC}; \cite[Section~3.36.5]{LangevinLean2026ULASC}.
\subsection{\texorpdfstring{Integrability of the score-\allowbreak{}test pairing}{Integrability of the score-test pairing}}
\label{report-step-UnadjustedLangevin-LSI-GaussianScorePairing-lean-4}


The component score and the state field \(V(y)\) are both square integrable under the actual pair law. Holder's inequality makes the product of their norms integrable, and \(|\langle S,V\rangle|\le\|S\|\|V\|\) proves integrability of their inner product. The exact pair-density identity then rewrites this law as \(K_t(x,y)\mu(dx)dy\). The weighted-density equivalence gives joint integrability of \(K_t\langle S_{t,x},V(y)\rangle\) under \(\mu\otimes dy\), which is the hypothesis needed for Fubini in the averaging calculation.



Cauchy–Schwarz on the pair probability bounds the absolute pairing integral by the product of component-score and state-field \(L^2\) norms. Hence the signed pairing is integrable before its order is exchanged.



Earlier results used here: \cite[Section~3.34.1]{LangevinLean2026ULASC}; \cite[Section~3.37.14]{LangevinLean2026ULASC}.
\subsection{\texorpdfstring{The weighted mixture-\allowbreak{}score identity}{The weighted mixture-score identity}}
\label{report-step-UnadjustedLangevin-LSI-GaussianScorePairing-lean-5}


The already constructed spatial derivative of the Gaussian mixture identifies its score as the posterior mean of the component score. The posterior density in \(x\) is \(K_t(x,y)/p_t(y)\), where \(p_t(y)=\int K_t(x,y)d\mu(x)>0\). Substitute this explicit weighted-density formula into the posterior integral and multiply by \(p_t(y)\). The common denominator cancels pointwise, giving \(p_t(y)\nabla\log p_t(y)=\int K_t(x,y)S_{t,x}(y)d\mu(x)\). The proof uses the actual Gaussian normalizers in both precision and step-size notation, rather than assuming an unproved Bayes identity.



Integrating the component gradient over the starting law gives \(\nabla p_t\). Dividing by the positive mixture identifies its averaged score with \(\nabla\log p_t\), the posterior-score identity in weighted form.



Earlier results used here: \cite[Section~3.37.12]{LangevinLean2026ULASC}; \cite[Section~3.37.4]{LangevinLean2026ULASC}.
\subsection{\texorpdfstring{Integrability of the kernel-\allowbreak{}weighted score pairing}{Integrability of the kernel-weighted score pairing}}
\label{report-step-UnadjustedLangevin-LSI-GaussianScorePairing-lean-6}


The component score is \(L^2\) under the actual pair probability law, hence is also \(L^1\). Identify that law with its explicit density \(K_t\) under \(\mu\otimes dy\). The vector-valued weighted-density integrability theorem gives \((x,y)\mapsto K_t(x,y)S_{t,x}(y)\in L^1(\mu\otimes dy)\). In particular, Fubini supplies integrability of the vector section in \(x\) for almost every output \(y\); this is needed before moving an inner product through its integral.



The Gaussian score has finite square norm and the tested state field has finite square norm. Their product is therefore integrable, supplying the precise Fubini hypothesis for the following averaging identity.



Earlier results used here: \cite[Section~3.37.4]{LangevinLean2026ULASC}.
\subsection{\texorpdfstring{Averaging the Gaussian component score}{Averaging the Gaussian component score}}
\label{report-step-UnadjustedLangevin-LSI-GaussianScorePairing-lean-7}


For almost every \(y\), vector-section integrability permits the linear functional \(w\mapsto\langle w,V(y)\rangle\) to commute with integration. The pointwise weighted-score identity then gives \(\int K_t(x,y)\langle S_{t,x}(y),V(y)\rangle d\mu(x)=p_t(y)\langle\nabla\log p_t(y),V(y)\rangle\).



The scalar pairing is jointly integrable by the preceding Holder argument, so Fubini exchanges the initial-state and output integrals. Integrate the displayed identity in \(y\). Finally \(p_t(y)dy=\rho_t(dy)\), with its explicit measurable finite density, changes the weighted Lebesgue integral into the actual-law integral. Thus \(\int K_t\langle S_{t,x},V\rangle d\mu\,dy=\int\langle\nabla\log p_t,V\rangle d\rho_t\).



Average the component score at a fixed observation, then integrate its pairing with the state test. The mixture-gradient identity converts the result into the pairing with the actual Lebesgue score under \(\rho_t\).

\subsection{\texorpdfstring{The score identity on the independent-\allowbreak{}noise space}{The score identity on the independent-noise space}}
\label{report-step-UnadjustedLangevin-LSI-GaussianScorePairing-lean-8}


Apply the measurable pair map \(T_t(x,z)=(x,Y_t(x,z))\) to the continuous pairing \(\langle S_{t,x}(y),V(y)\rangle\). The map-integral theorem expresses its actual pair-law integral as an integral over the fixed independent initial/\allowbreak{}noise space. The exact pair-density theorem expresses the same integral as the Gaussian-weighted Lebesgue pairing just evaluated.



Substitute \(Y_t=c_t(x)+\sqrt{2t}z\) in the component score. The residual becomes \(\sqrt{2t}z\), so the result is \(\mathbb E\langle-(\sqrt{2t}/(2t))Z,V(Y_t)\rangle=\mathbb E_{\rho_t}\langle\nabla\log p_t,V\rangle\). With \(V=\nabla\log r_t\), this identity turns the Gaussian-noise part of the raw entropy flux into the Lebesgue-score pairing needed for the full Fisher term.



Use the joint-law identity to rewrite the same integral on the original independent state-noise space. This lets the entropy flux be combined with the exact displacement and gradient-error estimates already proved there.



Earlier results used here: \cite[Section~3.34.1]{LangevinLean2026ULASC}; \cite[Section~3.36.5]{LangevinLean2026ULASC}; \cite[Section~3.37.4]{LangevinLean2026ULASC}.
\section{\texorpdfstring{Spend one quarter under Fisher information on drift error}{Spend one quarter under Fisher information on drift error}}
\label{proof-module-UnadjustedLangevin-LSI-DriftError-lean}


Young’s inequality \(\langle a,b\rangle\le\frac14\|a\|^2+\|b\|^2\) bounds the error pairing in entropy production. Weighted energy integrability yields an actual integral inequality \(H\prime\le-\frac34J+\mathcal E\).



The error term in entropy production pairs the drift discrepancy \(e\) with the relative score \(s\). The weighted square identity \((\|e\|-\|s\|/2)^2\ge0\) gives \(|\langle e,s\rangle|\le\|e\|^2+\|s\|^2/4\). Its absolute-value form first proves that the pairing is integrable from the two energy bounds. Its signed form then absorbs one quarter of Fisher information, leaving \(-3I/4\) as a strictly negative dissipation contribution.

\subsection{\texorpdfstring{Young absorption preserving three quarters of Fisher information}{Young absorption preserving three quarters of Fisher information}}
\label{report-step-UnadjustedLangevin-LSI-DriftError-lean-1}


Cauchy–Schwarz gives \(\langle e,s\rangle\le\|e\|\|s\|\). Expand the nonnegative square \((\|e\|-\|s\|/2)^2\ge0\) to get \(\|e\|\|s\|\le\|e\|^2+\|s\|^2/4\). This is the Young coefficient that preserves three quarters of Fisher information.



Expanding the square gives \(\|e\|^2-\|e\|\|s\|+\|s\|^2/4\ge0\). Thus the quarter coefficient is exact for this chosen absorption, and leaves a unit coefficient on the discrepancy energy.

\subsection{\texorpdfstring{Integrability of the weighted error-\allowbreak{}score pairing}{Integrability of the weighted error-score pairing}}
\label{report-step-UnadjustedLangevin-LSI-DriftError-lean-2}


Multiply the absolute Young bound by the nonnegative weight \(r\). It dominates \(|r\langle e,s\rangle|\) by \(r\|s\|^2/4+r\|e\|^2\). Both energy functions are integrable, so the assumed measurability and integrable domination prove integrability of the cross term.



The absolute-value domination is established before integrating the signed pairing. Measurability plus the two finite energy integrals therefore gives a legitimate ordinary cross-term integral.

\subsection{\texorpdfstring{The forced entropy-\allowbreak{}production inequality}{The forced entropy-production inequality}}
\label{report-step-UnadjustedLangevin-LSI-DriftError-lean-3}


First derive cross-term integrability from the two finite energies and its measurability. The forced weak identity tested by \(\log r+1\) gives production \(-I+\int r\langle e,s\rangle\). Insert the integrated Young bound to obtain \(-\frac34I+\int r\|e\|^2\). The entropy differential inequality is therefore proved from the stated spatial weak identity.



Insert the admissible entropy test only after its forcing pairing has been proved integrable. The resulting inequality is derived from the spatial weak equation and finite energies, not postulated as a scalar property of the entropy curve.

\subsection{\texorpdfstring{Absolute-\allowbreak{}value Young domination}{Absolute-value Young domination}}
\label{report-step-UnadjustedLangevin-LSI-DriftError-lean-4}


Use the absolute Cauchy–Schwarz inequality \(|\langle e,s\rangle|\le\|e\|\|s\|\), followed by the same square expansion. The absolute-value version controls integrability of the signed cross term before its integral is estimated.



Replacing the pairing by its absolute value is used for domination and Fubini. The same nonnegative-square calculation supplies the upper bound, so no additional moment beyond the two square energies is needed.

\subsection{\texorpdfstring{The integrated weighted Young inequality}{The integrated weighted Young inequality}}
\label{report-step-UnadjustedLangevin-LSI-DriftError-lean-5}


Multiply the pointwise Young inequality by \(r\ge0\), then integrate. The supplied integrability permits splitting the upper integral and pulling out \(1/4\). Thus \(\int r\langle e,s\rangle\le\frac14\int r\|s\|^2+\int r\|e\|^2\). No probability normalization of the weight is needed.



Nonnegativity of \(r\) preserves the pointwise upper bound. Integral linearity then separates the two finite energies, retaining exactly the Fisher fraction required for the subsequent LSI argument.

\section{\texorpdfstring{Prove the actual Fisher dissipation and sharp frozen-\allowbreak{}gradient forcing}{Prove the actual Fisher dissipation and sharp frozen-gradient forcing}}
\label{proof-module-UnadjustedLangevin-LSI-GaussianEntropyProduction-lean}


The actual Gaussian entropy chain rule gives an almost-everywhere derivative represented by a flux on the independent initial/\allowbreak{}noise probability space. Averaging the genuine Gaussian component score identifies this flux with \(-I_t+\mathbb E\langle\nabla U(Y_t)-\nabla U(X),\nabla\log r_t(Y_t)\rangle\). Young's inequality absorbs a quarter of \(I_t\); the already proved exact Gaussian displacement estimate bounds the remaining error. The result is the unconditional analytic production inequality \(H'_t\le-\tfrac34I_t+t^2L^2G_0+2tdL^2\) almost everywhere at positive times, with no interpolation regularity or PDE data hypotheses.



The actual entropy identity has one negative Fisher term and one drift-error pairing. The two fields have finite square energy on the common independent-noise space, so their pairing is integrable. Young's inequality bounds it by one quarter of Fisher information plus the drift-error energy. The displacement identity and gradient Lipschitzness then give \(H'(t)\le-3I_t/4+t^2L^2G_0+2tdL^2\) for almost every positive time, on one full-measure set suitable for time integration.

\subsection{\texorpdfstring{Justify inner-\allowbreak{}product integrability from two L\ensuremath{{}^{2}} fields}{Justify inner-product integrability from two L² fields}}
\label{report-step-UnadjustedLangevin-LSI-GaussianEntropyProduction-lean-1}


For vector fields \(F,G\in L^2(P)\), Hölder's inequality puts \(\|F\|\|G\|\) in \(L^1(P)\). Since \(|\langle F,G\rangle|\le\|F\|\|G\|\), their inner product is integrable.



The initial gradient, current gradient, Gaussian component score, and relative score all have the square energies proved above. Each needed pairing is therefore integrable by Cauchy–Schwarz. Their signed integrals may consequently be split and recombined in the entropy-flux algebra.



Cauchy–Schwarz proves absolute integrability of the pairing on the common independent-noise probability. This permits the exact entropy flux to be split into its Fisher and discrepancy contributions as ordinary integrals.

\subsection{\texorpdfstring{Identify relative-\allowbreak{}score energy on the independent noise space with actual Fisher information}{Identify relative-score energy on the independent noise space with actual Fisher information}}
\label{report-step-UnadjustedLangevin-LSI-GaussianEntropyProduction-lean-2}


Let \(S_t(X,Z)=\nabla\log r_t(Y_t)\) and \(P=\mu\otimes\mathcal N(0,I)\). The exact update-law identity gives

\[\mathbb E_P\|S_t\|^2=\int\|\nabla\log r_t\|^2\,d\rho_t=I_t.\]

Pushforward change of variables gives \(\mathbb E\|S_t(Y_t)\|^2=\int\|S_t\|^2d\rho_t=I_t\). The relative-score representative is the same one constructed by Gaussian smoothing.



Earlier results used here: \cite[Section~3.26.3]{LangevinLean2026ULASC}; \cite[Section~3.34.1]{LangevinLean2026ULASC}; \cite[Section~3.36.5]{LangevinLean2026ULASC}.
\subsection{\texorpdfstring{Match the measurable state-\allowbreak{}space score field to the same finite Fisher value}{Match the measurable state-space score field to the same finite Fisher value}}
\label{report-step-UnadjustedLangevin-LSI-GaussianEntropyProduction-lean-3}


The jointly measurable field used in time-space estimates agrees with the actual relative score at each positive time. Consequently

\[\int\|S_t(y)\|^2\,d\rho_t(y)=I_t.\]

The measurable state-space score equals the actual logarithmic gradient at positive time. Pushing its square energy through the actual proposal map preserves its integral. Thus this field, its independent-noise composition, and the score used in LSI all have the same finite Fisher energy.



At positive time the jointly measurable field equals the actual relative score pointwise. Its integrated energy therefore matches the finite Fisher quantity used in both LSI and production.



Earlier results used here: \cite[Section~3.26.3]{LangevinLean2026ULASC}; \cite[Section~3.36.5]{LangevinLean2026ULASC}.
\subsection{\texorpdfstring{Integrate Young absorption against the actual frozen-\allowbreak{}gradient discrepancy}{Integrate Young absorption against the actual frozen-gradient discrepancy}}
\label{report-step-UnadjustedLangevin-LSI-GaussianEntropyProduction-lean-4}


Write \(D_t=\nabla U(Y_t)-\nabla U(X)\). The pointwise inequality

\[\langle D_t,S_t\rangle\le\frac14\|S_t\|^2+\|D_t\|^2\]

is integrable because both squared norms have already been proved integrable. Their expectations satisfy \(\mathbb E\|S_t\|^2=I_t\) and

\[\mathbb E\|D_t\|^2\le L^2\bigl(t^2G_0+2td\bigr),\qquad G_0=\mathbb E\|\nabla U(X)\|^2.\]

Apply \(\langle e,s\rangle\le\|e\|^2+\|s\|^2/4\) pointwise and integrate on the independent-noise space. The score energy becomes \(I_t/4\), while the error energy retains its sharp Gaussian bound.



Earlier results used here: \cite[Section~3.26.3]{LangevinLean2026ULASC}; \cite[Section~3.36.5]{LangevinLean2026ULASC}.
\subsection{\texorpdfstring{Convert the actual Gaussian entropy flux into Fisher dissipation and gradient error}{Convert the actual Gaussian entropy flux into Fisher dissipation and gradient error}}
\label{report-step-UnadjustedLangevin-LSI-GaussianEntropyProduction-lean-5}


The raw entropy flux is \(-\mathbb E\langle g(X)-\sqrt{2t}Z/(2t),S_t\rangle\). The averaged Gaussian score identity gives

\[\mathbb E\left\langle-\frac{\sqrt{2t}}{2t}Z,S_t\right\rangle=\int\langle\nabla\log p_t,\nabla\log r_t\rangle\,d\rho_t.\]

Since \(\nabla\log r_t=\nabla U+\nabla\log p_t\), the right side equals \(I_t-\mathbb E\langle\nabla U(Y_t),S_t\rangle\). Substituting and collecting the two gradient terms yields

\[Q_t=-I_t+\mathbb E\langle\nabla U(Y_t)-\nabla U(X),S_t\rangle.\]

Average the component Gaussian score using its established pairing identity. The flux then becomes \(-I_t+\mathbb E\langle\nabla U(Y_t)-\nabla U(X),S_t(Y_t)\rangle\), with no unproved conditional-drift choice.



Earlier results used here: \cite[Section~3.26.3]{LangevinLean2026ULASC}; \cite[Section~3.34.1]{LangevinLean2026ULASC}; \cite[Section~3.36.5]{LangevinLean2026ULASC}.
\subsection{\texorpdfstring{Retain three quarters of actual Fisher dissipation}{Retain three quarters of actual Fisher dissipation}}
\label{report-step-UnadjustedLangevin-LSI-GaussianEntropyProduction-lean-6}


Insert the integrated Young estimate into the exact entropy-flux decomposition:

\[Q_t=-I_t+\mathbb E\langle D_t,S_t\rangle\le-\frac34I_t+t^2L^2G_0+2tdL^2.\]

The proof rewrites the flux with the decomposition of entropy flux into Fisher dissipation and the drift-discrepancy pairing, obtains the already justified numerical cross-term bound, and collects the scalar terms. The target's probability normalization and finite initial gradient energy suffice for this pointwise flux estimate. LSI is needed later to relate Fisher information to entropy, rather than to justify Young absorption.



Subtracting the absorbed quarter from the full negative Fisher term gives \(-3I_t/4\). The remaining discrepancy energy is at most \(t^2L^2G_0+2tdL^2\), giving the desired scalar forcing.



Earlier results used here: \cite[Section~3.26.3]{LangevinLean2026ULASC}.
\subsection{\texorpdfstring{Obtain the genuine positive-\allowbreak{}time entropy derivative inequality on one common full-\allowbreak{}measure set}{Obtain the genuine positive-time entropy derivative inequality on one common full-measure set}}
\label{report-step-UnadjustedLangevin-LSI-GaussianEntropyProduction-lean-7}


The actual KL entropy is absolutely continuous on every compact interval inside \((0,\infty)\), and its derivative equals the actual raw flux almost everywhere there. To use this fact on an entire finite range \((0,T)\), choose the countable intervals

\[\left(\frac1{n+1},T+1\right),\qquad n=0,1,2,\ldots.\]

Intersect their full-measure derivative sets. Every \(t\in(0,T)\) belongs to at least one such interval because \(1/(n+1)\to0\). At that time the actual derivative equals the raw flux, so the pointwise flux bound gives

\[H'_t\le-\frac34 I_t+t^2L^2G_0+2tdL^2\quad\text{for almost every }t\in(0,T).\]

The local entropy identity supplies its derivative almost everywhere, and the pointwise flux bound holds at every positive time with finite energies. Intersecting these full-measure time conditions yields one almost-everywhere derivative inequality suitable for integration.



Earlier results used here: \cite[Section~3.26.3]{LangevinLean2026ULASC}.
\section{\texorpdfstring{The full Gaussian ULA entropy theorem}{The full Gaussian ULA entropy theorem}}
\label{proof-module-UnadjustedLangevin-LSI-UnconditionalTheorem-lean}


For a differentiable potential with an \(L\)-Lipschitz gradient, its normalized Gibbs target satisfying LSI with constant \(\alpha>0\), and a finite-KL initial probability law, this module proves the actual ULA bound

\[\mathrm{KL}(\rho_k\mid\nu)\le e^{-\alpha hk}\mathrm{KL}(\rho_0\mid\nu)+8hdL^2/\alpha,\quad0<h\le\alpha/(4L^2).\]

The final conclusion includes finiteness of the actual divergence. All score moments, density calculus, weak equations, entropy endpoints, and numerical compatibility estimates are supplied by the preceding modules. The theorem applies also in dimension zero.



The actual-law proof begins with finite initial KL. Score exponential integrability supplies the initial gradient moment; Gaussian smoothing supplies positive \(C^1\) densities and finite Fisher energy at positive time; LSI then supplies finite KL. The bounded-test entropy identity, its initial majorant, and compatibility \(\alpha\le L\) prove the scalar one-step estimate. Induction preserves finite KL along the actual ULA iterates, so the final geometric recurrence represents their true entropies.

\subsection{\texorpdfstring{Handle the singleton state space and a zero smoothness parameter}{Handle the singleton state space and a zero smoothness parameter}}
\label{report-step-UnadjustedLangevin-LSI-UnconditionalTheorem-lean-1}


On a singleton state space every probability measure is equal: any measurable set either contains the unique point and is the whole space, or is empty. Measure extensionality and probability normalization therefore identify any two laws. In particular their KL divergence is zero.



For the substantive entropy calculation one works with \(L>0\), so that the step restriction and exponential score moment have their usual positive denominators. In positive dimension, \(\alpha\le L\) and \(\alpha>0\) guarantee this positivity. The singleton case has already been resolved directly with zero KL, so it needs no division by a smoothness constant.



When the state space is a singleton, every probability law is the same point mass and KL is zero. In the nondegenerate case, the compatibility argument rules out zero smoothness when the LSI constant is positive, permitting the subsequent divisions by \(L\).

\subsection{\texorpdfstring{Prove finiteness at every actual ULA iterate}{Prove finiteness at every actual ULA iterate}}
\label{report-step-UnadjustedLangevin-LSI-UnconditionalTheorem-lean-2}


Start from the supplied finite initial KL and induct on \(k\). If \(\rho_k\) has finite KL, the direct entropy/\allowbreak{}score estimate gives \(\nabla U\in L^2(\rho_k)\). For the actual positive-time Gaussian interpolation, the Lebesgue score energy is at most \(d/(2h)\), and the propagated gradient moment is finite. Thus relative Fisher information is finite. LSI with \(\alpha>0\) then gives finite KL of the output.



The endpoint law is exactly \(\rho_{k+1}\), so this argument closes the finite-KL induction. The finiteness conclusion needs only positivity of the step; its smaller quantitative restriction is used later for contraction and forcing absorption.



The induction begins with the assumed finite initial KL. Finite current KL supplies its score moment; the positive-time Gaussian transition has finite Fisher energy; LSI then gives finite next KL. Thus each actual entropy term is finite before the recurrence is iterated.



Earlier results used here: \cite[Section~3.26.3]{LangevinLean2026ULASC}.
\subsection{\texorpdfstring{Derive all scalar step restrictions}{Derive all scalar step restrictions}}
\label{report-step-UnadjustedLangevin-LSI-UnconditionalTheorem-lean-3}


In positive dimension, the bounded sinusoidal tests in compatibility give \(\alpha\le L\). Since \(\alpha>0\), also \(L>0\). Multiplying the original restriction by \(4L^2>0\) gives \(4hL^2\le\alpha\le L\), hence \(hL\le1/4\). Multiplying \(\alpha\le L\) by \(h>0\) also gives \(\alpha h\le1/4\).



The final theorem splits off the singleton case before invoking this positive-dimensional argument. No separate compatibility assumption is required.



From \(h\le\alpha/(4L^2)\) and \(\alpha\le L\), derive \(hL\le1/4\) and \(\alpha h\le1/4\). These numerical consequences are checked before invoking the scalar absorption and geometric-denominator lemmas.



Earlier results used here: \cite[Section~3.26.3]{LangevinLean2026ULASC}.
\subsection{\texorpdfstring{Apply LSI to the actual positive-\allowbreak{}time density}{Apply LSI to the actual positive-time density}}
\label{report-step-UnadjustedLangevin-LSI-UnconditionalTheorem-lean-4}


Finite initial KL supplies the initial gradient moment. The actual Gaussian mixture is positive C\ensuremath{{}^{1}} relative to Gibbs, has finite KL, and has finite Fisher information at every \(t>0\). Apply the original extended-valued LSI to this density. Since both sides are finite, order is preserved on passing to their real values:

\[2\alpha\,\mathrm{KL}(\rho_t\mid\nu)\le I_t.\]

The integral form identifies \(I_t\) with \(\int\|\mathcal S(t,y)\|^2d\rho_t(y)\), using the proved equality between the measurable score field and \(\nabla\log r_t\), and the equality \(r_t\nu=\rho_t\). The second version keeps the extended-Fisher notation for the entropy-production assembly.



At positive elapsed time, the constructed relative density is positive, \(C^1\), and normalized, with finite Fisher and KL. It therefore belongs to the original LSI test class, giving \(2\alpha H(t)\le I_t\) for the actual law.



Earlier results used here: \cite[Section~3.26.3]{LangevinLean2026ULASC}.
\subsection{\texorpdfstring{Supply the actual analytic inputs to scalar comparison}{Supply the actual analytic inputs to scalar comparison}}
\label{report-step-UnadjustedLangevin-LSI-UnconditionalTheorem-lean-5}


Fix any finite-KL probability law \(\mu\). The score-exponential argument gives the integrability of \(\|\nabla U\|^2\) and

\[G_0\le4L^2H_\nu(\mu)/\alpha+2dL.\]

The bounded-log entropy identity proves local absolute continuity of the actual entropy \(H(t)\) on every \([s,h]\) with \(s>0\). The actual Gaussian weak equation and Young estimate give \(H^{\prime}\le-3I/4+t^2L^2G_0+2tdL^2\) almost everywhere. The preceding LSI estimate supplies \(2\alpha H\le I\).



The explicit initial majorant \(M\) is continuous at zero, satisfies \(M(0)=H_\nu(\mu)\), and bounds \(H(t)\) for sufficiently small positive \(t\). These are exactly the premises of the integrating-factor comparison for a locally absolutely continuous entropy curve, giving

\[H_\nu(\mu P_h)\le e^{-\alpha h}H_\nu(\mu)+6h^2dL^2.\]

The singleton case instead uses equality of probability laws and nonnegativity of the right side. No density path, weak equation, entropy derivative, or moment estimate is a theorem argument.



The actual entropy curve has local absolute continuity, the almost-everywhere production bound, and an upper initial majorant continuous at zero. Together with the entropy-derived score moment, these supply every premise of the scalar one-step comparison.



Earlier results used here: \cite[Section~3.26.3]{LangevinLean2026ULASC}.
\subsection{\texorpdfstring{Derive the actual discrete recurrence}{Derive the actual discrete recurrence}}
\label{report-step-UnadjustedLangevin-LSI-UnconditionalTheorem-lean-6}


The law \(\rho_k\) is a probability measure by the proved Gaussian Markov construction. Its KL is finite by the induction above. Apply the single-transition estimate with \(\mu=\rho_k\). The interpolation endpoint is precisely \(\rho_{k+1}\), yielding

\[H_{k+1}\le e^{-\alpha h}H_k+6h^2dL^2.\]

entropy of the actual ula law is the real value of the actual KL of iteration of the gaussian transition. The recurrence is derived here; it is not supplied as a premise.



Use the exact interpolation endpoints to replace the scalar initial and final entropies by \(H_k\) and \(H_{k+1}\). The comparison then gives \(H_{k+1}\le e^{-\alpha h}H_k+6h^2dL^2\) without any separately assumed recurrence.



Earlier results used here: \cite[Section~3.26.3]{LangevinLean2026ULASC}.
\subsection{\texorpdfstring{Sum the recurrence and retain the finite-\allowbreak{}KL guard}{Sum the recurrence and retain the finite-KL guard}}
\label{report-step-UnadjustedLangevin-LSI-UnconditionalTheorem-lean-7}


Return the already-proved output finiteness as the first component of the conclusion. In positive dimension, the step estimates give \(\alpha h\le1/4\). The scalar geometric-sum theorem applied to the actual recurrence gives

\[H_k\le e^{-\alpha hk}H_0+\frac{6h^2dL^2}{1-e^{-\alpha h}}\le e^{-\alpha hk}H_0+\frac{8hdL^2}{\alpha}.\]

The final inequality uses \(1-e^{-x}\ge3x/4\) for \(0\le x\le1/4\). In dimension zero, equality of probability measures makes the left side zero and the result follows by nonnegativity.



The final hypotheses are differentiability, globally Lipschitz gradient, the Gibbs LSI with positive constant, a positive admissible step, and finite initial KL. Probability normalization of the target is included in LSI. Every density, moment, weak-equation, and entropy-regularity property used in the recurrence has been derived above from these assumptions.



Apply the geometric recurrence estimate to the now finite sequence of true KL values. The positive denominator bound converts the forcing sum into \(8hdL^2/\alpha\), yielding the final statement under only the displayed potential, LSI, step, and initial-entropy assumptions.



Earlier results used here: \cite[Section~3.26.3]{LangevinLean2026ULASC}.
\section{\texorpdfstring{Conditional LSI theorem for the actual ULA laws}{Conditional LSI theorem for the actual ULA laws}}
\label{proof-module-UnadjustedLangevin-LSI-Theorem-lean}


The actual iterated Gaussian ULA laws satisfy \(H_k\le e^{-\alpha hk}H_0+8hdL^2/\alpha\), once the displayed interpolation identification, finite KL, entropy/\allowbreak{}Fisher dissipation, LSI, and moment hypotheses are supplied. This endpoint does not construct all those analytic data from LSI and smoothness alone.



For a single step starting at \(\rho_k\), write \(H(t)=\operatorname{KL}(\rho_{k,t}\mid\nu)\), \(I(t)=I_\nu(\rho_{k,t})\), and \(G_k=\int\|\nabla U\|^2d\rho_k\). The supplied path hypotheses yield \(H'\le-3I/4+t^2L^2G_k+2tdL^2\), while LSI gives \(I\ge2\alpha H\). The endpoint identities connect the resulting continuous estimate to \(H_k\) and \(H_{k+1}\). This intermediate theorem is conditional on its displayed analytic path assumptions; the later actual-law theorem constructs the necessary positive-time inputs.

\subsection{\texorpdfstring{The target estimate and the conditional interpolation theorem}{The target estimate and the conditional interpolation theorem}}
\label{report-step-UnadjustedLangevin-LSI-Theorem-lean-1}


For every index, derive the one-step recurrence from the preceding interpolation theorem. The numerical conditions \(\alpha\le L\), \(h\le\alpha/(4L^2)\) imply \(\alpha h\le1/4\). Apply the geometric-series bound for the entropy recurrence and rewrite the initial value with the equality of the zeroth law with the initial law. This gives \(H_k\le e^{-\alpha hk}\operatorname{KL}(\mu\mid\pi_U)+8hdL^2/\alpha\), conditional on the displayed analytic path hypotheses.



The recurrence is applied to the actual sequence \(H_k=\operatorname{KL}(\rho_k\mid\nu)\), after finiteness has been proved from the supplied path. The upper restriction on \(h\) and \(\alpha\le L\) give \(\alpha h\le\alpha^2/(4L^2)\le1/4\). This supplies the denominator estimate needed to sum the discrete forcing into the stated floor. The result at this intermediate stage retains every displayed analytic path assumption.



Earlier results used here: \cite[Section~3.26.3]{LangevinLean2026ULASC}.
\subsection{\texorpdfstring{Why finite KL and Fisher information are separate obligations}{Why finite KL and Fisher information are separate obligations}}
\label{report-step-UnadjustedLangevin-LSI-Theorem-lean-2}


At time zero the supplied relative-density law identity gives \(r_{k,0}\nu=\rho_k\), because the interpolation identity at elapsed time zero returns its input law. Substitute this equality into finite KL of the time-zero density path. It yields finite actual KL of the algorithm law at every index.



Finite Fisher information at positive time controls entropy through LSI, whereas finite initial KL supplies the numerical initial value in the recurrence. These are different tasks. Identifying the time-zero relative law with \(\rho_k\) transfers the supplied finite-entropy property to the algorithm law, before its real-valued entropy is used in the estimate.

\subsection{\texorpdfstring{Combine production, LSI, and the initial gradient moment}{Combine production, LSI, and the initial gradient moment}}
\label{report-step-UnadjustedLangevin-LSI-Theorem-lean-3}


For the actual current law, write \(G_k=\int\|\nabla U(x)\|^2d\rho_k(x)\) once square-gradient integrability is known. The separate score-moment hypothesis ensures this integral is a finite energy. Its entropy-dependent upper bound is then used to control forcing over the next interpolation step.



The energy \(G_k\) concerns the gradient at the starting law of step \(k\). It is therefore constant as the elapsed interpolation time varies. Its later bound by \((4L^2/\alpha)H_k+2dL\) separates an entropy-dependent contribution from a dimension-dependent one, allowing them to be absorbed and accumulated differently in scalar comparison.

\subsection{\texorpdfstring{The direct endpoint constructs the recurrence from the path hypotheses}{The direct endpoint constructs the recurrence from the path hypotheses}}
\label{report-step-UnadjustedLangevin-LSI-Theorem-lean-4}


Identify the entropy path's endpoints using the interpolation identity at elapsed time zero and the identification of the elapsed-time endpoint with the next iterate: its time-zero entropy is \(H_k\), and its time-\(h\) entropy is \(H_{k+1}\). The law identity and probability preservation prove that every intermediate relative density is normalized.



Apply the real LSI inequality using the explicit positive smooth density and finite KL and Fisher information hypotheses. The scalar the scalar entropy comparison on one step then combines it with the supplied entropy derivative, production, and gradient-moment estimates. Rewriting the two endpoints yields \(H_{k+1}\le e^{-\alpha h}H_k+6h^2dL^2\). Those analytic interpolation premises remain explicit in this theorem.



For each \(k\), set \(H(0)=H_k\) and \(H(h)=H_{k+1}\) using the two exact law endpoints. The scalar comparison combines \(-3I/4\le-(3\alpha/2)H\) with the stated score and drift-energy bounds. Its entropy-dependent forcing is absorbed by the exponential-rate gap, leaving the uniform step cost \(6h^2dL^2\).



Earlier results used here: \cite[Section~3.26.3]{LangevinLean2026ULASC}.
\section{\texorpdfstring{Differentiate the entropy integrand under domination}{Differentiate the entropy integrand under domination}}
\label{proof-module-UnadjustedLangevin-LSI-EntropyDerivative-lean}


The scalar identity \((r\log r)\prime=r\prime(\log r+1)\), for \(r>0\), lifts to an integral derivative under the explicitly supplied common local integrable bound. Domination is a premise, not a consequence of Gaussian convolution in this module.



The scalar derivative of \(u\log u\) at a positive value is \(\log u+1\). For a time-dependent relative density this produces \(\partial_t r_t(\log r_t+1)\). Passing from this pointwise identity to a derivative of the entropy integral requires a single integrable spatial majorant valid throughout a time neighborhood. Pointwise differentiability alone does not justify the exchange of derivative and integral, particularly where a density becomes small and its logarithm is large.

\subsection{\texorpdfstring{Differentiating a positive density times its logarithm}{Differentiating a positive density times its logarithm}}
\label{report-step-UnadjustedLangevin-LSI-EntropyDerivative-lean-1}


If \(r(t)>0\), the product and logarithm chain rules give \((r\log r)'=\dot r\log r+r\dot r/r\). Positivity supplies \(r(t)\ne0\), so cancellation gives \(\dot r(\log r+1)\). The cancellation is justified by strict positivity of the density.



The additional \(+1\) comes from differentiating the density factor in \(r\log r\). Keeping this term until the mass derivative is accounted for avoids silently assuming normalization outside the time interval on which the path is known to be a probability.

\subsection{\texorpdfstring{Differentiating the entropy integral under domination}{Differentiating the entropy integral under domination}}
\label{report-step-UnadjustedLangevin-LSI-EntropyDerivative-lean-2}


For almost every spatial point, apply the scalar chain rule \((r\log r)\prime=r\prime(\log r+1)\) throughout a common time neighborhood. Its derivative is bounded in norm by the supplied integrable spatial envelope, uniformly over that neighborhood.



The dominated differentiation theorem uses this bound, integrability of \(r_t\log r_t\), and measurability of both integrands. It concludes both integrability of \(\dot r_t(\log r_t+1)\) and \(\frac d{dt}\int r_t\log r_t d\nu=\int\dot r_t(\log r_t+1)d\nu\). The derivative of the integral is a conclusion, not a hypothesis.



The domination is uniform over a time neighborhood, so difference quotients can be controlled by the same spatial integrable bound. The resulting derivative identity is a conclusion of dominated differentiation, rather than a restatement of assumed entropy differentiability.

\section{\texorpdfstring{Local regularity needed for actual entropy calculus}{Local regularity needed for actual entropy calculus}}
\label{proof-module-UnadjustedLangevin-LSI-Regularity-lean}


A local entropy regularity specification supplies positive densities, time derivatives, and local integrable domination. Its consequences prove finite KL, continuity, derivatives and a production inequality for the identified density entropy; existence of these packages is retained as a premise.



The local regularity assumptions are sufficient conditions for differentiating \(\int r_t\log r_t\,d\nu\): they specify positivity, pointwise time differentiability, and an integrable neighborhood majorant for the derivative integrand. Probability normalization then identifies that integral with the actual relative entropy. The spatial weak equation supplies the value of the derivative after this calculus step; it is not part of the conclusion of dominated differentiation itself. At a boundary, only the derivative from within the admissible time interval is used.

\subsection{\texorpdfstring{Local hypotheses for entropy differentiation}{Local hypotheses for entropy differentiation}}
\label{report-step-UnadjustedLangevin-LSI-Regularity-lean-1}


At a fixed time \(t\), store a common time neighborhood, an integrable spatial envelope for \(\dot r_u(\log r_u+1)\), local pointwise positivity/\allowbreak{}differentiability, and measurability/\allowbreak{}integrability of the entropy integrands. These are sufficient data for dominated differentiation. The structure contains no derivative of the entropy integral and no entropy-production estimate.



The local neighborhood may depend on the chosen time, but within it one spatial majorant controls all relevant differentiated integrands. This is sufficient for a local differentiation theorem; no global-in-time bound is stored in these hypotheses.

\subsection{\texorpdfstring{The derivative supplied by local entropy regularity}{The derivative supplied by local entropy regularity}}
\label{report-step-UnadjustedLangevin-LSI-Regularity-lean-2}


Pass the stored neighborhood, envelope, and pointwise derivative fields to the dominated differentiation formula for \(\int r_t\log r_t\,d\nu\). It concludes both integrability of the differentiated integrand and the derivative formula for \(\int r_t\log r_t d\nu\). This is the calculus theorem supplied by the local data, not a field assumed in them.



Apply the scalar chain rule pointwise and then the neighborhood domination theorem. The outcome includes integrability of the derivative integrand itself, allowing its integral to serve as the ordinary entropy derivative.

\subsection{\texorpdfstring{Production from an admissible forced weak equation}{Production from an admissible forced weak equation}}
\label{report-step-UnadjustedLangevin-LSI-Regularity-lean-3}


Use local dominated differentiation to compute \(H'(t)=\int\dot r_t(\log r_t+1)d\nu\). The forced weak equation, weighted Fisher integrability, and drift-energy integrability then permit the entropy test and integrated Young inequality. The result is \(H'(t)\le-3I(t)/4+E(t)\), with \(E(t)=\int r_t\|e_t\|^2d\nu\). No numerical bound on \(E(t)\) is assumed in this lemma.



The spatial weak equation identifies the derivative integral with a negative Fisher term plus drift forcing. Applying weighted Young gives \(-I+I/4+E=-3I/4+E\), retaining the exact error energy rather than assuming its eventual numerical bound.

\subsection{\texorpdfstring{Finite KL from local regularity}{Finite KL from local regularity}}
\label{report-step-UnadjustedLangevin-LSI-Regularity-lean-4}


The local regularity data include integrability of \(r_t\log r_t\). Positivity and \(C^1\) regularity give a measurable nonnegative density. Apply the finite-KL criterion from \(\int r|\log r|d\nu<\infty\) to obtain finite actual KL before using it as an ordinary finite entropy.



The regularity assumptions control the entropy integrand at the selected time. Together with positivity and normalization, they identify a genuine finite relative entropy before any right-derivative or scalar comparison is used.

\subsection{\texorpdfstring{Continuity of actual entropy on the interval}{Continuity of actual entropy on the interval}}
\label{report-step-UnadjustedLangevin-LSI-Regularity-lean-5}


At every point of the closed interval, dominated differentiation makes the ordinary entropy integral differentiable and hence continuous. Normalization and the finite-entropy density formula identify that integral with actual KL entropy at every interval point. congr transfers continuity to \(t\mapsto\operatorname{KL}(\rho_t\mid\nu)\).



Differentiability implies continuity of the entropy integral locally. Pointwise equality with actual KL on the whole stated interval transfers that continuity to the entropy path, including its required one-sided endpoint behavior.

\subsection{\texorpdfstring{The one-\allowbreak{}sided derivative of actual entropy}{The one-sided derivative of actual entropy}}
\label{report-step-UnadjustedLangevin-LSI-Regularity-lean-6}


The ordinary entropy integral has the derivative computed by local regularity. Its equality with actual KL is known only for normalized densities in \([0,T]\), so the proof works with right derivatives at \(t\in[0,T)\).



In a sufficiently small right neighborhood, \(u\ge t\ge0\) and \(u<T\), so the density formula identifies both functions there and at \(t\). equality of derivatives for functions agreeing on a neighborhood transfers the derivative to actual KL. This avoids requiring path normalization outside the interpolation interval.



The density law is only normalized on the admissible interval. A sufficiently small right neighborhood of an interior point remains there, so equality with KL can transfer a right derivative without asserting anything about densities outside the interval.

\subsection{\texorpdfstring{Production in terms of relative Fisher information}{Production in terms of relative Fisher information}}
\label{report-step-UnadjustedLangevin-LSI-Regularity-lean-7}


Convert squared-score integrability under \(\rho_t\) into weighted Fisher integrability under \(\nu\) by the weighted-density criterion. Apply the regular weak-equation derivative estimate, and identify its weighted integral with relative fisher information's real value. Pair that estimate with the actual-KL right derivative from the preceding theorem to get \(H'(t)\le-3I(t)/4+E(t)\) for the identified density law.



The energy under the tilted law and its weighted reference form are equal by change of density. This identifies the production bound with the same Fisher quantity used in LSI, ensuring the two inequalities may be combined directly.

\section{\texorpdfstring{Derive entropy calculus from regular interpolation hypotheses}{Derive entropy calculus from regular interpolation hypotheses}}
\label{proof-module-UnadjustedLangevin-LSI-RegularizedTheorem-lean}


This result replaces a supplied entropy derivative estimate with explicit regular density data, a forced weak identity and an error-energy bound. It then invokes the conditional ULA entropy endpoint. Constructing those regular interpolation data remains a displayed premise.



This route supplies a density path and an admissible weak equation, then derives the entropy derivative, production inequality, one-step recurrence, and convergence bound in that order. Its additional regularity assumptions make differentiation under the integral available directly. The conclusion concerns the actual ULA laws only after the supplied density is identified with their interpolation. This conditional organization retains the analytic construction and energy estimates as visible hypotheses.

\subsection{\texorpdfstring{Construct the entropy-\allowbreak{}calculus hypotheses and invoke the actual-\allowbreak{}law theorem}{Construct the entropy-calculus hypotheses and invoke the actual-law theorem}}
\label{report-step-UnadjustedLangevin-LSI-RegularizedTheorem-lean-1}


First use the density/\allowbreak{}law identity to prove probability normalization along each interpolation. Local regularity proves finite actual KL, score integrability proves finite Fisher, and continuity obtained from differentiating the entropy integral gives continuity of actual entropy.



Define the candidate derivative \(H'_k(t)=\int\dot r_{k,t}(\log r_{k,t}+1)d\pi_U\). The regular weak-equation theorem proves that this is the actual right derivative and that it is bounded by \(-3I/4+E\). Insert the supplied drift-energy bound \(E\le t^2L^2G_k+2tdL^2\), then invoke the entropy convergence theorem for an identified interpolation.



The final estimate is therefore derived without assuming an entropy derivative, a production inequality, or a discrete recurrence. This result still explicitly requires the regular density construction, its admissible weak equation, its drift-energy estimate, and the entropy-to-gradient moment bound; the result does not itself construct them from LSI.



The law identity is used before entropy calculus, ensuring that every density path represents the true numerical interpolation. Dominated differentiation then supplies the entropy derivative, the admissible weak equation supplies production, and the separate numerical energy estimates supply its forcing. Only after these steps is the scalar convergence theorem invoked. The construction hypotheses remain explicit in this conditional result.



Earlier results used here: \cite[Section~3.26.3]{LangevinLean2026ULASC}.
\section{\texorpdfstring{A reusable regular-\allowbreak{}interpolation interface}{A reusable regular-interpolation interface}}
\label{proof-module-UnadjustedLangevin-LSI-InterpolationData-lean}


The analytic interpolation specification bundles density paths identifying the actual ULA interpolation, regularity, weak-equation and error controls. Its methods derive path probability, finite KL, the one-step recurrence and the iterated entropy bound. Existence of this package is not asserted by its definition.



The interpolation data identify each density path with the actual law starting at \(\rho_k\). They also contain positivity, spatial regularity, a local domination condition for entropy differentiation, and finite score and error energies. These assumptions permit the entropy calculation to be performed; the derivative identity and scalar recurrence are derived from them. Keeping the law identity explicit prevents a bound for an unrelated smooth density curve from being mistaken for a statement about the algorithm.

\subsection{\texorpdfstring{The analytic interpolation hypotheses}{The analytic interpolation hypotheses}}
\label{report-step-UnadjustedLangevin-LSI-InterpolationData-lean-1}


Package the relative densities \(r_{k,t}\), their time derivatives, and drift fields together with the exact law identity, positivity, spatial regularity, local entropy domination, finite score/\allowbreak{}error energies, admissible weak equation, drift-energy bound, and gradient-moment bound. These remain explicit mathematical inputs. No entropy derivative, entropy recurrence, or convergence estimate is a assumption.



The relative density is linked to the algorithm at every interpolation time, not only at its endpoint. This permits the same LSI to act along the path and identifies the resulting scalar endpoints with successive actual entropies. The listed energy and domination assumptions are sufficient inputs for deriving the calculation.



Earlier results used here: \cite[Section~3.26.3]{LangevinLean2026ULASC}.
\subsection{\texorpdfstring{Probability normalization of the represented path}{Probability normalization of the represented path}}
\label{report-step-UnadjustedLangevin-LSI-InterpolationData-lean-2}


The law identity identifies the tilted density with the actual interpolation from \(\rho_k\). Probability preservation proves that interpolation is normalized. Rewrite by the law identity to obtain normalization of each density along the supplied path.



The path's weighted measure equals an actual Gaussian transition law. Because that transition preserves probability, its total mass is one, which proves normalization of the relative density without adding a separate integral-one axiom.



Earlier results used here: \cite[Section~3.26.3]{LangevinLean2026ULASC}.
\subsection{\texorpdfstring{Finite KL at every represented iterate}{Finite KL at every represented iterate}}
\label{report-step-UnadjustedLangevin-LSI-InterpolationData-lean-3}


Use local entropy regularity to prove finite KL for each relative density, then apply the time-zero law identification in the finite-entropy conclusion for the identified time-zero law. The target probability hypothesis is explicit; in the final theorem it is supplied by the LSI hypothesis.



Local entropy regularity gives finite KL for the represented density. Evaluating the exact law identification at elapsed time zero transfers this finiteness to the actual algorithm law at each step.



Earlier results used here: \cite[Section~3.26.3]{LangevinLean2026ULASC}.
\subsection{\texorpdfstring{Deriving the one-\allowbreak{}step estimate from the analytic data}{Deriving the one-step estimate from the analytic data}}
\label{report-step-UnadjustedLangevin-LSI-InterpolationData-lean-4}


Define \(H'_k(t)\) by the differentiated entropy integral. The regular weak-equation theorem proves its right-derivative identity and Fisher/\allowbreak{}drift production bound. The assumptions supply probability, KL/\allowbreak{}Fisher finiteness, continuity, numerical error-energy and gradient bounds. Feed these proved ingredients to the single-step entropy comparison for the identified law to derive the recurrence rather than assuming it in the data.



The candidate derivative is the differentiated entropy integral. Its equality with the genuine right derivative and its production bound are proved from the supplied regularity and weak equation, after which scalar comparison gives the one-step recurrence.



Earlier results used here: \cite[Section~3.26.3]{LangevinLean2026ULASC}.
\subsection{\texorpdfstring{The iterated bound from the analytic data}{The iterated bound from the analytic data}}
\label{report-step-UnadjustedLangevin-LSI-InterpolationData-lean-5}


Unpack every analytic field into the entropy convergence theorem for an identified regular density path, retaining LSI, potential smoothness, and the step regime as separate arguments. Its already proved differentiation, one-step estimate, and geometric iteration give the full displayed KL bound for the actual ULA laws represented by the data.



The full convergence estimate uses the already derived recurrence and its geometric sum. No recurrence is assumed among the data: the law, regularity, weak equation, and energy conditions are the assumptions from which it is constructed.



Earlier results used here: \cite[Section~3.26.3]{LangevinLean2026ULASC}.
\section{\texorpdfstring{Compare entropy using derivatives only at positive times}{Compare entropy using derivatives only at positive times}}
\label{proof-module-UnadjustedLangevin-LSI-InteriorEntropy-lean}


Gaussian regularization naturally supplies derivatives for positive times. The scalar comparison here uses ordinary derivatives on \(0<t<T\) and continuity on \([0,T]\); neither endpoint needs a derivative. A second theorem computes actual KL derivatives from local relative-density data only in that open interval. The existing global ULA endpoint retains its earlier interface; these lemmas provide a more flexible route for a future positive-time construction.



At an interior positive time, local entropy differentiation gives an ordinary derivative after the required integrable domination is checked. On a compact positive interval, multiplying the affine differential inequality by an integrating factor yields a comparison between its endpoints. This avoids requiring a derivative at either endpoint. Passing the lower endpoint to zero is a distinct analytic step, supplied later by a one-sided entropy majorant for the actual initial law.

\subsection{\texorpdfstring{Use an integrating factor without differentiating at the endpoints}{Use an integrating factor without differentiating at the endpoints}}
\label{report-step-UnadjustedLangevin-LSI-InteriorEntropy-lean-1}


Let \(a>0\), let \(H\) be continuous on \([0,T]\), and suppose it has derivative \(H'(t)\le-aH(t)+b\) for \(0<t<T\). Then for every \(t\in[0,T]\),

\[H(t)\le e^{-at}H(0)+\frac ba(1-e^{-at}).\]

Define \(F(t)=e^{at}(H(t)-b/a)\). The product and chain rules give

\[F'(t)=e^{at}\bigl(a(H(t)-b/a)+H'(t)\bigr)\le0.\]

Consequently \(F(t)\le F(0)=H(0)-b/a\). Multiplication by the positive factor \(e^{-at}\), together with \(e^{-at}e^{at}=1\), yields the stated affine comparison. The final arithmetic rearrangement is nlinarith. Endpoint continuity is essential in this proof; endpoint differentiability is unnecessary.



For \(a<t<b\), differentiate \(e^{ct}H(t)\) only at interior times. Integrating its upper bound and using continuity at \(a,b\) gives the endpoint comparison. Thus no derivative at either endpoint is needed.

\subsection{\texorpdfstring{Differentiate actual KL at an interior interpolation time}{Differentiate actual KL at an interior interpolation time}}
\label{report-step-UnadjustedLangevin-LSI-InteriorEntropy-lean-2}


Fix \(0<t<T\). Assume that, at every interior time \(u\), \(r_u\) is positive, continuously differentiable, normalized relative to the probability target \(\nu\), and has the local dominated entropy regularity recorded in local hypotheses for differentiating entropy. Then

\[\frac d{dt}D_{\mathrm{KL}}(r_t\nu\Vert\nu)=\int \dot r_t(x)\bigl(\log r_t(x)+1\bigr)\,d\nu(x).\]

equality of derivatives for functions agreeing on a neighborhood transfers the already proved integral derivative to actual KL. The statement requests no density or local domination information at time zero. To combine it with the preceding scalar theorem, continuity of the actual entropy at the two endpoints still has to be justified separately.



At a strictly positive time, the local regularity neighborhood stays within positive times. The dominated entropy derivative formula therefore applies to the actual smoothed density, with its normalization converting the integral into genuine KL.

\section{\texorpdfstring{Transfer the Gibbs score moment through a transport coupling}{Transfer the Gibbs score moment through a transport coupling}}
\label{proof-module-UnadjustedLangevin-LSI-GradientMoment-lean}


For a coupling \((X,Y)\) of \(\mu\) with Gibbs, \(\|\nabla U(X)\|^2\le2L^2\|X-Y\|^2+2\|\nabla U(Y)\|^2\). Integrating first proves score integrability under \(\mu\), then bounds it by a coupling cost and the Gibbs score moment \(Ld\). Taking the transport infimum and inserting a supplied entropy-to-transport bound produces the stated entropy-dependent coefficient. This route keeps that transport premise explicit; the final actual-law proof also has a direct exponential-score route that needs no such premise.



Earlier results used here: \cite[Section~3.17]{LangevinLean2026ULASC}; \cite[Section~3.27]{LangevinLean2026ULASC}.
\subsection{\texorpdfstring{Derive both integrability and the coupling moment estimate}{Derive both integrability and the coupling moment estimate}}
\label{report-step-UnadjustedLangevin-LSI-GradientMoment-lean-1}


Let \(c\) be a joint measure with marginals \(\mu\) and \(\nu\). Assume \(\int\|x-y\|^2dc<\infty\), that the gradient is \(L\)-Lipschitz, and that \(\int\|\nabla U\|^2d\nu<\infty\). The theorem proves square integrability under \(\mu\) and

\[\int\|\nabla U\|^2d\mu\le2L^2\int\|x-y\|^2\,dc(x,y)+2\int\|\nabla U\|^2d\nu.\]

The Lipschitz gradient gives \(\|\nabla U(X)\|\le L\|X-Y\|+\|\nabla U(Y)\|\). Squaring with \((a+b)^2\le2a^2+2b^2\) and integrating proves both finite energy and the coupling bound before optimizing over couplings.



Earlier results used here: \cite[Section~3.17.1]{LangevinLean2026ULASC}.
\subsection{\texorpdfstring{Insert the derived target score bound}{Insert the derived target score bound}}
\label{report-step-UnadjustedLangevin-LSI-GradientMoment-lean-2}


For the normalized Gibbs probability \(\nu=Z^{-1}e^{-U}dx\), differentiability of \(U\) and an \(L\)-Lipschitz gradient imply the target estimate \(\int\|\nabla U\|^2d\nu\le nL\). Consequently every finite-cost coupling gives

\[\int\|\nabla U\|^2d\mu\le2L^2\int\|x-y\|^2dc+2nL.\]

The proof first derives \(Z<\infty\) from probability normalization, using finiteness of the partition forced by target probability normalization. The shared theorem square-integrability of the gradient under Gibbs then supplies target square integrability; its \(L^2\) conclusion is converted into the integrability input for the preceding coupling theorem. The sharp target moment itself comes from the target gradient-energy bound \(\int\|\nabla U\|^2d\nu\le dL\).



Those shared Gibbs results use translation invariance and a sub-Gaussian score argument under first-order smoothness, so strong convexity is unnecessary here. Substituting their target bound into the coupling estimate proves the stated result and retains the newly derived square integrability under \(\mu\).



The target-gradient term is evaluated under the Gibbs marginal of the coupling, so the score estimate makes it at most \(Ld\). This removes dependence on the chosen coupling from that part of the energy bound.



Earlier results used here: \cite[Section~3.17.1]{LangevinLean2026ULASC}; \cite[Section~3.26.3]{LangevinLean2026ULASC}; \cite[Section~3.27.6]{LangevinLean2026ULASC}; \cite[Section~3.27.7]{LangevinLean2026ULASC}.
\subsection{\texorpdfstring{Obtain the exact entropy-\allowbreak{}dependent coefficient from transport control}{Obtain the exact entropy-dependent coefficient from transport control}}
\label{report-step-UnadjustedLangevin-LSI-GradientMoment-lean-3}


Suppose the preceding finite-cost coupling additionally satisfies

\[\int\|x-y\|^2dc(x,y)\le\frac2\alpha D_{\mathrm{KL}}(\mu\Vert\nu),\qquad \alpha>0,\quad D_{\mathrm{KL}}(\mu\Vert\nu)<\infty.\]

Substitution into the coupling bound gives

\[\int\|\nabla U\|^2d\mu\le\frac{4L^2}{\alpha}D_{\mathrm{KL}}(\mu\Vert\nu)+2nL.\]

The supplied transport estimate has a specified entropy constant. Inserting it into the factor \(2L^2\) from the coupling calculation produces exactly \(4L^2H/\alpha\), while the target-score contribution remains \(2dL\).



Earlier results used here: \cite[Section~3.17.1]{LangevinLean2026ULASC}; \cite[Section~3.26.3]{LangevinLean2026ULASC}.

\begin{thebibliography}{ULA-LSI}
\bibitem[ULA-SC]{LangevinLean2026ULASC}
LangevinLean Project. \emph{Unadjusted Langevin under Strong Convexity}. LL-01, 2026. Draft report, content revision 8d6335dce333. \url{https://odakyildiz.github.io/LangevinLean/docs/pdf/UnadjustedLangevin/FullProof.pdf}.
\end{thebibliography}

\end{document}
